\pdfoutput=1
\documentclass[11pt]{amsart}

\usepackage[margin=1.1in]{geometry}
\usepackage{amsmath,amssymb,amsthm,mathtools,mathrsfs}
\usepackage{array,booktabs,longtable}
\usepackage{enumitem}
\usepackage{float}
\usepackage{microtype}
\usepackage{aliascnt}
\usepackage{tikz}
\usetikzlibrary{arrows.meta,positioning,shapes.geometric}
\usepackage{hyperref}
\usepackage[nameinlink,capitalize]{cleveref}
\newcolumntype{L}[1]{>{\raggedright\arraybackslash}p{#1}}

\hypersetup{
  colorlinks=true,
  linkcolor=blue,
  citecolor=blue,
  urlcolor=blue,
  pdftitle={Covariant Stretching--Diffusion Rigidity and Global Regularity for Navier--Stokes},
  pdfauthor={Guillem Duran Ballester},
  pdfsubject={A covariant equality and material-rigidity approach to Navier--Stokes regularity},
  pdfkeywords={Navier--Stokes equations, global regularity, material deformation, vortex stretching, anisotropic diffusion, entropy}
}
\raggedbottom
\emergencystretch=2em
\numberwithin{equation}{section}

\theoremstyle{plain}
\newtheorem{theorem}{Theorem}[section]
\newaliascnt{proposition}{theorem}
\newtheorem{proposition}[proposition]{Proposition}
\aliascntresetthe{proposition}
\newaliascnt{lemma}{theorem}
\newtheorem{lemma}[lemma]{Lemma}
\aliascntresetthe{lemma}
\newaliascnt{corollary}{theorem}
\newtheorem{corollary}[corollary]{Corollary}
\aliascntresetthe{corollary}
\theoremstyle{definition}
\newaliascnt{definition}{theorem}
\newtheorem{definition}[definition]{Definition}
\aliascntresetthe{definition}

\theoremstyle{remark}
\newaliascnt{remark}{theorem}
\newtheorem{remark}[remark]{Remark}
\aliascntresetthe{remark}

\newcommand{\R}{\mathbb R}
\newcommand{\dd}{\,\mathrm d}
\newcommand{\Pdiv}{\mathbb P}
\newcommand{\diver}{\operatorname{div}}
\newcommand{\esssup}{\operatorname*{ess\,sup}}
\newcommand{\supp}{\operatorname{supp}}
\newcommand{\Energy}{\mathcal E}
\newcommand{\Action}{\mathcal A}
\newcommand{\ProbeDirichlet}{\mathcal G}
\newcommand{\Power}{\mathcal P}
\newcommand{\Align}{\operatorname{Align}}
\newcommand{\RayDefect}{\mathscr D^{\mathrm{ray}}}
\newcommand{\Dcov}{\mathfrak D}
\newcommand{\Tgeom}{\mathcal T}
\newcommand{\EntropyNS}{\boldsymbol\mu_{\mathrm{NS}}}
\newcommand{\ProductionNS}{\mathfrak D_{\mathrm{NS}}}
\newcommand{\OscEnergy}{\mathcal E^{\circ}}
\newcommand{\OscMass}{\mathcal C^{\circ}}
\newcommand{\MaterialF}{\mathbf F}
\newcommand{\MaterialMetric}{\mathbf C}
\newcommand{\MaterialVort}{\boldsymbol\Omega}
\newcommand{\MaterialEntropy}{\mathscr H_{\mathrm{cov}}}
\newcommand{\SingT}[1]{\mathfrak S_{\mathrm{sing},T}\!\left(#1\right)}
\title[Covariant Stretching--Diffusion Rigidity]
{Covariant Stretching--Diffusion Rigidity and Global Regularity\\
for the Three-Dimensional Navier--Stokes Equations}

\author{Guillem Duran Ballester}
\email{guillem@fragile.tech}
\date{}

\subjclass[2020]{35Q30, 35B44, 35B40, 35K55, 76D05}
\keywords{Navier--Stokes equations, global regularity, suitable weak solutions,
material deformation, stretching--diffusion duality, covariant dissipation,
concentration compactness, rigidity, blow-up coordinates}

\begin{document}

\begin{abstract}
We give a covariant proof of global regularity for the three-dimensional
incompressible Navier--Stokes equations organized around one normalized
observer approaching a putative singularity.  Translation, parabolic scale,
constant Galilean velocity, and pressure gauge are fixed coherently, while a
backward adjoint localization retains a scale-critical quotient oscillation.
Every retained unit window has a uniform positive viscous price in normalized
critical units.

The moving-root local-energy identity records exactly how that price is
financed.  Its signed current contains pressure, nonlinear transport,
dilation, translation, and cutoff motion.  Summation over successive windows
shows that sublinear cumulative current cannot support infinitely many active
events.  Persistent current is retained together with the observer as marked
PDE data.  Loss of compactness in center, pressure, scale, or defect variables
produces a classified local output; repeated regeneration of active centers
produces a sequence-realized logarithmic critical tail.  The diffuse,
homogeneous, defect, log-periodic, and aperiodic tail alternatives are rigid,
so every persistent-current output is trivial.  This contradicts the retained
oscillation and excludes finite-time blow-up.

The proof is developed through five structural theorems: canonical
observation, a positive critical dissipation price, the exact financing
identity, persistent-current compactification, and critical-tail rigidity.
The appendices derive these statements directly from the suitable local-energy
inequality, pressure reconstruction, parabolic compactness, and
\(\varepsilon\)-regularity.  Material coordinates provide the geometric
interpretation: vortex stretching and viscous diffusion are governed by the
dual tensors \(\nabla X\) and
\((\nabla X^{T}\nabla X)^{-1}\), and zero production is spatially homogeneous
modulo the Galilean quotient.
\end{abstract}

\maketitle
\enlargethispage{2pt}
\tableofcontents

\section{Introduction}

Consider the incompressible Navier--Stokes equations on \(\R^3\),
\begin{equation}\label{eq:NS}
  \partial_t u+(u\cdot\nabla)u+\nabla p=\nu\Delta u,
  \qquad \diver u=0,
  \qquad \nu>0.
\end{equation}
Three geometric observations organize the argument.  Arnold's
description of Euler flow treats incompressible transport as conservative
motion on the group of volume-preserving diffeomorphisms
\cite{Arnold1966,EbinMarsden1970}.  Onsager's variational description of
irreversible evolution identifies a dissipative current with a gradient in a
mobility metric \cite{Onsager1931}.  For Navier--Stokes these statements fit
together exactly:
\begin{equation}\label{eq:intro-covariant-flow}
  \underbrace{\partial_t u+\Pdiv((u\cdot\nabla)u)}_{
    \text{Euler-covariant derivative}}
  =-\underbrace{\nu(-\Pdiv\Delta)}_{
    \text{Stokes mobility}}\,D\Energy(u),
  \qquad
  \Energy(u)=\frac12\|u\|_2^2.
\end{equation}
The conservative term is orthogonal to \(D\Energy(u)=u\), while the squared
metric norm of the right-hand side is \(\nu\|\nabla u\|_2^2\).  Thus the
energy inequality is the energy--dissipation identity associated with
\eqref{eq:intro-covariant-flow}.

The third observation is specific to incompressible vorticity.  In material
coordinates the deformation tensor \(\MaterialF=\nabla_aX\) transports
vorticity, while viscosity is pulled back by the inverse Cauchy--Green metric
\(\MaterialMetric^{-1}=(\MaterialF^T\MaterialF)^{-1}\).  Since
\(\det\MaterialF=1\), stretching and enhanced diffusion are dual aspects of
the same volume-preserving deformation.  This is the structural mechanism
developed in \cref{sec:material-duality}.

\begin{center}
\small
\begin{tabular}{@{}L{0.42\textwidth}L{0.48\textwidth}@{}}
\toprule
Ricci-flow viewpoint & Navier--Stokes realization \\
\midrule
Canonical observer near a singularity & Scale, center, Galilean mode, and pressure gauge \\
Conjugate heat kernel & Backward adjoint localization \\
Entropy or reduced-volume monotonicity & Covariant flux--dissipation identity \\
Noncollapsing & Retained scale-critical quotient oscillation \\
Blow-up limit & Ancient or recurrent normalized flow \\
Soliton at equality & Zero-production covariant equality state \\
Soliton rigidity & Spatial homogeneity modulo the Galilean quotient \\
Singular-scale analysis & Marked logarithmic current and critical-tail analysis \\
\bottomrule
\end{tabular}
\end{center}

The physical Stokes action alone does not exclude concentration because its
dissipation on a normalized event at scale \(r\) is weighted by \(r\).  A normalized
gap and a physical lower bound are therefore distinct assertions.  In
particular, cancellation of fluxes between nested cells cannot by itself
replace this factor by a scale-independent constant.  The required quantity
is instead the net dissipation in a scale-critical localized-energy balance.
The proof has four steps.
\begin{enumerate}[label=\textup{\arabic*.},leftmargin=2.1em]
\item \emph{Canonical observation.}  Terminal loss generates one coherent
scale--center path.  Translation, dilation, Galilean velocity, and pressure
gauge are fixed, while a backward adjoint localization retains a uniform
quotient-oscillation floor.
\item \emph{Critical dissipation price.}  Compactness and zero-production
rigidity turn retained nonconstant oscillation into
\[
 \nu\int_{I_j}\int\zeta^2|\nabla V|^2\ge d_*>0
 \]
on every active window.  This is scale invariant; the corresponding physical
energy is only \(r_jd_*\).
\item \emph{Exact financing.}  The moving-root local-energy equality is
\[
 K_j^+-K_j^-+D_j+M_j=H_j,
 \]
where \(H_j\) contains the complete signed pressure, transport, dilation,
translation, cutoff, and transition current.  After summation, bounded
endpoint storage and the fixed price rule out sublinear cumulative current.
\item \emph{Recurrence--tail rigidity.}  Persistent current either loses a
first center, pressure, scale, exterior, trace, or defect coordinate, or
returns through another active center.  The first failed coordinate is
retained at its first parabolic scale.  In the recurrent case the observer
and current are retained as marks.  Both constructions produce a realized
logarithmic critical tail, and critical-tail rigidity makes that tail
target-null.
\end{enumerate}

Thus the entire mechanism is
\begin{equation}\label{eq:structural-proof-spine}
\boxed{\;
\begin{gathered}
\text{retained critical activity}
\Longrightarrow \text{positive critical price}\\
\Longrightarrow \text{persistent signed current}
\Longrightarrow \text{recurrent critical tail}
\Longrightarrow \text{rigidity}.
\end{gathered}
\;}
\end{equation}

This separation starts from the localized Navier--Stokes energy identity.
Its convection, pressure, cutoff-motion, chart, dissipation, and defect terms
all depend linearly on the localization test.  A compatible parent--child
partition therefore prevents double counting but does not, by itself, cancel
only the signed terms.  Selective cancellation is a statement about oriented
normal traces on a disjoint spacetime partition.  Viscous dissipation and the
suitable-energy defect then retain their sign on the disjoint interiors.
Section~\ref{sec:dissipation-reduction} develops this local energy decomposition entirely within
the Navier--Stokes equation and isolates the target-visible zero-dissipation
problem.

The rigidity in the fourth step is the covariant equality calculation.
Ordered active cylinders admit one exact chart with nonincreasing scale.  On
an autonomous compact recurrent limit, covariant spatial averaging removes
convection, translation, and every nonzero spatial pressure mode, leaving
\[
 \nu\langle|\nabla V|^2\rangle
 -a_*\langle|V-\Pi_0V|^2\rangle=0.
\]
For \(a_*\le0\) this forces spatial constancy and hence loss of the
mean-subtracted singular target.  This pressure-cancellation argument is
proved in \cref{sec:recurrent-covariant-rigidity}.

The proof uses one coherent observer and one accounting identity.  The
observer retains scale-critical quotient oscillation; the identity says that
its fixed viscous price must be financed by signed current.  Persistent current
either appears as a classified loss of compactness or returns through smaller
scales and generates a logarithmic critical tail.  Tail rigidity supplies the
final contradiction.

\begin{theorem}[Global regularity]
\label{thm:intro-main}
Let \(m>5/2\), \(\nu>0\), and let
\(u_0\in H^m(\R^3;\R^3)\) be divergence free.  Then \eqref{eq:NS} has a
unique global classical solution
\[
 u\in C([0,\infty);H^m)\cap C^\infty(\R^3\times(0,\infty)).
\]
Every Leray--Hopf solution with initial datum \(u_0\) agrees with this
classical solution.
\end{theorem}

The proof of \cref{thm:intro-main} is given in
\cref{sec:global-exclusion}.  It invokes the following five structural
theorems, each proved in this paper:
\begin{enumerate}[label=\textup{(\roman*)},leftmargin=2.2em]
\item canonical observation, \cref{thm:canonical-critical-orbit};
\item the positive critical price, \cref{thm:uniform-production-gap};
\item the exact financing identity,
\cref{thm:exact-critical-financing};
\item persistent-current compactification,
\cref{thm:persistent-current-compactification};
\item critical-tail rigidity, \cref{thm:critical-tail-rigidity}.
\end{enumerate}

\begin{remark}[Logical scope]
\label{rem:logical-scope}
No regularity theorem from another manuscript is an input.  The appendices
derive the five structural theorems from the suitable local-energy inequality,
pressure reconstruction, parabolic compactness, \(\varepsilon\)-regularity,
and the ancient-solution rigidity statements cited there.  The finite-depth
source decomposition proves that persistent-current compactification is
exhaustive; each of its outputs is resolved inside critical-tail rigidity.
\end{remark}

\subsection*{The five structural theorems}

\begin{center}
\small
\begin{tabular}{@{}L{0.28\textwidth}L{0.25\textwidth}L{0.35\textwidth}@{}}
\toprule
Theorem & Geometric content & Derivation in this paper \\
\midrule
Canonical observation & One scale--center--gauge path retains quotient activity & Appendix~\ref{app:critical-orbit} \\
Critical price & Every retained normalized window pays \(d_*>0\) & \Cref{sec:entropy-gap,app:material-rigidity} \\
Exact financing & Stored energy plus dissipation and defect equals complete signed current & \Cref{sec:dissipation-reduction,app:scale-space-current} \\
Persistent-current compactification & Sublinear current is impossible; persistent current yields a classified loss or recurrent tail & \Cref{sec:finite-depth-PDE-decomposition,app:marked-regeneration,app:collapsed-residual-closure} \\
Critical-tail rigidity & Every realized diffuse, defect, homogeneous, log-periodic, or aperiodic tail is target-null & \Cref{app:terminal-record-exclusion,app:material-rigidity} \\
\bottomrule
\end{tabular}
\end{center}

\subsection*{Proof architecture}

Rectangles are analytic states, diamonds are exhaustive alternatives, and
solid ellipses are contradictions.  Every edge belongs to the single
observer--price--current--tail argument.

\clearpage
\noindent\textbf{Proof-dependency diagram.}\par\smallskip

\begin{figure}[H]
\centering
\begin{tikzpicture}[
 scale=0.92,transform shape,
 box/.style={rectangle,rounded corners,draw,align=center,inner sep=3pt,text width=4.0cm},
 dec/.style={diamond,draw,aspect=2.35,align=center,inner sep=1.5pt,text width=3.15cm},
 term/.style={ellipse,draw,align=center,inner sep=3pt,text width=3.8cm},
 arrow/.style={-{Latex[length=2mm]},thick},
 every node/.style={font=\small}
]
\node[box] (n1) at (0,9.2) {\textbf{[1]} Putative finite-time singularity};
\node[box] (n2) at (0,7.5) {\textbf{[2]} Canonical observer with retained quotient activity};
\node[box] (n3) at (0,5.8) {\textbf{[3]} Fixed positive critical price on every active window};
\node[box] (n4) at (0,4.1) {\textbf{[4]} Exact moving-root financing identity};
\node[dec] (n5) at (0,1.8) {\textbf{[5]} Cumulative signed current};
\node[term] (n6) at (-5.0,-0.7) {\textbf{[6]} Sublinear current contradicts the accumulated price};
\node[box] (n7) at (4.9,-0.5) {\textbf{[7]} Persistent current with complete observer marks};
\node[dec] (n8) at (4.9,-3.0) {\textbf{[8]} Compactness loss or active recurrence};
\node[box] (n9) at (0,-5.2) {\textbf{[9]} First failed coordinate produces a realized tail};
\node[box] (n10) at (5.4,-5.2) {\textbf{[10]} Sequence-realized logarithmic critical tail};
\node[term] (n11) at (2.7,-7.2) {\textbf{[11]} Critical-tail rigidity contradicts retained activity};
\draw[arrow] (n1)--(n2);
\draw[arrow] (n2)--(n3);
\draw[arrow] (n3)--(n4);
\draw[arrow] (n4)--(n5);
\draw[arrow] (n5)--node[above left]{\scriptsize sublinear}(n6);
\draw[arrow] (n5)--node[above right]{\scriptsize persistent}(n7);
\draw[arrow] (n7)--(n8);
\draw[arrow] (n8)--node[above left]{\scriptsize first failure}(n9);
\draw[arrow] (n8)--node[right]{\scriptsize recurrence}(n10);
\draw[arrow] (n9)--(n11);
\draw[arrow] (n10)--(n11);
\end{tikzpicture}
\caption{The single covariant proof spine.  Persistent current cannot be lost:
its first failed coordinate is retained as a classified PDE output, while
continued active recurrence produces a realized critical tail.}
\label{fig:proof-flow-marked-current}
\end{figure}

\noindent\textbf{Proof-state source table.}\par\smallskip

\begin{center}
\small
\begin{tabular}{@{}L{0.09\textwidth}L{0.32\textwidth}L{0.46\textwidth}@{}}
\toprule
State & Analytic content & Exact source and successor \\
\midrule\relax
[1] & Terminal loss & \Cref{thm:velocity-concentration}; continue to [2]. \\\relax
[2] & Coherent normalization & \Cref{thm:canonical-critical-orbit}; continue to [3]. \\\relax
[3] & Critical price & \Cref{thm:uniform-production-gap}; continue to [4]. \\\relax
[4] & Exact accounting & \Cref{thm:exact-critical-financing}; continue to [5]. \\\relax
[5] & Current growth test & \Cref{lem:sublinear-critical-current-closure}; outputs [6] or [7]. \\\relax
[6] & Insufficient financing & The left side grows at least like \(Nd_*\), while the right side is \(o(N)\); terminal. \\\relax
[7] & Persistent current & \Cref{lem:persistent-current-support-alternative}; continue to [8]. \\\relax
[8] & Exhaustive compactification & \Cref{thm:persistent-current-compactification}; outputs [9] or [10]. \\\relax
[9] & First failed coordinate & \Cref{thm:finite-depth-NS-source-alternative,prop:first-bad-scale-descendant}; continue to [11]. \\\relax
[10] & Recurrent marked path & \Cref{thm:sparse-moving-root-path-exclusion}; produces a realized tail and continues to [11]. \\\relax
[11] & Tail rigidity & \Cref{thm:critical-tail-rigidity}; terminal. \\
\bottomrule
\end{tabular}
\end{center}

The facts carried from [2] onward are the fixed quotient-oscillation floor,
realization by the same physical terminal sequence, compatible pressure
gauges, the exact scale--center composition law, and the complete signed
local-energy current.  Thus neither a subsequence nor a recurrent limit is
permitted to forget the moving root or its pressure contribution.

\medskip\noindent\textbf{Exhaustion table.}\par\smallskip

\begin{center}
\small
\begin{tabular}{@{}L{0.14\textwidth}L{0.29\textwidth}L{0.43\textwidth}@{}}
\toprule
Case & Retained output & Closure used by \cref{thm:intro-main} \\
\midrule
Sublinear current & The complete signed current is \(o(N)\) over \(N\) retained events & \Cref{lem:sublinear-critical-current-closure}. \\
Compactness loss & Pressure, center, scale, exterior, trace, or defect coordinate first fails & \Cref{thm:finite-depth-NS-source-alternative,prop:first-bad-scale-descendant,thm:critical-tail-rigidity}. \\
Active recurrence & Persistent current repeatedly reaches another retained center & \Cref{lem:persistent-current-support-alternative,thm:sparse-moving-root-path-exclusion,thm:critical-tail-rigidity}. \\
\bottomrule
\end{tabular}
\end{center}

\medskip\noindent\textbf{First-failure source table.}\par\smallskip

The following table expands state [9].  A failed coordinate is retained as PDE
data and passed to the critical-tail theorem; it is never treated as an
unproved terminal alternative.

\begin{center}
\small
\begin{tabular}{@{}L{0.10\textwidth}L{0.34\textwidth}L{0.43\textwidth}@{}}
\toprule
Output & PDE state & Status \\
\midrule
F1 & Homogeneous velocity in the Galilean quotient & Target-null by definition of \(\OscMass_\Gamma\). \\
F2 & Compact nonexpanding recurrence & Zero-production rigidity by \Cref{thm:recurrent-nonexpanding-rigidity}. \\
F3 & Pressure, exterior, scale, trace, or defect loss & First visible parabolic scale gives a sequence-realized tail by \Cref{thm:finite-depth-NS-source-alternative,prop:first-bad-scale-descendant}. \\
F4 & Same-origin scale collapse with persistent current & Eventwise critical price and marked recurrence give a realized tail by \Cref{thm:collapsed-same-origin-residual-exclusion}. \\
\bottomrule
\end{tabular}
\end{center}

\section{Suitable solutions and the singular target}
\label{sec:singular-target}

\subsection{Suitable weak solutions}

We use the standard finite-energy suitable class
\cite{Leray1934,CaffarelliKohnNirenberg1982,Lin1998}.  Thus
\begin{equation}\label{eq:suitable-class}
 u\in L^\infty(0,T;L^2_\sigma(\R^3))
 \cap L^2(0,T;\dot H^1_\sigma(\R^3)),
 \qquad p\in L^{3/2}_{\mathrm{loc}}(\R^3\times(0,T)),
\end{equation}
the equation holds in distributions, the global Leray--Hopf energy inequality
holds, and the local energy inequality holds for every nonnegative compactly
supported smooth test function.  Here the subscript \(\sigma\) denotes
divergence-free fields.  A terminal point
\((x_0,T)\) is regular when \(u\) is locally bounded in a backward
parabolic neighborhood of that point.  We write
\begin{equation}\label{eq:singular-set}
 \Sigma_u(T):=\{x_0\in\R^3:(x_0,T)\text{ is not regular}\}.
\end{equation}

For \(z_0=(x_0,t_0)\) and \(r>0\), write
\(Q_r(z_0)=B_r(x_0)\times(t_0-r^2,t_0)\) and define
\begin{align}
 A(z_0,r)
 &:=r^{-1}\esssup_{t_0-r^2<t<t_0}
       \int_{B_r(x_0)}|u(x,t)|^2\dd x,\label{eq:A-def}\\
 E(z_0,r)
 &:=r^{-1}\int_{Q_r(z_0)}|\nabla u|^2\dd x\dd t,\label{eq:E-def}\\
 C(z_0,r)
 &:=r^{-2}\int_{Q_r(z_0)}|u|^3\dd x\dd t,\label{eq:C-def}\\
 D(z_0,r)
 &:=r^{-2}\int_{t_0-r^2}^{t_0}\int_{B_r(x_0)}
   |p-(p)_{B_r(x_0)}(t)|^{3/2}\dd x\dd t.\label{eq:D-def}
\end{align}
These quantities are invariant under the Navier--Stokes scaling.  The
pressure mean in \(D\) fixes the local pressure gauge.

\begin{theorem}[Velocity concentration at a singular point]
\label{thm:velocity-concentration}
There is a threshold \(\varepsilon_v=\varepsilon_v(\nu)>0\) such that every
\(x_0\in\Sigma_u(T)\) admits radii \(r_j\downarrow0\) for which
\begin{equation}\label{eq:velocity-concentration}
 C((x_0,T),r_j)\ge\varepsilon_v
 \qquad\text{for every }j.
\end{equation}
\end{theorem}

\begin{proof}
Use the one-scale velocity regularity criterion of
Gustafson--Kang--Tsai \cite{GustafsonKangTsai2007} with spatial and temporal
exponents equal to three.  In the present normalization it gives a constant
\(\varepsilon_{\mathrm{reg}}>0\) such that
\(C(z_0,r)<\varepsilon_{\mathrm{reg}}\) implies boundedness in
\(Q_{r/2}(z_0)\).  The statement is invariant under
\eqref{eq:NS-scaling}.  Therefore a singular point satisfies
\[
 \limsup_{r\downarrow0}C((x_0,T),r)
 \ge\varepsilon_{\mathrm{reg}}.
\]
Take \(\varepsilon_v=\varepsilon_{\mathrm{reg}}/2\) and choose a sequence
converging to zero on which \(C\ge\varepsilon_v\).
\end{proof}

Appendix~\ref{app:velocity-concentration} records the additional pressure
normalization and oscillation-retention statements needed by the later
compactness construction; these do not enter the preceding concentration
lemma.

\subsection{Scale-critical concentration and local kinetic energy}

The target of the variational interpretation is the set of germs retaining
the velocity concentration forced by singularity.

\begin{definition}[Solution-specific singular concentration target]
\label{def:singular-target}
Let \(\varepsilon_v(\nu)>0\) be the threshold in the velocity
\(\varepsilon\)-regularity theorem and put \(\eta_*:=\varepsilon_v/2\).
For a fixed solution \(u\) and terminal time \(T\), let \(\SingT{u}\) be the
set of pairs
\((x_0,\{r_j\}_{j\ge1})\), with \(r_j\downarrow0\), such that
\begin{equation}\label{eq:singular-target}
 C((x_0,T),r_j)\ge\eta_*
 \qquad\text{for every }j.
\end{equation}
Every \(x_0\in\Sigma_u(T)\) occurs in at least one element of
\(\SingT{u}\) by \cref{thm:velocity-concentration}.  Conversely, if \(u\) is bounded in a
backward neighborhood of \((x_0,T)\), then
\(C((x_0,T),r)=O(r^3)\), so no sequence based at \(x_0\) belongs to
\(\SingT{u}\).
\end{definition}

The following estimate explains why localized kinetic probes are sufficient
to measure the energy side of this target.

\begin{lemma}[Local interpolation]
\label{lem:local-interpolation}
For every suitable weak solution and every interior cylinder \(Q_r(z_0)\),
\begin{equation}\label{eq:local-interpolation}
 C(z_0,r)
 \le c\bigl(A(z_0,r)^{3/4}E(z_0,r)^{3/4}
            +A(z_0,r)^{3/2}\bigr),
\end{equation}
where \(c\) is universal.
\end{lemma}

\begin{proof}
After translating the cylinder, the Sobolev and interpolation inequalities on
\(B_r\) give, for almost every time,
\[
 \int_{B_r}|u|^3
 \le c
 \left(\int_{B_r}|u|^2\right)^{3/4}
 \left(\int_{B_r}|\nabla u|^2\right)^{3/4}
 +c r^{-3/2}\left(\int_{B_r}|u|^2\right)^{3/2}.
\]
Integrate over a time interval of length \(r^2\), apply H\"older's
inequality to the first term, and divide by \(r^2\).  Substitution of
\eqref{eq:A-def} and \eqref{eq:E-def} yields
\eqref{eq:local-interpolation}.
\end{proof}

\begin{corollary}[Concentration requires store or dissipation]
\label{cor:store-or-dissipation}
Suppose \(C(z_j,r_j)\ge\eta>0\).  Then
\[
 \max\{A(z_j,r_j),E(z_j,r_j)\}\ge c\eta^{2/3}>0.
\]
If \(E(z_j,r_j)\to0\), then
\(\liminf_j A(z_j,r_j)\ge c\eta^{2/3}\).
\end{corollary}

\begin{proof}
Both assertions follow directly from \eqref{eq:local-interpolation}.
\end{proof}

Thus a singular target retains local kinetic energy, viscous dissipation, or
both.  The canonical orbit constructed in \cref{sec:critical-orbit} sharpens
this alternative by subtracting the weighted velocity mean.  Pressure remains
part of the local regularity criterion and of the conservative flux; viscosity
supplies the dissipative current.

\section{Euler transport and the Stokes gradient}
\label{sec:gradient-flow}

\subsection{The solenoidal Gelfand triple}

Let
\[
 \mathcal V_\sigma:=\dot H^1_\sigma(\R^3),
 \qquad
 \mathcal H_\sigma:=L^2_\sigma(\R^3),
 \qquad
 \mathcal V_\sigma^*:=\dot H^{-1}_\sigma(\R^3).
\]
The Stokes operator \(A_{\mathrm S}:\mathcal V_\sigma\to
\mathcal V_\sigma^*\) is defined weakly by
\begin{equation}\label{eq:stokes-operator}
 \langle A_{\mathrm S}v,\varphi\rangle
 =\int_{\R^3}\nabla v:\nabla\varphi\dd x.
\end{equation}
On smooth fields, \(A_{\mathrm S}=-\Pdiv\Delta\).  Put
\begin{equation}\label{eq:mobility}
 K:=\nu A_{\mathrm S}.
\end{equation}
This is the Riesz isomorphism associated with the Dirichlet form multiplied
by \(\nu\).  For \(g,h\in\mathcal V_\sigma^*\), define
\begin{equation}\label{eq:inverse-metric}
 (g,h)_{K^{-1}}:=\langle g,K^{-1}h\rangle,
 \qquad \|g\|_{K^{-1}}^2=(g,g)_{K^{-1}}.
\end{equation}

The projected convection operator is
\begin{equation}\label{eq:B-def}
 B(v,w):=\Pdiv((v\cdot\nabla)w),
 \qquad
 \langle B(v,w),\varphi\rangle
 =-\int_{\R^3}v\otimes w:\nabla\varphi\dd x.
\end{equation}
For sufficiently regular divergence-free \(v\),
\begin{equation}\label{eq:B-skew}
 \langle B(v,v),v\rangle=0.
\end{equation}

\begin{definition}[Euler-covariant derivative]
\label{def:euler-covariant}
Along a divergence-free curve \(u(t)\), define
\begin{equation}\label{eq:euler-covariant}
 \Dcov_t^{\mathrm E}u:=\partial_t u+B(u,u).
\end{equation}
This is the material derivative after projecting away the pressure
constraint.
\end{definition}

\begin{proposition}[Exact covariant gradient identity]
\label{prop:covariant-gradient}
Let \(u\) be a smooth solution of \eqref{eq:NS}.  For
\(\Energy(u)=\tfrac12\|u\|_2^2\),
\begin{equation}\label{eq:covariant-gradient}
 \Dcov_t^{\mathrm E}u=-K D\Energy(u)=-Ku.
\end{equation}
If \(J_\nu:=-Ku\), then
\begin{equation}\label{eq:action-is-dissipation}
 \|J_\nu\|_{K^{-1}}^2
 =\langle Ku,u\rangle
 =\nu\|\nabla u\|_2^2,
\end{equation}
and
\begin{equation}\label{eq:smooth-energy-identity}
 \frac{\dd}{\dd t}\Energy(u(t))
 =-\|J_\nu(t)\|_{K^{-1}}^2.
\end{equation}
\end{proposition}

\begin{proof}
Apply \(\Pdiv\) to \eqref{eq:NS}.  Since \(\Pdiv\nabla p=0\), this gives
\(\partial_t u+B(u,u)+Ku=0\), which is
\eqref{eq:covariant-gradient}.  The definition of the inverse metric gives
\[
 \|Ku\|_{K^{-1}}^2=\langle Ku,u\rangle
 =\nu\int_{\R^3}|\nabla u|^2\dd x.
\]
Pairing \eqref{eq:covariant-gradient} with \(u\) and using
\eqref{eq:B-skew} proves \eqref{eq:smooth-energy-identity}.
\end{proof}

\subsection{The finite energy inequality}

For a Leray--Hopf solution, choose its standard weakly continuous
representative and set
\begin{equation}\label{eq:energy-balance}
 M(t):=\Energy(u(t))+\nu\int_0^t\|\nabla u(s)\|_2^2\dd s.
\end{equation}
The global energy inequality says that \(M\) is nonincreasing after changing
it on a null set.  Its distributional derivative therefore defines a
nonnegative Radon measure.

\begin{theorem}[Energy dissipation for suitable weak solutions]
\label{thm:weak-energy-balance}
There is a nonnegative measure \(\mu_{\mathrm{en}}\) on \((0,T)\) such that,
at every pair of times \(0\le t_1<t_2<T\) for which the strong energy
inequality is represented,
\begin{equation}\label{eq:weak-energy-balance}
 \Energy(u(t_2))
 +\int_{t_1}^{t_2}\|J_\nu(t)\|_{K^{-1}}^2\dd t
 +\mu_{\mathrm{en}}((t_1,t_2])
 =\Energy(u(t_1)).
\end{equation}
In particular,
\begin{equation}\label{eq:finite-action}
 \int_0^T\|J_\nu(t)\|_{K^{-1}}^2\dd t
 \le\Energy(u(0)).
\end{equation}
For every classical solution,
\(\mu_{\mathrm{en}}=0\) on each compact interval of classical existence.
\end{theorem}

\begin{proof}
Define \(\mu_{\mathrm{en}}:=-D_tM\) in distributions.  Monotonicity of \(M\) gives
\(\mu_{\mathrm{en}}\ge0\), and the increment formula for \(M\) is exactly
\eqref{eq:weak-energy-balance}.  The bound \eqref{eq:finite-action} follows by
discarding the two nonnegative terminal terms.  Smooth solutions satisfy the
energy equality \eqref{eq:smooth-energy-identity}, so their defect measure
vanishes.
\end{proof}

Equation \eqref{eq:weak-energy-balance} is the finite dissipation bound used in
the rest of the paper.  The available quantity is the initial kinetic energy,
and irreversible motion is measured by the Stokes action
\(\|J_\nu\|_{K^{-1}}^2\).

\section{Material deformation and stretching--diffusion duality}
\label{sec:material-duality}

Let \(u\) be smooth on a time interval and let \(X(a,t)\) be its material
flow,
\begin{equation}\label{eq:material-flow}
 \partial_tX(a,t)=u(X(a,t),t),
 \qquad X(a,t_0)=a.
\end{equation}
Write
\begin{equation}\label{eq:material-tensors}
 \MaterialF:=\nabla_aX,
 \qquad
 \MaterialMetric:=\MaterialF^T\MaterialF,
 \qquad
 \MaterialVort(a,t):=\omega(X(a,t),t),
 \qquad \omega:=\nabla\times u.
\end{equation}
The tensor \(\MaterialMetric\) is the pullback of the Euclidean metric by the
volume-preserving map \(X\).  Denote the Eulerian strain by
\(S=\tfrac12(\nabla u+\nabla u^T)\) and put
\(S_L=S\circ X\).

\begin{proposition}[Exact material pullback]
\label{prop:material-pullback}
For every smooth solution and every time on which \eqref{eq:material-flow}
is defined,
\begin{align}
 \det\MaterialF&=1,
 &\partial_t\MaterialF&=(\nabla u\circ X)\MaterialF,
 \label{eq:F-evolution}\\
 \partial_t\MaterialMetric&=2\MaterialF^TS_L\MaterialF,
 &\det\MaterialMetric&=1.                 \label{eq:g-evolution}
\end{align}
For every smooth scalar field \(q\),
\begin{equation}\label{eq:laplacian-pullback}
 (\Delta_xq)\circ X
 =\partial_{a_i}\!\left(\MaterialMetric^{ij}
             \partial_{a_j}(q\circ X)\right),
\end{equation}
where \((\MaterialMetric^{ij})=\MaterialMetric^{-1}\).  The pulled-back
vorticity satisfies
\begin{equation}\label{eq:material-vorticity}
 \partial_t\MaterialVort
 =S_L\MaterialVort
  +\nu\partial_{a_i}\!\left(\MaterialMetric^{ij}
                  \partial_{a_j}\MaterialVort\right).
\end{equation}
Equivalently, the viscous Cauchy invariant
\(\xi:=\MaterialF^{-1}\MaterialVort\) satisfies
\begin{equation}\label{eq:viscous-cauchy-invariant}
 \partial_t\xi
 =\nu\MaterialF^{-1}\partial_{a_i}\!\left(
       \MaterialMetric^{ij}\partial_{a_j}(\MaterialF\xi)\right).
\end{equation}
\end{proposition}

\begin{proof}
Differentiating \eqref{eq:material-flow} in \(a\) gives the evolution of
\(\MaterialF\).  Jacobi's formula and incompressibility give
\[
 \partial_t\det\MaterialF
 = (\diver u)\circ X\,\det\MaterialF=0,
\]
and the initial determinant is one.  Differentiating
\(\MaterialMetric=\MaterialF^T\MaterialF\) proves
\eqref{eq:g-evolution}.  Formula \eqref{eq:laplacian-pullback} is the
Laplace--Beltrami formula with volume density
\(\sqrt{\det\MaterialMetric}=1\).

The vorticity equation is
\[
 (\partial_t+u\cdot\nabla)\omega
 = (\omega\cdot\nabla)u+\nu\Delta\omega.
\]
The antisymmetric part of \(\nabla u\) annihilates \(\omega\), so
\((\omega\cdot\nabla)u=S\omega\).  Composition with \(X\) and
\eqref{eq:laplacian-pullback} give \eqref{eq:material-vorticity}.
Finally,
\(\partial_t\MaterialF^{-1}=-\MaterialF^{-1}(\nabla u\circ X)\);
the inviscid terms cancel after substituting
\(\MaterialVort=\MaterialF\xi\), which proves
\eqref{eq:viscous-cauchy-invariant}.
\end{proof}

\begin{corollary}[Material enstrophy balance]
\label{cor:material-enstrophy}
If the fields decay sufficiently at infinity, then
\begin{equation}\label{eq:material-enstrophy-balance}
 \frac12\frac{\dd}{\dd t}\int_{\R^3}|\MaterialVort|^2\dd a
 +\nu\int_{\R^3}\MaterialMetric^{ij}
       \partial_{a_i}\MaterialVort\cdot
       \partial_{a_j}\MaterialVort\dd a
 =\int_{\R^3}\MaterialVort\cdot S_L\MaterialVort\dd a.
\end{equation}
\end{corollary}

\begin{proof}
Pair \eqref{eq:material-vorticity} with \(\MaterialVort\), integrate in
\(a\), and integrate the divergence term by parts.  The material volume
element is \(\dd a=\dd x\) because \(\det\MaterialF=1\).
\end{proof}

The duality in \eqref{eq:material-vorticity} is algebraic.  If the singular
values of \(\MaterialF\) are
\(\sigma_1\ge\sigma_2\ge\sigma_3>0\), then
\begin{equation}\label{eq:singular-value-duality}
 \sigma_1\sigma_2\sigma_3=1,
 \qquad
 \lambda(\MaterialMetric^{-1})=
 \{\sigma_1^{-2},\sigma_2^{-2},\sigma_3^{-2}\},
 \qquad
 \sigma_3^{-2}\ge\sigma_1.
\end{equation}
Thus strong material stretching produces a strongly diffusive compressed
direction.  This observation alone is not a coercive estimate: the pulled-back
vorticity could avoid variation in that direction.  The equality mechanism is
therefore geometric.  Avoiding the enhanced direction requires simultaneous
alignment of vorticity, strain, and the material eigenframe.  The central
rigidity problem is to show that a target-retaining orbit cannot sustain this
alignment indefinitely unless it approaches a stretching-depleted,
lower-dimensional state.

For the regularity argument, this geometric alignment is tested by the
Eulerian critical production.  Equality means
\(\nabla V=0\) on the retained connected core, and therefore spatial
homogeneity modulo the Galilean quotient; see
\cref{thm:zero-production-rigidity}.  Thus the material duality explains the
sign of the production, while the compact Eulerian equality theorem supplies
the rigorous rigidity statement used in the proof.

\section{Covariance under scale--translation charts}
\label{sec:charts}

\subsection{The repaired equation}

Let \(X(\tau)\in\R^3\) and \(\Lambda(\tau)>0\) be absolutely continuous,
and define
\begin{equation}\label{eq:chart}
 x=X(\tau)+\Lambda(\tau)y,
 \qquad \frac{\dd t}{\dd\tau}=\Lambda(\tau)^2,
\end{equation}
\begin{equation}\label{eq:chart-fields}
 u(x,t)=\Lambda(\tau)^{-1}V(y,\tau),
 \qquad p(x,t)=\Lambda(\tau)^{-2}P(y,\tau).
\end{equation}
Put
\begin{equation}\label{eq:ab-def}
 a(\tau)=\frac{\Lambda'(\tau)}{\Lambda(\tau)},
 \qquad b(\tau)=\frac{X'(\tau)}{\Lambda(\tau)}.
\end{equation}

\begin{proposition}[Covariant equation in a repaired chart]
\label{prop:chart-equation}
The physical equation \eqref{eq:NS} is equivalent, in distributions on the
image of the chart, to
\begin{equation}\label{eq:renormalized-NS}
 \partial_\tau V+(V\cdot\nabla)V+\nabla P
 =\nu\Delta V+a(V+y\cdot\nabla V)+b\cdot\nabla V,
 \qquad \diver V=0.
\end{equation}
Consequently the chart-covariant derivative
\begin{equation}\label{eq:chart-covariant-derivative}
 \Dcov_\tau^{X,\Lambda}V
 :=\partial_\tau V+B(V,V)
   -a(V+y\cdot\nabla V)-b\cdot\nabla V
\end{equation}
satisfies
\begin{equation}\label{eq:chart-gradient-flow}
 \Dcov_\tau^{X,\Lambda}V=-\nu A_{\mathrm S}V.
\end{equation}
\end{proposition}

\begin{proof}
The change of variables in \eqref{eq:chart}--\eqref{eq:chart-fields} gives
the same factor \(\Lambda^{-3}\) in every term of the physical equation.
After removing that factor, the time derivative contributes
\(a(V+y\cdot\nabla V)+b\cdot\nabla V\), giving
\eqref{eq:renormalized-NS}.  Applying the Helmholtz projection gives
\eqref{eq:chart-gradient-flow}.
\end{proof}

For the centered Type~I variables
\[
 t=-e^{-\tau},\qquad \Lambda(\tau)=e^{-\tau/2},\qquad X=0,
\]
one has
\begin{equation}\label{eq:typeI-sign}
 a=-\frac12,\qquad b=0,
\end{equation}
and hence
\begin{equation}\label{eq:typeI-centered-equation}
 \partial_\tau V+(V\cdot\nabla)V+\nabla P
 =\nu\Delta V-\frac12(V+y\cdot\nabla V).
\end{equation}
This is the standard centered Type~I convention.  In a repaired Type~II
chart, the required coefficient regularity is
\begin{equation}\label{eq:typeII-ab-regularity}
 a,b\in L^1_{\mathrm{loc}}(\dd\tau).
\end{equation}
Whenever a compact-window argument requires bounded modulation, that bound is
supplied by the observer construction; failure of the bound is retained as an
observer-current output.

\subsection{Covariant energy and pressure}

On a whole-space chart, physical kinetic energy and action take the form
\begin{align}
 \Energy(u(t(\tau)))
 &=\frac{\Lambda(\tau)}2\int_{\R^3}|V(y,\tau)|^2\dd y,
 \label{eq:chart-energy}\\
 \|J_\nu(t)\|_{K^{-1}}^2\dd t
 &=\Lambda(\tau)\nu\int_{\R^3}|\nabla V(y,\tau)|^2\dd y\dd\tau.
 \label{eq:chart-action}
\end{align}

\begin{proposition}[Energy identity in moving coordinates]
\label{prop:chart-energy}
For a smooth whole-space represented solution,
\begin{equation}\label{eq:chart-energy-identity}
 \frac{\dd}{\dd\tau}
 \left(\frac{\Lambda}2\|V\|_2^2\right)
 =-\Lambda\nu\|\nabla V\|_2^2.
\end{equation}
Thus \eqref{eq:smooth-energy-identity} and
\eqref{eq:chart-energy-identity} are the same scalar identity in different
coordinates.
\end{proposition}

\begin{proof}
Pair \eqref{eq:renormalized-NS} with \(V\).  Convection, translation, and
pressure give zero.  Moreover,
\[
 \int_{\R^3}V\cdot(V+y\cdot\nabla V)\dd y
 =-\frac12\|V\|_2^2.
\]
Since \(\Lambda'=a\Lambda\), the derivative of the prefactor in
\eqref{eq:chart-energy} cancels this dilation term.  The remaining term is
the right-hand side of \eqref{eq:chart-energy-identity}.
\end{proof}

On compact cylinders, derivatives of the cutoff produce convective, pressure,
scale, and translation fluxes in the local energy inequality.  Their exact
form is recorded in Appendix~\ref{app:gauge-pressure}; these are the local
terms placed explicitly in the moving error density
\eqref{eq:moving-cutoff-error}.

The fixed- and moving-cutoff bounds follow from
\cref{prop:pressure-normalization-compactness,prop:covariant-adjoint-localization}
and the suitable local-energy inequality on the normalized windows.

Pressure remains the multiplier for incompressibility.  On an inner ball use
the decomposition
\begin{equation}\label{eq:pressure-decomposition}
 P=P_{\mathrm{loc}}+H,
 \qquad
 -\Delta P_{\mathrm{loc}}
 =\partial_i\partial_j(\zeta V_iV_j),
 \qquad \Delta H=0,
\end{equation}
with a spatial mean fixing \(H\).  The harmonic term can depend on space and
time.

The local Calder\'on--Zygmund estimate and the harmonic interior estimate,
with constants uniform on the normalized compact orbit, are the required
conclusion of \cref{app:local-pressure-decomposition}.

\begin{proposition}[Same-current covariance]
\label{prop:same-current-covariance}
Let two admissible scale--translation charts represent the same physical
solution on a common cylinder.  The fields, test weights, and Stokes pairings
transform by the corresponding pullback, while the physical current
\begin{equation}\label{eq:fixed-current}
 J_\nu=-\nu A_{\mathrm S}u
\end{equation}
is unchanged.  In particular, \eqref{eq:chart-action} is invariant after the
time change.  Optimizing a concentration probe over admissible charts changes
the probe but does not change \(J_\nu\).
\end{proposition}

\begin{proof}
Both charts are obtained from \eqref{eq:chart} by composing translations and
positive dilations.  The tensorial change of variables gives
\eqref{eq:chart-energy}--\eqref{eq:chart-action} in each chart.  Since both
right-hand sides equal the same physical energy and physical dissipation,
their values agree.  The current \eqref{eq:fixed-current} is determined by
\(u\) before either chart is selected.
\end{proof}

\section{Monotone charts and recurrent pressure cancellation}
\label{sec:recurrent-covariant-rigidity}

The scale and translation terms in \eqref{eq:renormalized-NS} have two
different roles.  Their oscillation along a particular normalization is a
choice of coordinates, whereas their limiting value on a recurrent compact
class is part of the equation.  The first result below removes artificial
positive scale motion.  The second identifies the recurrent equality class
without estimating pressure on individual windows.

\subsection{A monotone interpolation through ordered active cylinders}

\begin{proposition}[Monotone scale interpolation]
\label{prop:monotone-scale-interpolation}
Let
\[
 J_j=[s_j,t_j],\qquad
 t_j<s_{j+1},\qquad
 |J_j|=r_j^2,
 \qquad r_{j+1}<r_j,
 \]
be an infinite ordered family of physical time intervals, and let
\(x_j\in\R^3\).  There are locally absolutely continuous functions
\[
 \widehat\Lambda:(s_1,T)\to(0,\infty),
 \qquad
 \widehat X:(s_1,T)\to\R^3,
 \]
where \(T:=\lim_jt_j\), with the following properties:
\begin{enumerate}[label=\textup{(\roman*)},leftmargin=2.2em]
\item \(\widehat\Lambda=r_j\) and \(\widehat X=x_j\) on \(J_j\);
\item \(\widehat\Lambda\) is nonincreasing;
\item after the time change
\begin{equation}\label{eq:monotone-chart-clock}
 \widehat\tau(t):=
 \int_{s_1}^{t}\widehat\Lambda(q)^{-2}\dd q,
\end{equation}
the coefficients of the corresponding chart satisfy
\begin{equation}\label{eq:monotone-chart-coefficients}
 \widehat a:=\frac{\dd}{\dd\widehat\tau}
                 \log\widehat\Lambda\le0
 \quad\hbox{a.e.},
 \qquad
 \widehat b:=\frac{\dd\widehat X/\dd\widehat\tau}
                  {\widehat\Lambda}
 \in L^1_{\mathrm{loc}};
\end{equation}
\item every \(J_j\) is a unit interval in the
\(\widehat\tau\)-clock and \(\widehat\tau(t)\to\infty\) as
\(t\uparrow T\).
\end{enumerate}
On each active interval the represented fields are exactly the fixed-scale
parabolic fields
\begin{equation}\label{eq:monotone-active-fields}
 \widehat V_j(y,\sigma)
 =r_j u(x_j+r_jy,s_j+r_j^2\sigma),
 \qquad 0\le\sigma\le1,
\end{equation}
and therefore satisfy \eqref{eq:renormalized-NS} with
\(\widehat a=\widehat b=0\) there.
\end{proposition}

\begin{proof}
Set \(\widehat\Lambda=r_j\) and \(\widehat X=x_j\) on \(J_j\).
On each gap \([t_j,s_{j+1}]\), interpolate
\(\log\widehat\Lambda\) affinely from \(\log r_j\) to
\(\log r_{j+1}\), and interpolate \(\widehat X\) affinely from
\(x_j\) to \(x_{j+1}\).  The resulting functions are locally absolutely
continuous.  Since \(r_{j+1}<r_j\), the scale is nonincreasing.

The map \(t\mapsto\widehat\tau(t)\) is strictly increasing and locally
absolutely continuous.  Its inverse satisfies
\[
 \frac{\dd t}{\dd\widehat\tau}=\widehat\Lambda^2.
\]
Consequently
\[
 \widehat a
 =\widehat\Lambda\frac{\dd\widehat\Lambda}{\dd t}\le0,
 \qquad
 \widehat b
 =\widehat\Lambda\frac{\dd\widehat X}{\dd t},
\]
which proves \eqref{eq:monotone-chart-coefficients}.  Both coefficients are
locally integrable because the interpolants are absolutely continuous and
the scale is bounded above and below on every compact subinterval of
\((s_1,T)\).

On \(J_j\), the scale and center are constant, so
\[
 \widehat\tau(t)-\widehat\tau(s_j)
 =\frac{t-s_j}{r_j^2}.
\]
Thus \(J_j\) has represented length one and the change of variables is
exactly \eqref{eq:monotone-active-fields}.  Every active interval contributes
one to \eqref{eq:monotone-chart-clock}.  Since the intervals are disjoint and
infinite in number, \(\widehat\tau(t)\to\infty\) as \(t\uparrow T\).
Finally, \cref{prop:chart-equation} gives the represented equation.  On
\(J_j\), both chart derivatives vanish, so \(\widehat a=\widehat b=0\).
\end{proof}

\begin{remark}[What monotone interpolation removes]
\label{rem:monotone-interpolation-scope}
Once the singular extraction has produced chronologically ordered active
cylinders with decreasing scales, positive scale drift is unnecessary.  It
arises only if one insists on continuing a moment normalization through the
inactive gaps.  The monotone chart preserves the normalized fields on every
active cylinder and represents the same physical solution in the gaps, but
it does not by itself prove compactness of those gap segments or persistence
of the target away from the active cylinders.  Those are concentration--
compactness questions, not modulation estimates.
\end{remark}

\subsection{The invariant mean identity}

Fix constants \(a_*\in\R\) and \(b_*\in\R^3\), and consider entire smooth
solutions of
\begin{equation}\label{eq:autonomous-covariant-NS}
 \partial_\tau V+(V\cdot\nabla)V+\nabla P
 =\nu\Delta V+a_*(V+y\cdot\nabla V)+b_*\cdot\nabla V,
 \qquad \diver V=0.
\end{equation}
For \(c\in\R^3\) and \(s\in\R\), define
\begin{align}
 (T_cV)(y,\tau)&:=V(y+e^{-a_*\tau}c,\tau),
 & (T_cP)(y,\tau)&:=P(y+e^{-a_*\tau}c,\tau),
 \label{eq:autonomous-covariant-translation}\\
 (\Theta_sV)(y,\tau)&:=V(y,\tau+s),
 & (\Theta_sP)(y,\tau)&:=P(y,\tau+s).
 \label{eq:autonomous-time-translation}
\end{align}
Direct substitution shows that both maps preserve
\eqref{eq:autonomous-covariant-NS} and that
\begin{equation}\label{eq:autonomous-semidirect-law}
 \Theta_sT_c\Theta_{-s}=T_{e^{-a_*s}c}.
\end{equation}
They therefore generate the solvable, and hence amenable, group
\(G_{a_*}=\R^3\rtimes\R\).

\begin{lemma}[Pressure does no work outside the spatial zero mode]
\label{lem:recurrent-pressure-cancellation}
Let \(\mathcal M\) be a compact \(G_{a_*}\)-invariant family of bounded
smooth entire solutions of \eqref{eq:autonomous-covariant-NS}, compact in a
local topology in which the velocity and pressure-gradient evaluations below
are continuous.  Let \(\mu\) be a \(G_{a_*}\)-invariant Borel probability
measure on \(\mathcal M\).  In \(L^2(\mathcal M,\mu)\), put
\[
 u_i(V):=V_i(0,0),
 \qquad p_i(V):=\partial_iP(0,0),
\]
and let \(\Pi_0\) be the orthogonal projection onto the subspace invariant
under all \(T_c\).  If \(u_i^0:=\Pi_0u_i\) and
\(p_i^0:=\Pi_0p_i\), then
\begin{equation}\label{eq:recurrent-pressure-work}
 \int_{\mathcal M}u_i p_i\dd\mu
 =\int_{\mathcal M}u_i^0p_i^0\dd\mu.
\end{equation}
\end{lemma}

\begin{proof}
Let \(D_k\) be the skew-adjoint generator of the unitary action induced by
\(T_{he_k}\).  Since the translation in
\eqref{eq:autonomous-covariant-translation} equals the ordinary translation
by \(he_k\) at \(\tau=0\),
\[
 D_ku_i=\partial_kV_i(0,0),
 \qquad D_iu_i=0.
\]
Taking the divergence of \eqref{eq:autonomous-covariant-NS} gives
\begin{equation}\label{eq:recurrent-pressure-poisson}
 -\Delta P=\partial_i\partial_j(V_iV_j).
\end{equation}
The scale and translation terms have zero divergence.  Hence, with
\(q_i:=(I-\Pi_0)p_i\) and
\(A:=-\sum_kD_k^2\), the commuting spectral resolution of the spatial
translation action gives
\[
 q_k
 =D_kA^{-1}(I-\Pi_0)D_iD_j(u_iu_j)
\]
on the nonzero spatial spectral subspace.  More precisely, the right side is
the \(L^2\)-limit as \(\varepsilon\downarrow0\) of
\[
 q_{k,\varepsilon}
 :=D_k(\varepsilon+A)^{-1}(I-\Pi_0)D_iD_j(u_iu_j).
\]
The pressure curl relation determines the vector from its divergence and
excludes an additional nonzero-frequency harmonic part.  Skew-adjointness
and incompressibility give
\[
 \int_{\mathcal M}u_kq_{k,\varepsilon}\dd\mu
 =-\int_{\mathcal M}(D_ku_k)
    (\varepsilon+A)^{-1}(I-\Pi_0)D_iD_j(u_iu_j)\dd\mu
 =0.
\]
Pass to the \(L^2\)-limit.  Since
\(u_i-u_i^0\perp\operatorname{ran}\Pi_0\), one also has
\(\int(u_i-u_i^0)p_i^0\dd\mu=0\).  These two identities prove
\eqref{eq:recurrent-pressure-work}.  Only \(\nabla P\) occurs, so the
conclusion is independent of the additive pressure gauge.
\end{proof}

\begin{lemma}[Pressure work under activity-biased rooting]
\label{lem:rooted-pressure-current-identity}
Assume the hypotheses and notation of
\cref{lem:recurrent-pressure-cancellation}.  On the nonzero spatial spectral
subspace write
\begin{equation}\label{eq:pressure-potential-on-hull}
 q_k:=(I-\Pi_0)p_k=D_kh,
 \qquad
 h:=A^{-1}(I-\Pi_0)D_iD_j(u_iu_j),
 \qquad A=-\sum_kD_k^2,
\end{equation}
with the identity understood through the resolvent approximation used in the
proof of \cref{lem:recurrent-pressure-cancellation}.  Let
\(w\ge0\) be bounded, belong to the form domains of the spatial generators,
and satisfy
\begin{equation}\label{eq:root-weight-normalization}
 Z_w:=\int_{\mathcal M}w\dd\mu>0,
 \qquad
 h_\varepsilon u_kD_kw\longrightarrow h u_kD_kw
 \quad\hbox{in }L^1(\mu),
\end{equation}
For the rooted probability \(\dd\mu_w=Z_w^{-1}w\dd\mu\),
\begin{equation}\label{eq:rooted-pressure-current-identity}
 \sum_k\int_{\mathcal M}u_kq_k\dd\mu_w
 =-\frac1{Z_w}\int_{\mathcal M}h\,u_kD_kw\dd\mu.
\end{equation}
Thus activity-biased rooting cancels the nonzero pressure modes precisely when
the root-current pairing on the right of
\eqref{eq:rooted-pressure-current-identity} vanishes.
\end{lemma}

\begin{proof}
Use first the regularized potentials
\[
 h_\varepsilon
 :=(\varepsilon+A)^{-1}(I-\Pi_0)D_iD_j(u_iu_j),
 \qquad q_{k,\varepsilon}=D_kh_\varepsilon.
\]
Skew-adjointness of \(D_k\), the generator product rule, and
incompressibility give
\begin{align*}
 \sum_k\int wu_kD_kh_\varepsilon\dd\mu
 &=-\sum_k\int h_\varepsilon D_k(wu_k)\dd\mu\\
 &=-\sum_k\int h_\varepsilon u_kD_kw\dd\mu
   -\int h_\varepsilon w\sum_kD_ku_k\dd\mu\\
 &=-\sum_k\int h_\varepsilon u_kD_kw\dd\mu.
\end{align*}
The resolvent convergence in
\cref{lem:recurrent-pressure-cancellation} gives
\(q_{k,\varepsilon}\to q_k\) in \(L^2(\mu)\).  Approximation of \(w\) in
the common generator form domain passes the left side to its limit, while
\eqref{eq:root-weight-normalization} passes the right side to its stated
\(L^1\) limit.  Divide by \(Z_w\) to obtain
\eqref{eq:rooted-pressure-current-identity}.
\end{proof}

\begin{remark}[The retained root current]
\label{rem:retained-root-current}
The right side of \eqref{eq:rooted-pressure-current-identity} is absent from
the unweighted invariant calculation and need not have a sign.  In an
event-indexed spacetime decomposition it is represented by pressure and
transport across the moving boundary of the rooted cell.  A nonzero value is
therefore a transport or boundary coordinate in
\eqref{eq:five-PDE-source-families}; setting it to zero without proving the
pairing or the corresponding oriented boundary trace vanishes would lose a
possible successor.
\end{remark}

\begin{theorem}[Rigidity of a recurrent nonexpanding covariant class]
\label{thm:recurrent-nonexpanding-rigidity}
Let \(\mathcal M\) be a nonempty compact minimal
\(G_{a_*}\)-invariant family satisfying the hypotheses of
\cref{lem:recurrent-pressure-cancellation}.  If \(a_*\le0\), every element
of \(\mathcal M\) is spatially constant.  Equivalently, every continuous
Galilean-invariant local oscillation observable vanishes on \(\mathcal M\).
\end{theorem}

\begin{proof}
Amenability of \(G_{a_*}\) gives an invariant probability measure \(\mu\)
on \(\mathcal M\).  Its support is a nonempty closed invariant subset, so
minimality makes \(\supp\mu=\mathcal M\).

The spatially invariant subspace is preserved by time translation because
the spatial subgroup is normal by \eqref{eq:autonomous-semidirect-law}; hence
\(\Pi_0\) commutes with the time-translation group.  Evaluate
\eqref{eq:autonomous-covariant-NS} at \((0,0)\), move every term to the
left, and project onto that subspace.  Spatial derivatives, including the
nonlinearity written as \(D_j(u_ju_i)\), vanish after projection.  The term
\(y\cdot\nabla V\) vanishes at \(y=0\).  Thus
\begin{equation}\label{eq:recurrent-zero-mode-equation}
 \partial_\tau u_i^0-a_*u_i^0+p_i^0=0
 \quad\hbox{in }L^2(\mathcal M,\mu).
\end{equation}
Pairing with \(u_i^0\), summing in \(i\), and using invariance under
\(\Theta_s\) gives
\begin{equation}\label{eq:recurrent-zero-mode-pressure}
 \int_{\mathcal M}u_i^0p_i^0\dd\mu
 =a_*\int_{\mathcal M}|u^0|^2\dd\mu.
\end{equation}

Next take the scalar product of the equation at \((0,0)\) with \(u\) and
average.  The time derivative has zero mean.  Integration by parts with the
spatial generators gives
\[
 -\nu\int u\cdot\Delta V(0,0)\dd\mu
 =\nu\int|\nabla V(0,0)|^2\dd\mu.
\]
The nonlinear and constant-translation terms are spatial divergences and
have zero mean.  Using
\cref{lem:recurrent-pressure-cancellation,eq:recurrent-zero-mode-pressure}
for the pressure term yields
\begin{equation}\label{eq:recurrent-covariant-energy-identity}
 \nu\int_{\mathcal M}|\nabla V(0,0)|^2\dd\mu
 -a_*\int_{\mathcal M}|u-u^0|^2\dd\mu=0.
\end{equation}
Here we used that multiplication by the spatially invariant constant
\(a_*\) commutes with \(\Pi_0\).  Both terms in
\eqref{eq:recurrent-covariant-energy-identity} are nonnegative when
\(a_*\le0\); hence the first vanishes.  Full support of \(\mu\) and
continuity of the gradient evaluation imply
\(\nabla V(0,0)=0\) for every \(V\in\mathcal M\).  The group action moves
\((0,0)\) to every point of spacetime, so \(\nabla V\equiv0\) for every
element of \(\mathcal M\).  Such velocities are spatially constant, and
every mean-subtracted local oscillation is zero.
\end{proof}

The constant-coefficient formulation is convenient for minimal autonomous
hulls, but a monotone scale interpolation naturally produces a stationary
law of coefficient paths rather than a single constant.  The same identity
holds at that level.

\begin{definition}[Stationary monotone covariant ensemble]
\label{def:stationary-monotone-ensemble}
A stationary monotone covariant ensemble is a Borel probability space
\((\mathcal M,\mu)\) of entire smooth represented solutions
\((V,P,a,b)\) of \eqref{eq:renormalized-NS} with the following properties.
\begin{enumerate}[label=\textup{(\roman*)},leftmargin=2.2em]
\item The local velocity, pressure-gradient, and coefficient observables used
below are integrable, and the velocity observables are in the domains of the
spatial and time generators.  This holds, for example, on a compact locally
smooth class with a uniform local coefficient bound.
\item The law is invariant under time translation of the four coefficient
paths.
\item For \(c\in\R^3\), let \(q_c\) solve
\begin{equation}\label{eq:variable-covariant-translation-path}
 q_c'(\tau)=-a(\tau)q_c(\tau),
 \qquad q_c(0)=c,
\end{equation}
and set
\[
 T_c(V,P,a,b)
 :=\bigl(V(\,\cdot+q_c(\tau),\tau),
         P(\,\cdot+q_c(\tau),\tau),a,b\bigr).
\]
The class and the law are invariant under every \(T_c\).
\item The scale coefficient satisfies \(a\le0\) almost everywhere for
\(\mu\)-almost every path.
\end{enumerate}
The integrability in item~\textup{(i)} includes
\[
 \int_{\mathcal M}
 \bigl(|\nabla V(0,0)|^2+|a(0)|\,|V(0,0)|^2\bigr)\dd\mu<\infty,
\]
with time-zero values understood through the stationary measurable
representative.  Equivalently, one may formulate every identity below after
Steklov averaging in time and then pass to almost every time.
\end{definition}

\begin{theorem}[Rigidity of a stationary monotone covariant ensemble]
\label{thm:stationary-monotone-ensemble-rigidity}
Let \((\mathcal M,\mu)\) be a stationary monotone covariant ensemble.  Let
\(u_i=V_i(0,0)\), and let \(u_i^0\) be its orthogonal projection in
\(L^2(\mathcal M,\mu)\) onto the spatially invariant observables.  Then
\begin{equation}\label{eq:variable-recurrent-energy-identity}
 \nu\int_{\mathcal M}|\nabla V(0,0)|^2\dd\mu
 +\int_{\mathcal M}(-a(0))|u-u^0|^2\dd\mu=0.
\end{equation}
Consequently \(\nabla V=0\) at every spacetime point for
\(\mu\)-almost every represented solution.  In particular, every
mean-subtracted local oscillation observable vanishes \(\mu\)-almost surely.
\end{theorem}

\begin{proof}
Equation \eqref{eq:variable-covariant-translation-path} is precisely the
condition that makes \(T_c\) preserve \eqref{eq:renormalized-NS}.  At
\(\tau=0\), its infinitesimal generators are the ordinary spatial
derivatives.  The pressure Poisson equation remains
\eqref{eq:recurrent-pressure-poisson}, because both modulation fields are
divergence free.  The spectral proof of
\cref{lem:recurrent-pressure-cancellation} therefore gives
\begin{equation}\label{eq:variable-pressure-cancellation}
 \int u_ip_i\dd\mu=\int u_i^0p_i^0\dd\mu.
\end{equation}

Although time translation no longer conjugates \(T_c\) by a constant
parameter, it preserves every spatial orbit.  Indeed,
\[
 \Theta_sT_c(V,P,a,b)
 =T_{c_s}\Theta_s(V,P,a,b),
 \qquad
 c_s=c\exp\!\left(-\int_0^sa(q)\dd q\right).
\]
Thus the spatially invariant sigma-algebra is time invariant, and its
orthogonal projection commutes with time translation.  The functions
\(a(0)\) and \(b(0)\) are themselves spatially invariant.

Projecting the equation at \((0,0)\) onto the spatial zero mode now gives
\[
 \partial_\tau u_i^0-a(0)u_i^0+p_i^0=0.
\]
Pair with \(u_i^0\), sum in \(i\), and use stationarity to obtain
\begin{equation}\label{eq:variable-zero-mode-pressure}
 \int u_i^0p_i^0\dd\mu
 =\int a(0)|u^0|^2\dd\mu.
\end{equation}
Testing the full equation at the origin against \(u\), averaging, and using
spatial and time invariance removes the time derivative, convection, and the
\(b\)-transport term.  Diffusion gives
\(\nu\int|\nabla V(0,0)|^2\dd\mu\).  Combining
\eqref{eq:variable-pressure-cancellation} and
\eqref{eq:variable-zero-mode-pressure} gives
\[
 \nu\int|\nabla V(0,0)|^2\dd\mu
 -\int a(0)\bigl(|u|^2-|u^0|^2\bigr)\dd\mu=0.
\]
Multiplication by the spatially invariant observable \(a(0)\) preserves the
zero-mode subspace.  Orthogonality therefore gives
\[
 \int a(0)\bigl(|u|^2-|u^0|^2\bigr)\dd\mu
 =\int a(0)|u-u^0|^2\dd\mu,
\]
which proves \eqref{eq:variable-recurrent-energy-identity}.  Both terms are
nonnegative.  Hence the gradient vanishes almost surely at the origin.
Time and covariant spatial invariance propagate this conclusion to every
spacetime point.  The Steklov formulation mentioned in
\cref{def:stationary-monotone-ensemble} justifies the same calculation when
the coefficient paths are merely locally integrable.
\end{proof}

\begin{corollary}[No target-retaining recurrent limit]
\label{cor:no-target-recurrent-nonexpanding-limit}
Let \(\mathcal K\) be a compact invariant class of autonomous limits with
\(a_*\le0\), and let \(\mathcal O\) be a continuous nonnegative observable
on \(\mathcal K\) which vanishes on spatially constant velocities.  If every
minimal invariant component of \(\mathcal K\) retains
\(\mathcal O\ge\eta>0\), then \(\mathcal K\) is empty.
\end{corollary}

\begin{proof}
Every nonempty compact invariant class contains a compact minimal invariant
component.  Apply \cref{thm:recurrent-nonexpanding-rigidity} on that
component.  It gives \(\mathcal O=0\), contrary to the retained lower bound.
\end{proof}

The preceding rigidity statement becomes quantitative on a compact successor
class through the following elementary ergodic lemma.  This is the point at
which recurrence replaces a list of increasingly sharp estimates.

\begin{lemma}[Uniform block deficit or a saturation law]
\label{lem:uniform-block-deficit}
Let \(K\) be a compact metric space, let \(S:K\to K\) be continuous, and
let \(h\in C(K)\).  Write \(\mathcal P_S(K)\) for the set of
\(S\)-invariant Borel probability measures and put
\begin{equation}\label{eq:maximal-invariant-deficit-average}
 \beta:=\max_{\mu\in\mathcal P_S(K)}\int_Kh\dd\mu.
\end{equation}
Then exactly one of the following alternatives holds.
\begin{enumerate}[label=\textup{(\roman*)},leftmargin=2.2em]
\item \(\beta<0\), and there exist \(N\ge1\) and \(\delta>0\) such that
\begin{equation}\label{eq:uniform-negative-block}
 \sum_{k=0}^{N-1}h(S^kw)\le-\delta
 \qquad(w\in K);
\end{equation}
\item \(\beta\ge0\), and an invariant probability measure \(\mu_*\)
attains \(\beta\).
\end{enumerate}
In particular, if a PDE balance gives
\(\int h\dd\mu\le0\) for every invariant measure and every equality law
is target-null, while a continuous target observable satisfies
\(\mathcal O\ge\eta>0\) on \(K\), then alternative~\textup{(i)} holds.
\end{lemma}

\begin{proof}
The set \(\mathcal P_S(K)\) is nonempty and weak-* compact, and the map
\(\mu\mapsto\int h\dd\mu\) is continuous.  Hence the maximum in
\eqref{eq:maximal-invariant-deficit-average} is attained.

Assume \(\beta<0\) and that \eqref{eq:uniform-negative-block} fails for
every \(N\).  Then for each \(n\) there is \(w_n\in K\) such that
\[
 \sum_{k=0}^{n-1}h(S^kw_n)>-n^{-1}.
\]
For the empirical measures
\[
 \mu_n:=\frac1n\sum_{k=0}^{n-1}\delta_{S^kw_n}
\]
one consequently has \(\liminf_n\int h\dd\mu_n\ge0\).  Pass to a
weak-* convergent subsequence.  For every \(f\in C(K)\),
\[
 \int(f\circ S-f)\dd\mu_n
 =\frac{f(S^nw_n)-f(w_n)}n\longrightarrow0,
\]
so the limit \(\mu\) is \(S\)-invariant.  It satisfies
\(\int h\dd\mu\ge0\), contradicting \(\beta<0\).  Thus a block length
\(N\) exists; compactness of \(K\) makes the negative maximum of the
continuous block sum a number \(-\delta<0\), proving
\eqref{eq:uniform-negative-block}.

For the final assertion, an invariant measure with zero average would be
supported on the target-null equality class by the stated PDE rigidity, while
\(\mathcal O\ge\eta\) on all of \(K\).  Hence equality is impossible and
the maximal invariant average is strictly negative.
\end{proof}

\begin{corollary}[Recurrent equality is the only zero-gap obstruction]
\label{cor:recurrent-equality-only-obstruction}
Let a compact target-retaining successor class have a continuous signed
local-energy excess \(h\), and suppose its finite block sums are the exact
aggregate inward feeding minus viscous and suitable-defect production, up to
bounded endpoint storage and true exterior input of sublinear block growth.
If every invariant saturation law is realized either as a stationary
monotone covariant ensemble in the sense of
\cref{def:stationary-monotone-ensemble} or as a zero-Dirichlet ordinary
window, then the class has a uniform strictly dissipative block deficit.
\end{corollary}

\begin{proof}
Signed additivity removes the internal currents before the invariant average
is taken.  The endpoint storage and exterior input vanish after division by
the block length, so every invariant law has nonpositive mean excess.  In the
equality case, \cref{thm:stationary-monotone-ensemble-rigidity} makes the
first realization spatially constant, while
\cref{thm:zero-production-rigidity} makes the second oscillation-removable.
Both contradict the retained target floor.  Apply
\cref{lem:uniform-block-deficit}.
\end{proof}

\begin{remark}[Interface with noncompact extraction]
\label{rem:remaining-recurrent-extraction}
The theorem closes an \emph{ordinary compact recurrent} scale-collapse or
scale-rigid limit, including the limiting case \(a_*=0\).  Rough, exterior,
diffuse, and noncompact sequences are retained at their first failed
coordinate by \cref{thm:persistent-current-compactification}; repeated active
recentering is treated in the marked path space and does not require a
covariant hull retaining a fixed root.
\end{remark}

\section{Scale-critical quotient geometry}
\label{sec:quotient-geometry}

The physical energy identity and the scale-critical concentration problem
live at different levels.  This section separates them and fixes the quotient
structure used by the critical orbit.

\subsection{Physical action and critical action}

The Navier--Stokes scaling based at \((x_0,t_0)\) is
\begin{equation}\label{eq:NS-scaling}
 u^{(\lambda)}(y,s)
 =\lambda u(x_0+\lambda y,t_0+\lambda^2s),
 \qquad
 p^{(\lambda)}(y,s)
 =\lambda^2p(x_0+\lambda y,t_0+\lambda^2s).
\end{equation}
The quantities \(C,D\) are invariant under \eqref{eq:NS-scaling}.  The
localized Dirichlet quantity is invariant after the parabolic factor is
included:
\begin{equation}\label{eq:critical-dirichlet-scaling}
 r^{-1}\int_{Q_r(z_0)}|\nabla u|^2\dd x\dd t
 =\int_{Q_1}|\nabla u^{(r)}|^2\dd y\dd s.
\end{equation}
This is the dissipation seen by a normalized unit window.

By contrast, the physical Stokes action transforms according to
\eqref{eq:chart-action} and carries the factor \(\Lambda\):
\begin{equation}\label{eq:physical-versus-critical-action}
 \int_{t(\tau_1)}^{t(\tau_2)}
 \nu\|\nabla u(t)\|_2^2\dd t
 =\int_{\tau_1}^{\tau_2}\Lambda(\tau)
 \nu\|\nabla V(\tau)\|_2^2\dd\tau.
\end{equation}
Consequently, finite physical action does not control
\begin{equation}\label{eq:renormalized-critical-action}
 \int_{\tau_1}^{\tau_2}\nu
 \int_{B_R}|\nabla V|^2\dd y\dd\tau
\end{equation}
uniformly along a collapsing orbit.  For example, in centered Type~I
coordinates \(\Lambda=e^{-\tau/2}\).  A represented orbit with order-one
Dirichlet dissipation on every unit interval has finite contribution to
\eqref{eq:physical-versus-critical-action} but infinite contribution to
\eqref{eq:renormalized-critical-action}.  The global energy inequality alone
therefore cannot supply the required contradiction.

The corrected moving-corrected localized energy is designed precisely at the level of
\eqref{eq:renormalized-critical-action}.  Its lower bound replaces the missing
unweighted global bound in represented time.  The production gap then says
that a retained core must spend a fixed amount of this critical action, or of
the other covariant production terms, on every target-retaining window.

\subsection{Galilean and pressure quotients}

For a constant vector \(c\in\R^3\), the Galilean transform
\begin{equation}\label{eq:galilean-transform}
 u^c(x,t):=u(x-ct,t)+c,
 \qquad
 p^c(x,t):=p(x-ct,t)
\end{equation}
preserves \eqref{eq:NS}.  Addition of a function \(\pi(t)\) to the pressure
also preserves the equation.  These modes have no bearing on regularity but
can change an uncentered local kinetic integral.  The variance
\eqref{eq:oscillatory-probe} is invariant under the induced constant-velocity
action:
\begin{equation}\label{eq:variance-galilean-invariance}
 \frac12\int m|(V+c)-\overline{(V+c)}_m|^2
 =\frac12\int m|V-\bar V_m|^2.
\end{equation}
For \(c\ne0\), \eqref{eq:galilean-transform} does not preserve the global
finite-energy class on \(\R^3\).  It is used only as a symmetry of local
profiles and observables; the physical finite-energy solution and its Stokes
current are never replaced by \(u^c\).
Likewise, \(D\) and the pressure atlas depend only on the pressure class
modulo functions of time.

Fix nonzero nonnegative radial cutoffs
\(\zeta_0,\zeta_1\in C_c^\infty(B_R)\), a level \(\Theta_0>0\), and write
\(V_{\lambda,q}(y)=\lambda V(q+\lambda y)\).  For a nonzero cutoff \(\xi\),
set
\begin{equation}\label{eq:fixed-cutoff-mean}
 \langle W\rangle_\xi
 :=\frac{\int\xi W\dd y}{\int\xi\dd y},
 \qquad W_\xi^\circ:=W-\langle W\rangle_\xi.
\end{equation}
Translation and dilation are fixed by the four-component, Galilean-invariant
gauge map
\begin{equation}\label{eq:moment-map}
 \begin{aligned}
 \mathfrak F_0(\lambda,q;V)
   &:=\int\zeta_0
       |(V_{\lambda,q})_{\zeta_0}^\circ|^3\dd y-\Theta_0,\\
 \mathfrak F_j(\lambda,q;V)
   &:=\int y_j\zeta_1
       |(V_{\lambda,q})_{\zeta_1}^\circ|^2\dd y,
   \qquad 1\le j\le3.
 \end{aligned}
\end{equation}
Indeed, replacing \(V\) by \(V+c\) adds the constant \(\lambda c\) to
\(V_{\lambda,q}\), which is removed by both means in
\eqref{eq:fixed-cutoff-mean}.  Thus the gauge selects only scale and spatial
center; it does not select a velocity representative.
At a zero for which
\begin{equation}\label{eq:gauge-jacobian}
 D_{(\lambda,q)}\mathfrak F(1,0;V)\in\R^{4\times4}
 \quad\text{is invertible},
\end{equation}
the implicit-function theorem selects the scale and center locally.
Differentiating the four constraints produces \((a,b)\).  Degeneration of
\eqref{eq:gauge-jacobian} is tested during canonical-observer construction;
no smooth continuation through a degenerate point is assumed.

\subsection{The quotient topology}

On a fixed represented cylinder \(B_R\times I\), put
\begin{equation}\label{eq:window-state-space}
 \mathcal X_R(I)
 :=\bigl(L^\infty_I L^2(B_R)\cap L^2_IH^1(B_R)
          \cap L^3(B_R\times I)\bigr)
   \times L^{3/2}(B_R\times I)
   \times L^1(I;\R^4),
\end{equation}
where the factors represent \(V\), \(P\), and \((a,b)\).  The compactness
topology consists of strong local \(L^3\) convergence for velocity, weak local
\(L^2H^1\) convergence for its gradient, strong local
\(L^{3/2}\) convergence for the normalized pressure, and \(L^1\) convergence
of modulation coefficients on the selected windows.  Adjoint densities and
cutoffs are added with the topology in
\eqref{eq:adjoint-compactness-hypotheses} and
\eqref{eq:weight-convergence}.

Let \(\mathcal G_{\mathrm{ch}}\) denote the local chart action generated by
translations, positive dilations, and pressure gauge.  Near a zero satisfying
\eqref{eq:gauge-jacobian}, the normalized slice
\eqref{eq:canonical-moments}, together with
\eqref{eq:canonical-pressure-gauge}, intersects the corresponding
\(\mathcal G_{\mathrm{ch}}\)-orbit locally and transversely.  Global
continuation and consistency of these local slices are part of
\cref{thm:canonical-critical-orbit}.  Constant-velocity Galilean symmetry
is not fixed by a time-dependent mean condition; instead, every scalar used
below is invariant under that symmetry by construction.  A compact family of
normalized windows may therefore be studied on the chart slice and then
quotiented by the residual Galilean action.

\begin{proposition}[Covariance of the reduced data]
\label{prop:reduced-data-covariance}
On the nondegenerate normalized slice, the following objects depend only on
the quotient state:
\begin{enumerate}[label=\textup{(\roman*)},leftmargin=2.2em]
\item the oscillatory kinetic energy \(\OscEnergy_\Gamma\) and concentration
\(\OscMass_\Gamma\);
\item the physical Stokes current and its inverse-metric norm;
\item the ray distance after the target probe is pulled back with the chart;
\item the normalized pressure class and, for a fixed padded terminal problem,
the localized adjoint probability measure \(m\dd y\).
\end{enumerate}
\end{proposition}

\begin{proof}
The first item follows from change of variables and
\eqref{eq:variance-galilean-invariance}.  The second and third follow from
\cref{prop:same-current-covariance} and covariance of orthogonal projection in
the Stokes metric.  Pressure means are removed by
\eqref{eq:canonical-pressure-gauge}.  For a fixed padded terminal problem, the
normalized adjoint equation pulls back covariantly when its domain, cutoff,
and terminal datum are pulled back with the chart, and normalization fixes the
mass.  No comparison between distinct terminal problems is asserted or needed
for the moving-corrected localized energy.
\end{proof}

Compactness and coherent continuation of the observer are supplied by
\cref{thm:canonical-critical-orbit}.  Moving-cutoff errors are retained in
the signed current rather than imposed as a separate global condition.

\section{Backward adjoint localizations and entropy}
\label{sec:adjoint}

The adjoint density supplies scale-relative localization.  Its entropy is
kept separate from the kinetic-energy gradient identity; this separation
prevents pressure and critical norms from acquiring artificial variational
equations.

\subsection{The normalized finite-horizon adjoint problem}

Fix a repaired cylinder \(B_R\times(\tau_0,\sigma)\), let
\begin{equation}\label{eq:transport-drift}
 c(y,\tau):=V(y,\tau)-a(\tau)y-b(\tau),
 \qquad \diver c=-3a,
\end{equation}
and choose \(\chi\in C_c^\infty(B_R)\), \(\chi\ge0\), together with a
nonnegative terminal datum \(h\in C_c^\infty(B_R)\), \(h\not\equiv0\).
First solve the linear terminal-boundary problem
\begin{equation}\label{eq:unnormalized-adjoint}
 -\partial_\tau\widetilde\rho-\nu\Delta\widetilde\rho
 -\diver(c\widetilde\rho)=0,
 \qquad
 \widetilde\rho|_{\partial B_R}=0,
 \qquad
 \widetilde\rho(\sigma)=h.
\end{equation}
Define
\begin{equation}\label{eq:adjoint-normalization}
 Z(\tau):=\int_{B_R}\chi^2\widetilde\rho\dd y,
 \qquad
 \rho:=\frac{\widetilde\rho}{Z},
 \qquad
 m:=\chi^2\rho.
\end{equation}
We work on intervals on which \(Z>0\).  Then \(m\ge0\) and
\begin{equation}\label{eq:m-normalized}
 \int_{B_R}m(y,\tau)\dd y=1.
\end{equation}

For a retained unit window \(I=[s,s+1]\), the adjoint problem is always
solved on a padded interval ending at
\begin{equation}\label{eq:adjoint-buffer}
 \sigma_I=s+1+\kappa_*,
 \qquad \kappa_*>0.
\end{equation}
Consequently
\begin{equation}\label{eq:theta-buffer}
 \kappa_*e^{-2M_*}\le\theta_I(\tau)
 \le(1+\kappa_*)e^{2M_*}
 \qquad(\tau\in I).
\end{equation}
Here \(M_*\) is a uniform bound for
\(\int_s^{\sigma_I}|a(r)|\dd r\) on the normalized padded windows, and
\(\theta_I\) is the unique solution of the covariant terminal-value problem
\begin{equation}\label{eq:covariant-heat-scale}
 \theta_I'+2a\theta_I=-1,
 \qquad \theta_I(\sigma_I)=0.
\end{equation}
Indeed,
\begin{equation}\label{eq:covariant-heat-scale-formula}
 \theta_I(\tau)
 =\int_\tau^{\sigma_I}
   \exp\!\left(2\int_\tau^r a(q)\dd q\right)\dd r,
\end{equation}
which proves \eqref{eq:theta-buffer}.  The Gaussian entropy is evaluated only
on the retained window, where \(\theta_I>0\), not at the terminal time.  It is
a covariant diagnostic on each padded window; the global monotonicity balance
constructed below does not compare or reset finite-horizon adjoint densities.

\begin{lemma}[Normalized adjoint equation]
\label{lem:normalized-adjoint}
Set
\begin{equation}\label{eq:kappa}
 \kappa(\tau):=\frac{Z'(\tau)}{Z(\tau)}
 =\int_{B_R}
 \bigl(-\nu\Delta\chi^2+c\cdot\nabla\chi^2\bigr)\rho\dd y.
\end{equation}
Then
\begin{equation}\label{eq:normalized-adjoint}
 -\partial_\tau\rho-\nu\Delta\rho-\diver(c\rho)
 =\kappa\rho
\end{equation}
in distributions, and \eqref{eq:m-normalized} holds at every time at which
the weak solution is represented continuously in \(L^2\).
\end{lemma}

\begin{proof}
Differentiate \(Z\) using \eqref{eq:unnormalized-adjoint} and integrate by
parts.  The support of \(\chi\) is contained in \(B_R\), so
\[
 Z'=\int_{B_R}
 \bigl(-\nu\Delta\chi^2+c\cdot\nabla\chi^2\bigr)
 \widetilde\rho\dd y.
\]
Division by \(Z\) gives \eqref{eq:kappa}.  Substituting
\(\widetilde\rho=Z\rho\) into
\eqref{eq:unnormalized-adjoint} gives
\eqref{eq:normalized-adjoint}.  The last statement follows from the
definition of \(Z\).
\end{proof}

This construction solves the linear equation before normalization.  The
coefficient \(\kappa\) is therefore determined by the solution and does not
enter the linear terminal problem.

\begin{lemma}[Weak adjoint construction]
\label{lem:weak-adjoint}
Assume
\begin{equation}\label{eq:weak-drift-class}
 V\in L^2_\tau H^1_y(B_R),
 \qquad a,b\in L^1_\tau,
 \qquad \diver c=-3a.
\end{equation}
Then \eqref{eq:unnormalized-adjoint} has a nonnegative weak solution with
\[
 \widetilde\rho\in L^\infty_\tau L^2_y(B_R)
 \cap L^2_\tau H^1_{0,y}(B_R).
\]
It is obtained as the weak limit of solutions with smooth drifts
\(c_\varepsilon\).  On every compact subinterval where \(Z\ge z_0>0\), the
normalized densities satisfy the same bounds, and
\(\sqrt\rho\in L^2_\tau H^1_{\mathrm{loc},y}\) whenever their Fisher
information is uniformly bounded.
\end{lemma}

\begin{proof}
Reverse time in \eqref{eq:unnormalized-adjoint} and approximate \(V,a,b\) by
smooth functions preserving \(\diver c_\varepsilon=-3a_\varepsilon\).
Galerkin approximation gives smooth weak solutions.  Testing by
\(\widetilde\rho_\varepsilon\) makes the transport contribution
\(-\tfrac12\int(\diver c_\varepsilon)
|\widetilde\rho_\varepsilon|^2\), so Gronwall's inequality uses only
\(a_\varepsilon\in L^1\).  This yields uniform
\(L^\infty L^2\cap L^2H^1_0\) bounds.  Choose the approximating drifts to
converge strongly in \(L^2L^6+L^1L^\infty\).  The first summand follows from
\(V\in L^2H^1\), and the affine part belongs to the second summand on
\(B_R\).  The bounds on \(\widetilde\rho_\varepsilon\) then allow
\(c_\varepsilon\widetilde\rho_\varepsilon\) to pass to the limit against
every smooth test function.  Weak compactness gives a solution of
\eqref{eq:unnormalized-adjoint}; testing the negative part proves
nonnegativity.  Division by \(Z\ge z_0\) preserves the estimates.  Finally,
\[
 \int\frac{|\nabla\rho|^2}{\rho+\varepsilon}\dd y
 =4\int\left|\nabla\sqrt{\rho+\varepsilon}\right|^2\dd y,
\]
and weak lower semicontinuity gives the last assertion as
\(\varepsilon\downarrow0\).
\end{proof}

\subsection{Transport identity and the Gaussian equality model}

For smooth positive \(\rho\), write
\begin{equation}\label{eq:f-def}
 \rho=(4\pi\nu\theta)^{-3/2}e^{-f},
 \qquad \theta>0.
\end{equation}
Equation \eqref{eq:normalized-adjoint} gives
\begin{equation}\label{eq:f-equation}
 \partial_\tau f+\nu\Delta f-\nu|\nabla f|^2+c\cdot\nabla f
 =\kappa-\frac{3\theta'}{2\theta}+\diver c.
\end{equation}

\begin{lemma}[Localized adjoint transport identity]
\label{lem:adjoint-transport}
For every smooth scalar \(F\),
\begin{align}
 \frac{\dd}{\dd\tau}\int_{B_R}\chi^2F\rho\dd y
 ={}&\int_{B_R}\chi^2
 (\partial_\tau F-\nu\Delta F+c\cdot\nabla F)\rho\dd y
 \notag\\
 &+\int_{B_R}
 \left[F(-\nu\Delta\chi^2+c\cdot\nabla\chi^2)
       -2\nu\nabla\chi^2\cdot\nabla F\right]\rho\dd y
 \notag\\
 &-\kappa\int_{B_R}\chi^2F\rho\dd y.
 \label{eq:adjoint-transport}
\end{align}
\end{lemma}

\begin{proof}
Differentiate the integral, substitute
\eqref{eq:normalized-adjoint}, and integrate the Laplacian and divergence
terms by parts.  All spatial boundary terms vanish because
\(\chi\Subset B_R\).
\end{proof}

The three lines of \eqref{eq:adjoint-transport} separate interior evolution,
cutoff transport, and normalization.  Taking
\begin{equation}\label{eq:adjoint-W-density}
 W_f:=\nu\theta|\nabla f|^2+f-3,
 \qquad
 \mathcal W_{\mathrm{ad},\chi}(\tau)
 :=\int_{B_R}\chi^2W_f\rho\dd y,
\end{equation}
gives the exact localized derivative by setting \(F=W_f\) in
\eqref{eq:adjoint-transport} and using \eqref{eq:f-equation}.  In
particular, derivatives of \(\chi\) and the scalar \(\kappa\) remain visible
in the last two lines.  The velocity drift contributes through
\(\nabla c\), and its compression contributes through \(\diver c=-3a\).
No sign is assigned to those local terms without an additional estimate.

The equality model is most transparent without localization.

\begin{proposition}[Adjoint entropy production with a variable heat scale]
\label{prop:adjoint-entropy-production}
Let \(\rho\) be a smooth positive probability density on \(\R^3\) satisfying
\[
 -\partial_\tau\rho-\nu\Delta\rho-\diver(c\rho)=0
\]
with sufficient decay.  With \(f\) as in \eqref{eq:f-def}, for every
positive absolutely continuous \(\theta\) one has
\begin{align}
 \frac{\dd}{\dd\tau}
 \int_{\R^3}(\nu\theta|\nabla f|^2+f-3)\rho\dd y
 ={}&2\nu^2\theta\int_{\R^3}
 \left|\nabla^2f-\frac{1}{2\nu\theta}I\right|^2\rho\dd y
 \notag\\
 &+\int_{\R^3}(\diver c)\rho\dd y
 \notag\\
 &+2\nu\theta\int_{\R^3}
 \left(\nabla f\cdot\nabla\diver c
 -(\nabla f)^{T}(\nabla c)\nabla f\right)\rho\dd y
 \notag\\
 &+(\theta'+1)
 \left(\nu\int_{\R^3}|\nabla f|^2\rho\dd y
       -\frac{3}{2\theta}\right).
 \label{eq:adjoint-entropy-production}
\end{align}
For \(c=0\) and \(\theta'=-1\), the derivative is nonnegative and vanishes precisely when
\(\nabla^2f=(2\nu\theta)^{-1}I\) on the support of \(\rho\).
\end{proposition}

\begin{proof}
Let \(H(\rho)=-\int\rho\log\rho\) and
\(I(\rho)=\int|\nabla\log\rho|^2\rho\).  Direct differentiation with the
backward adjoint equation and integration by parts give
\[
 \frac{\dd H}{\dd\tau}
 =-\nu I+\int(\diver c)\rho
\]
and
\[
 \frac{\dd I}{\dd\tau}
 =2\nu\int|\nabla^2\log\rho|^2\rho
 -2\int\left(
 \nabla\log\rho\cdot\nabla\diver c
 +(\nabla\log\rho)^{T}(\nabla c)\nabla\log\rho
 \right)\rho.
\]
Differentiate
\[
 \nu\theta I+H-\frac32\log(4\pi\nu\theta)-3
\]
in \(\tau\), use \(\nabla\log\rho=-\nabla f\), and first collect the
terms corresponding to \(\theta'=-1\).  Completing their square gives the
first three lines of \eqref{eq:adjoint-entropy-production}.  The difference
between general \(\theta'\) and \(-1\) is exactly its last line.
\end{proof}

For the repaired Navier--Stokes drift \(c=V-ay-b\), one has
\(\nabla\diver c=0\) and
\(\nabla c=\nabla V-aI\).  Hence
the choice \eqref{eq:covariant-heat-scale} cancels every scalar scale term in
\eqref{eq:adjoint-entropy-production} and gives
\begin{align}
 \frac{\dd}{\dd\tau}\mathcal W_{\mathrm{ad}}
 ={}&2\nu^2\theta\int_{\R^3}
 \left|\nabla^2f-\frac{1}{2\nu\theta}I\right|^2\rho\dd y
 \notag\\
 &-2\nu\theta\int_{\R^3}
 (\nabla f)^{T}
 (\nabla V)_{\mathrm{sym}}
 \nabla f\,\rho\dd y.
 \label{eq:adjoint-entropy-repaired-drift}
\end{align}
Here \(\mathcal W_{\mathrm{ad}}\) denotes the whole-space entropy in
\cref{prop:adjoint-entropy-production}, and
\((\nabla V)_{\mathrm{sym}}
:=\tfrac12(\nabla V+(\nabla V)^{T})\).  Formula
\eqref{eq:adjoint-entropy-repaired-drift} isolates the positive Gaussian
square from velocity strain.  The translation coefficient \(b\) drops out
because it has no spatial derivative, while \(a\) drops out because the heat
scale is transported by \eqref{eq:covariant-heat-scale}.  The strain term has
no sign by itself, but it completes a square with viscous production below.

In the unmodulated case \(a=b=0\), for which \(\theta'=-1\), the Euclidean
heat kernel is
\begin{equation}\label{eq:heat-kernel}
 \rho_*(y,\tau)=(4\pi\nu\theta)^{-3/2}
 \exp\left(-\frac{|y|^2}{4\nu\theta}\right),
 \qquad f_*=\frac{|y|^2}{4\nu\theta},
\end{equation}
and satisfies
\[
 \nabla^2f_*=\frac{1}{2\nu\theta}I,
 \qquad
 \int_{\R^3}(\nu\theta|\nabla f_*|^2+f_*-3)\rho_*\dd y=0.
\]
Hence both the entropy and its production vanish.  This is the Euclidean
normalization of Perelman's \(\mathcal W\)-entropy \cite{Perelman2002}.

For weak densities, the Fisher term is defined by
\begin{equation}\label{eq:weak-fisher}
 \int\chi^2\frac{|\nabla\rho|^2}{\rho}\dd y
 :=4\int\chi^2|\nabla\sqrt\rho|^2\dd y,
\end{equation}
with value \(+\infty\) when \(\sqrt\rho\notin H^1_{\mathrm{loc}}\).
The entropy uses \(0\log0=0\).  At weak regularity the Fisher information has
its usual lower-semicontinuous extension.  The Shannon term in
\(\mathcal W_{\mathrm{ad},\chi}\) has the opposite semicontinuity orientation,
so its passage to a limit requires an explicit strong-convergence and
uniform-integrability argument, as recorded in
\cref{app:normalized-adjoint-solutions}; no blanket
lower-semicontinuity claim is made for the full entropy.  No pointwise equation for
\(f=-\log\rho+\text{constant}\) is asserted on the zero set of \(\rho\).
Smooth positive adjoint approximations justify the localized transport
formulae needed below.

\section{Stokes-metric projection onto localized kinetic variations}
\label{sec:alignment}

\subsection{Scale-critical kinetic probes}

On a compact represented-time interval \(I_\tau\), an admissible probe
consists of a repaired scale--translation chart \((X,\Lambda)\) and a
nonnegative function
\begin{equation}\label{eq:probe-m}
 \begin{gathered}
 m\in W^{1,1}(I_\tau;W^{1,\infty}(B_R))
     \cap L^\infty(I_\tau;W^{2,\infty}(B_R)),\\
 \supp m(\cdot,\tau)\subset B_{R_0}\quad(0<R_0<R),
 \qquad \int_{B_R}m(y,\tau)\dd y=1.
 \end{gathered}
\end{equation}
The smooth adjoint weights \(m=\chi^2\rho\) from \cref{sec:adjoint}, and
regularized weak adjoint weights with the same normalization, are the
principal examples.  The physical multiplier is
\begin{equation}\label{eq:physical-weight}
 w_\Gamma(x,t(\tau))
 :=\Lambda(\tau)^{-1}
 m\left(\frac{x-X(\tau)}{\Lambda(\tau)},\tau\right).
\end{equation}
Its physical mass and weighted velocity mean are
\begin{equation}\label{eq:physical-weighted-mean}
 M_\Gamma(t):=\int_{\R^3}w_\Gamma\dd x=\Lambda(\tau)^2,
 \qquad
 \bar u_\Gamma(t):=M_\Gamma(t)^{-1}
       \int_{\R^3}w_\Gamma u\dd x.
\end{equation}
The Galilean-reduced localized kinetic probe is
\begin{equation}\label{eq:kinetic-probe}
 F_\Gamma(u(t)):=\frac12\int_{\R^3}w_\Gamma(x,t)
       |u(x,t)-\bar u_\Gamma(t)|^2\dd x.
\end{equation}
On the represented cylinder,
\begin{equation}\label{eq:probe-chart-form}
 F_\Gamma(u(t(\tau)))
 =\frac12\int_{B_R}m(y,\tau)|V(y,\tau)-\bar V_m(\tau)|^2\dd y,
 \qquad
 \bar V_m:=\int_{B_R}mV\dd y.
\end{equation}
Thus \(F_\Gamma\) is invariant under the Navier--Stokes scaling and is a
Galilean-invariant version of the local quantity \(A\).  Denote its chart
expression by
\begin{equation}\label{eq:oscillatory-probe}
 \OscEnergy_\Gamma(V(\tau))
 :=\frac12\int_{B_R}m|V-\bar V_m|^2\dd y.
\end{equation}
For variations of \(V\) with \(m\) fixed,
\begin{equation}\label{eq:oscillatory-gradient}
 D\OscEnergy_\Gamma(V)[\varphi]
 =\int_{B_R}m(V-\bar V_m)\cdot\varphi\dd y.
\end{equation}
Indeed, the variation of \(\bar V_m\) pairs with
\(\int m(V-\bar V_m)=0\).  Thus the target gradient below is the Stokes
gradient of the Galilean-reduced variance in every representative.  No
time-dependent velocity-centering condition is imposed on the chart.

Let
\begin{equation}\label{eq:q-def}
 q_\Gamma:=\Pdiv\bigl(w_\Gamma(u-\bar u_\Gamma)\bigr).
\end{equation}
For variations \(\varphi\in\mathcal H_\sigma\),
\[
 D F_\Gamma(u)[\varphi]=(q_\Gamma,\varphi)_{L^2}.
\]
Since \(q_\Gamma\in\mathcal V_\sigma\) for almost every time, this derivative
extends to currents \(J\in\mathcal V_\sigma^*\) by
\(D F_\Gamma(u)[J]:=\langle J,q_\Gamma\rangle\).
The Stokes gradient current of the probe is
\begin{equation}\label{eq:target-gradient-current}
 G_\Gamma:=Kq_\Gamma.
\end{equation}
Define the viscous action, the Stokes norm of the localized variation, and
the signed localized power by
\begin{align}
 \Action_\nu(u)
 &:=\|J_\nu\|_{K^{-1}}^2
 =\nu\|\nabla u\|_2^2,\label{eq:viscous-action}\\
 \ProbeDirichlet_\Gamma(u)
 &:=\|G_\Gamma\|_{K^{-1}}^2
 =\nu\|\nabla q_\Gamma\|_2^2,\label{eq:localized-variation-energy}\\
 \Power_\Gamma(u)
 &:=(J_\nu,G_\Gamma)_{K^{-1}}
 =D F_\Gamma(u)[J_\nu].\label{eq:target-power}
\end{align}

\begin{lemma}[Explicit viscous power]
\label{lem:explicit-power}
For smooth \(u\) and \(w_\Gamma\),
\begin{equation}\label{eq:explicit-power}
 \Power_\Gamma
 =-\nu\int_{\R^3}w_\Gamma|\nabla u|^2\dd x
 +\frac\nu2\int_{\R^3}|u-\bar u_\Gamma|^2
       \Delta w_\Gamma\dd x.
\end{equation}
The first term is local viscous dissipation.  The second is the diffusive
flux across the spatial variation of the probe.
\end{lemma}

\begin{proof}
Since \(A_{\mathrm S}u\) is solenoidal and \(\Pdiv\) is self-adjoint,
\[
 \Power_\Gamma
 =-\nu\int\nabla u:\nabla q_\Gamma
 =-\nu\int\nabla u:\nabla\bigl(w_\Gamma(u-\bar u_\Gamma)\bigr).
\]
Expand the last gradient and integrate
\(-\nu\int \partial_j u_i(u_i-\bar u_{\Gamma,i})
\partial_jw_\Gamma\) by parts.  This gives
\eqref{eq:explicit-power}.
\end{proof}

In the represented variables, the power measure has the scale-free form
\begin{equation}\label{eq:chart-power}
 \Power_\Gamma(u(t(\tau)))\,\dd t
 =\widehat{\Power}_\Gamma(V(\tau))\,\dd\tau,
 \qquad
 \widehat{\Power}_\Gamma(V)
 :=-\nu\int_{B_R}m|\nabla V|^2\dd y
 +\frac{\nu}{2}\int_{B_R}|V-\bar V_m|^2\Delta m\dd y.
\end{equation}

For smooth solutions the total derivative of the probe splits as
\begin{equation}\label{eq:probe-balance}
 \frac{\dd}{\dd t}F_\Gamma(u(t))
 =\Tgeom_\Gamma(u,p)+\Power_\Gamma(u),
\end{equation}
where
\begin{equation}\label{eq:geometric-transport}
 \Tgeom_\Gamma(u,p)
 :=\frac12\int|u-\bar u_\Gamma|^2
       (\partial_t w_\Gamma+u\cdot\nabla w_\Gamma)\dd x
 +\int p\,(u-\bar u_\Gamma)\cdot\nabla w_\Gamma\dd x.
\end{equation}
The term \(\Tgeom_\Gamma\) contains conservative transport, pressure flux,
and the motion of the scale--center coordinates.  The term
\(\Power_\Gamma\) contains viscosity alone.  For suitable weak solutions,
define
\begin{equation}\label{eq:local-defect-measure}
 \mu_{\mathrm{loc}}
 :=-\left[
 \partial_t\frac{|u|^2}{2}
 +\diver\left(\left(\frac{|u|^2}{2}+p\right)u\right)
 -\nu\Delta\frac{|u|^2}{2}
 +\nu|\nabla u|^2
 \right].
\end{equation}
The local energy inequality says precisely that this distribution is a
nonnegative Radon measure.  Combining it with the weak momentum identity for
the weighted mean gives
\begin{equation}\label{eq:weak-probe-balance}
 F_\Gamma(t_2)-F_\Gamma(t_1)
 =\int_{t_1}^{t_2}(\Tgeom_\Gamma+\Power_\Gamma)\dd t
 -\mu_{\mathrm{loc}}(w_\Gamma;[t_1,t_2]),
\end{equation}
where \(\mu_{\mathrm{loc}}(w_\Gamma;[t_1,t_2])\) denotes the measure in
\eqref{eq:local-defect-measure} paired with the nonnegative test weight on the
indicated time interval.  The identity holds at the corresponding local
energy times; arbitrary compact subintervals follow by one-sided
representatives.

\subsection{The Stokes-metric projection inequality}

\begin{definition}[Viscous alignment]
\label{def:alignment}
Set
\begin{equation}\label{eq:alignment}
 \Align_\Gamma(u)
 :=
 \begin{cases}
 \displaystyle
 \frac{(\Power_\Gamma(u)_+)^2}
 {\Action_\nu(u)\ProbeDirichlet_\Gamma(u)},
 &\Action_\nu(u)\ProbeDirichlet_\Gamma(u)>0,\\[3mm]
 0,&\text{otherwise}.
 \end{cases}
\end{equation}
Positive power means that the viscous current points toward increasing the
localized kinetic probe.  Negative power means that viscosity lowers that
probe.
\end{definition}

\begin{theorem}[Stokes-metric projection inequality]
\label{thm:stokes-projection}
For almost every time,
\begin{equation}\label{eq:stokes-projection-pointwise}
 (\Power_\Gamma(u)_+)^2
 \le\Action_\nu(u)\ProbeDirichlet_\Gamma(u),
 \qquad 0\le\Align_\Gamma(u)\le1.
\end{equation}
For every interval \(I\Subset(0,T)\),
\begin{equation}\label{eq:stokes-projection-integrated}
 \int_I\Power_\Gamma(u(t))_+\dd t
 \le
 \left(\int_I\Action_\nu(u(t))\dd t\right)^{1/2}
 \left(\int_I\ProbeDirichlet_\Gamma(u(t))\dd t\right)^{1/2}.
\end{equation}
The first factor is at most \(\Energy(u(0))^{1/2}\).
\end{theorem}

\begin{proof}
The definitions give
\[
 \Power_\Gamma=(J_\nu,G_\Gamma)_{K^{-1}}.
\]
Cauchy--Schwarz in the \(K^{-1}\)-metric proves
\eqref{eq:stokes-projection-pointwise}.  A second application of
Cauchy--Schwarz in time proves \eqref{eq:stokes-projection-integrated}.  The
energy bound is \eqref{eq:finite-action}.
\end{proof}

The inequality separates the amount of irreversible motion from its
orientation.  Equal action can produce zero or maximal positive power,
depending on the angle with the target gradient.

\begin{definition}[Covariant ray-alignment defect]
\label{def:covariant-W}
For a sufficiently regular divergence-free comparison curve \(v\), set
\begin{equation}\label{eq:equation-defect}
 R(v):=\Dcov_t^{\mathrm E}v+Kv.
\end{equation}
For an admissible probe \(\Gamma\) and interval \(I\), define
\begin{equation}\label{eq:covariant-W}
 \begin{aligned}
 \RayDefect_{\Gamma,I}(v)
 :={}&\frac12\int_I\|R(v)\|_{K^{-1}}^2\dd t\\
 &+\frac12\int_I
 \left(
 \Action_\nu(v)
 -\frac{(\Power_\Gamma(v)_+)^2}{\ProbeDirichlet_\Gamma(v)}
 \right)\dd t,
 \end{aligned}
\end{equation}
where the quotient is defined as zero when
\(\ProbeDirichlet_\Gamma=0\).
\end{definition}

\begin{proposition}[Metric meaning of the ray defect]
\label{prop:W-distance}
The functional \(\RayDefect_{\Gamma,I}\) is nonnegative.  Pointwise in time,
\begin{equation}\label{eq:ray-distance}
 \Action_\nu
 -\frac{(\Power_\Gamma{}_+)^2}{\ProbeDirichlet_\Gamma}
 =\operatorname{dist}_{K^{-1}}^2
 \bigl(J_\nu,\{\alpha G_\Gamma:\alpha\ge0\}\bigr).
\end{equation}
For an exact Navier--Stokes solution, the equation-defect term in
\eqref{eq:covariant-W} vanishes.  Equality
\(\RayDefect_{\Gamma,I}(u)=0\) then means that the actual viscous current lies on
the nonnegative target-gradient ray almost everywhere on \(I\).  Moreover,
\begin{equation}\label{eq:W-alignment-identity}
 2\RayDefect_{\Gamma,I}(u)
 =\int_I\Action_\nu(u(t))
   \bigl(1-\Align_\Gamma(u(t))\bigr)\dd t.
\end{equation}
\end{proposition}

\begin{proof}
Orthogonal projection of a vector \(J\) onto the closed ray generated by a
nonzero vector \(G\) has coefficient
\((J,G)_+ /\|G\|^2\).  Its squared residual norm is
\(\|J\|^2-(J,G)_+^2/\|G\|^2\).  Apply this identity in the
\(K^{-1}\)-metric.  When \(G=0\), the distance is \(\|J\|\), in agreement
with the convention in \eqref{eq:covariant-W}.  The equation-defect statement
is \eqref{eq:covariant-gradient}.  Substitution of
\eqref{eq:alignment} gives \eqref{eq:W-alignment-identity}.
\end{proof}

Let \(\mathfrak G_{\mathrm{adm}}(I)\) denote repaired probes with
scale-relative cutoff bounds and normalized adjoint weights as in
\cref{sec:adjoint}.  Let \(\mathfrak G_{\mathrm{adm}}(t)\) be their time-\(t\)
slices.  Define
\begin{equation}\label{eq:optimized-alignment}
 \Align_*(t):=\sup_{\Gamma\in\mathfrak G_{\mathrm{adm}}(t)}
 \Align_\Gamma(t),
 \qquad
 \RayDefect_*(I):=\inf_{\Gamma\in\mathfrak G_{\mathrm{adm}}(I)}
 \RayDefect_{\Gamma,I}(u).
\end{equation}
Suprema and infima are understood through approximating sequences; no
attainment is required.  The Stokes mobility \(K\) and current \(J_\nu\) are
fixed throughout these optimizations.

\begin{theorem}[Local energy identities and projection bound]
\label{thm:local-monotonicity}
For every admissible probe and interval \([t_1,t_2]\Subset(0,T)\), the
following three statements hold simultaneously:
\begin{enumerate}[label=\textup{(\roman*)},leftmargin=2.2em]
\item global energy decreases by viscous action and the nonnegative energy
defect according to \eqref{eq:weak-energy-balance};
\item the localized probe satisfies \eqref{eq:weak-probe-balance}, separating
coordinate/transport flux from viscous target power;
\item positive viscous target power satisfies
\eqref{eq:stokes-projection-integrated} and is therefore controlled by the finite
global action.
\end{enumerate}
These conclusions remain valid along an optimizing sequence for
\eqref{eq:optimized-alignment}.
\end{theorem}

\begin{proof}
Items \textup{(i)--(iii)} are
\cref{thm:weak-energy-balance}, the suitable local energy inequality tested by
\(w_\Gamma\), and \cref{thm:stokes-projection}, respectively.  Each inequality
holds for every admissible probe, so it holds termwise along any approximating
sequence in \eqref{eq:optimized-alignment}.
\end{proof}

\subsection{Alignment, coercivity, and coherent probe selection}

The alignment is an angle in the Stokes metric.  It contains no information
about the amount of action unless it is paired with \(\Action_\nu\), and it
contains no compactness information unless the probe belongs to a normalized
family.  The following elementary observation fixes the equality orientation.

\begin{lemma}[Positive-ray equality for a kinetic probe]
\label{lem:positive-ray-equality}
Let \(w\ge0\) have positive finite mass, put
\(\bar u_w=(\int w)^{-1}\int wu\), let
\(q=\Pdiv(w(u-\bar u_w))\), and suppose
\(u\in L^2_\sigma\cap\dot H^1_\sigma\).  If
\begin{equation}\label{eq:positive-ray-equality}
 -Ku=\alpha Kq
 \qquad\text{for some }\alpha\ge0,
\end{equation}
then \(u=0\).
\end{lemma}

\begin{proof}
Apply \(K^{-1}\) in the solenoidal Gelfand triple and pair the resulting
identity \(-u=\alpha q\) with \(u\).  Since \(u\) is solenoidal,
\[
 -\|u\|_2^2
 =\alpha(q,u)_{L^2}
 =\alpha\int w|u-\bar u_w|^2\dd x\ge0.
\]
Both sides vanish, so \(u=0\).
\end{proof}

For a finite-energy whole-space state, vanishing ray defect is therefore a
rigid equality.  Along a singular rescaling, however, the physical ray defect
is weighted by \(\Lambda\) exactly as in
\eqref{eq:physical-versus-critical-action}.  It may have finite physical
integral even when its scale-critical chart density remains positive on every
unit window.  This explains why the ray defect is only a covariant diagnostic;
the exclusion proof instead controls the local Dirichlet chart density through
the moving-cutoff energy.

Probe optimization also has a precise order.  The physical solution first
determines \(J_\nu\) and the canonical chart.  A fixed terminal-data rule then
selects the finite-horizon adjoint problem on the padded window
\eqref{eq:adjoint-buffer}.  Only after these choices are fixed may one take a
supremum in \(\Align\) or an infimum in \(\RayDefect\).  Replacing the current
by the gradient of an optimizing probe would change the equation and is not
an admissible variation.

\begin{definition}[Compatible family of localized probes]
\label{def:compatible-probe-family}
A compatible family of localized probes on a represented orbit is a family
\(\{\Gamma_j\}_{j\ge1}\) on padded retained windows such that:
\begin{enumerate}[label=\textup{(\roman*)},leftmargin=2.2em]
\item every \(\Gamma_j\) is constructed from the same physical solution by a
fixed terminal-data rule and has the buffer \eqref{eq:theta-buffer};
\item chart transitions satisfy the composition law in
\cref{thm:canonical-critical-orbit};
\item the current paired with every target is the pullback of
\(J_\nu=-\nu A_{\mathrm S}u\);
\item approximating probes converge in the compact topology required by the
local adjoint construction.
\end{enumerate}
\end{definition}

The terminal adjoint equation is a testing device.  Coherence here means that
all probes test the same physical current.  It does not identify distinct
finite-horizon adjoint densities with one global density, and no such
identification is used by the corrected localized energy below.

\begin{proposition}[Scale-critical window form of the ray defect]
\label{prop:critical-ray-defect}
Let \(I_\tau\) be a represented interval and \(I_t\) its physical image.  For
an exact Navier--Stokes solution, the pullback of the physical ray defect has
the form
\begin{equation}\label{eq:critical-ray-defect}
 2\RayDefect_{\Gamma,I_t}(u)
 =\int_{I_\tau}\Lambda(\tau)
 \widehat{\mathfrak D}_{\mathrm{ray},\Gamma}(V(\tau))\dd\tau,
\end{equation}
where \(\widehat{\mathfrak D}_{\mathrm{ray},\Gamma}\) is invariant under a
constant parabolic rescaling of the represented window.  Equivalently, it is
defined by the Radon--Nikodym relation
\begin{equation}\label{eq:critical-ray-radon-nikodym}
 \widehat{\mathfrak D}_{\mathrm{ray},\Gamma}(\tau)\dd\tau
 :=\Lambda(\tau)^{-1}
 \operatorname{dist}_{K^{-1}}^2
 \bigl(J_\nu,\R_+G_\Gamma\bigr)\dd t.
\end{equation}
\end{proposition}

\begin{proof}
Both vectors in the Stokes pairing transform by the pullback of the same
chart.  Their squared inverse-metric distance has the scaling of viscous
action.  Formula \eqref{eq:chart-action} shows that its physical measure is
\(\Lambda\) times the dimensionless represented measure, which is exactly
\eqref{eq:critical-ray-radon-nikodym}.
\end{proof}

The hatted density exhibits the same physical-versus-critical distinction as
\eqref{eq:physical-versus-critical-action} and
\eqref{eq:renormalized-critical-action}.  It is retained below as a
covariant diagnostic, but the local coercive floor does not use it.

\section{Canonical critical orbits}
\label{sec:critical-orbit}

The covariant argument is carried by one normalized orbit rather than a list
of asymptotic regimes.  This section specifies the state space and the exact
compactness conclusion required by the entropy proof.

\subsection{Symmetry reduction and target-retaining windows}

Fix a normalized probe radius \(R_*>1\) and a nonnegative inner cutoff
\(\zeta\in C_c^\infty(B_{R_*})\),
\(0\le\zeta\le1\), \(\int\zeta^3>0\), for which
\(\{\zeta>0\}\) is connected and whose support
lies where the adjoint
cutoff equals one.  Put
\begin{equation}\label{eq:fixed-inner-mean}
 \bar V_\zeta(\tau)
 :=\frac{\int\zeta^3V(y,\tau)\dd y}{\int\zeta^3\dd y}.
\end{equation}
For a unit represented interval \(I=[s,s+1]\), define
\begin{equation}\label{eq:oscillatory-mass}
 \OscMass_\Gamma(V;I)
 :=\int_I\int_{B_{R_*}}
       \zeta^3|V-\bar V_\zeta|^3\dd y\dd\tau.
\end{equation}
This observable is independent of the adjoint terminal datum, continuous
under strong local \(L^3\) convergence, and invariant under constant velocity
shifts.

The relevant symmetries are spatial translation, positive parabolic dilation,
constant-velocity Galilean transformation, and addition of a function of time
to the pressure.  The repaired variables \((X,\Lambda,a,b)\) encode scale and
center.  We fix those representatives by
\begin{equation}\label{eq:canonical-moments}
 \mathfrak F(1,0;V(\tau))=0,
 \qquad
 \det D_{(\lambda,q)}\mathfrak F(1,0;V(\tau))\ne0
\end{equation}
at almost every active time, with absolutely continuous continuation supplied
by the canonical-observer construction.  Galilean symmetry is not fixed by
\eqref{eq:canonical-moments}; the gauge itself is invariant under it, and the
scalar observables remove it through \eqref{eq:fixed-inner-mean} and
\eqref{eq:physical-weighted-mean}.
The pressure gauge is
\begin{equation}\label{eq:canonical-pressure-gauge}
 \int_{B_1}P(y,\tau)\dd y=0.
\end{equation}
The nondegeneracy and absolute continuity of this gauge are conclusions of the
canonical-observer theorem, not extra conditions imposed on an arbitrary weak
solution.

\begin{definition}[Normalized target-retaining window]
\label{def:normalized-target-window}
Fix \(\eta_0>0\) and a collection of compact-cylinder bounds
\(\mathbf M=(M_R)_{R\ge1}\).  A represented unit window
\(\mathfrak w=(V,P,a,b;I)\) is an
\((\eta_0,\mathbf M)\)-normalized target-retaining window when:
\begin{enumerate}[label=\textup{(\roman*)},leftmargin=2.2em]
\item \((V,P)\) is a suitable solution of \eqref{eq:renormalized-NS} on a
neighborhood of \(B_R\times I\) for every fixed \(R\);
\item the gauges \eqref{eq:canonical-moments} and
\eqref{eq:canonical-pressure-gauge} hold, and the inverse of the gauge
Jacobian in \eqref{eq:gauge-jacobian} is uniformly bounded on the window;
\item the local suitable norms satisfy, on each fixed \(B_R\times I\),
\begin{equation}\label{eq:target-window-bounds}
 \|V\|_{L^\infty_I L^2(B_R)}
 +\|V\|_{L^2_I H^1(B_R)}
 +\|V\|_{L^3(B_R\times I)}
 +\|P\|_{L^{3/2}(B_R\times I)}
 +\|a\|_{L^1(I)}+\|b\|_{L^1(I)}
 \le M_R;
\end{equation}
\item the oscillatory floor
\begin{equation}\label{eq:target-window-floor}
 \OscMass_\Gamma(V;I)\ge\eta_0
\end{equation}
holds for the selected inner probe.
\end{enumerate}
\end{definition}

The subtraction in \eqref{eq:oscillatory-mass} is essential.  A spatially
constant represented velocity can have positive uncentered \(L^3\) mass, but
it has no spatial oscillation and cannot by itself support a singular point;
constant Galilean shifts are a special case.  The canonical-observer theorem
therefore retains concentration only after passing to this quotient.

\subsection{Canonical observation}

\begin{theorem}[Canonical observer at terminal loss]
\label{thm:canonical-critical-orbit}
Let \(u\) be a classical finite-energy solution of \eqref{eq:NS} on
\([0,T)\), where \(T<\infty\) is its maximal classical existence time, and
let \((u,p)\) denote any suitable continuation through \(T\).  Then there
exist \(\eta_0=\eta_0(\nu)>0\), compact bounds \(\mathbf M\), an absolutely
continuous repaired chart \((X,\Lambda)\), and a represented suitable orbit
\[
 (V,P,a,b)(\tau),\qquad \tau\in[0,\infty),
\]
with unit intervals \(I_j=[\tau_j,\tau_j+1]\) having pairwise disjoint
interiors and \(\tau_j\to\infty\), such that:
\begin{enumerate}[label=\textup{(\roman*)},leftmargin=2.2em]
\item after passing once to a subsequence and relabeling, with
\(x_j:=X(\tau_j)\) and \(r_j:=\Lambda(\tau_j)\), one has
\(r_j\downarrow0\), and the physical time images of the windows converge to
\(T\); the centers may remain bounded or escape to spatial infinity;
\item every restricted state
\(\mathfrak w_j=(V,P,a,b;I_j)\) is an
\((\eta_0,\mathbf M)\)-normalized target-retaining window;
\item in the product topology specified in \cref{sec:quotient-geometry} and
made explicit here, the sequential closure of the selected windows is
sequentially compact.  More explicitly, after translating
each time interval to \([0,1]\), every sequence in that closure has a
subsequence for which
\begin{equation}\label{eq:critical-orbit-compactness}
 V_j\to V_\infty\ \text{ strongly in }L^3_{\mathrm{loc}},
 \qquad
 \nabla V_j\rightharpoonup\nabla V_\infty
 \ \text{ in }L^2_{\mathrm{loc}},
\end{equation}
 on every compact spatial cylinder in \(\R^3\times[0,1]\), and in particular
 strongly in \(L^3(\supp\zeta\times[0,1])\),
 with strong local \(L^{3/2}\) convergence of normalized pressures and
 \(L^1\) convergence of \((a,b)\);
\item chart changes obey the composition law and represent the same physical
solution.
\end{enumerate}
\end{theorem}

\begin{proof}
See Appendix~\ref{app:critical-orbit}, in particular
\cref{thm:record-point-ancient-extraction,prop:first-bad-scale-descendant,prop:coherent-observer-interpolation,prop:canonical-orbit-compactness-retention}.
\end{proof}

\begin{remark}[Why one orbit is required]
The corrected energy decrease is accumulated along one scalar function of represented time.
Independent rescaling on each target-retaining window would provide unrelated endpoint
values and would not support the telescoping argument.  Clause~\textup{(iv)}
is therefore part of the construction rather than a convention made after
extraction.
\end{remark}

\subsection{The two classical scale regimes}

Let \(\ell(t)=\sqrt{T-t}\) denote the parabolic terminal scale.  Along the
canonical orbit, the dimensionless ratio
\begin{equation}\label{eq:scale-ratio}
 \mathfrak r(\tau):=\frac{\Lambda(\tau)}{\ell(t(\tau))}
\end{equation}
records the classical distinction.  A Type~I window has
\(\mathfrak r\) bounded above and below after normalization; in the centered
choice \(t=-e^{-\tau}\), \(\Lambda=e^{-\tau/2}\), one has
\(a=-1/2\).  A Type~II sequence contains windows on which
\(\mathfrak r\to0\), and its coefficient \(a=\Lambda'/\Lambda\) must be
controlled through the repaired gauge \eqref{eq:moment-map}.

The covariant equation \eqref{eq:chart-gradient-flow} is identical in the two
regimes.  The difference is carried entirely by \((a,b)\) and by the
asymptotic path of the normalized orbit.  The moving-cutoff alternative is
consequently formulated for locally integrable \(a,b\), without fixing
\(a=-1/2\) or assuming convergence of \(a\).  A centered ancient orbit, a
rapid scale collapse, and a recurrent repaired orbit are either reduced to
the corresponding local estimates or tested by the same moving production measure.
The gap argument uses only the coherent orbit, compactness, the oscillatory
floor, and the summable moving-cutoff case.
Thus the argument uses precisely the compactness in
\cref{thm:canonical-critical-orbit} and introduces no additional rescaled
solution.

There remains one compatibility issue.  A window selected at time
\(\tau_j\) must extend to the next selected window without replacing the
underlying solution.  Let
\(g_{j+1\leftarrow j}\in\mathcal G_{\mathrm{ch}}\) be the accumulated canonical
transition from the terminal representative of \(I_j\) to the initial
representative of \(I_{j+1}\).  Coherence means
\begin{equation}\label{eq:chart-transition-composition}
 g_{k\leftarrow j}
 =g_{k\leftarrow k-1}\circ\cdots\circ g_{j+1\leftarrow j},
 \qquad j<k,
\end{equation}
and every composition pulls back the same physical pair \((u,p)\).  The
backward adjoint localization and its cutoff are transported on this one
orbit.  There is no cross-window reset in the exclusion argument.

\section{Exact critical financing}
\label{sec:dissipation-reduction}

The material identities identify the mechanism but do not by themselves
produce a scale-critical finite measure.  The required statement is a
dissipation-or-reduction alternative: signed transport between nested regions is internal, whereas
viscous dissipation and true exterior leakage are counted once.  Its precise
interface with the compact orbit is derived below from the localized energy
equation.

Let \(I=[s,t]\) be a represented interval and let
\(\psi\in C^\infty(I;C_c^\infty(\R^3))\), \(\psi\ge0\).  Set
\begin{equation}\label{eq:signed-local-energy}
 K_\psi(\tau):=\frac12\int_{\R^3}|V(y,\tau)|^2\psi(y,\tau)\dd y,
 \qquad
 D_I(\psi):=\nu\int_I\int_{\R^3}|\nabla V|^2\psi\dd y\dd\tau,
\end{equation}
and define the signed localization current
\begin{align}
 J_I(\psi):={}&\int_I\bigg[
 \frac12\int |V|^2(\partial_\tau\psi+\nu\Delta\psi)
 +\int\left(\frac12|V|^2+P\right)V\cdot\nabla\psi \notag\\
 &\hspace{2.8em}
 -\frac a2\int|V|^2(\psi+y\cdot\nabla\psi)
 -\frac12 b\cdot\int|V|^2\nabla\psi
 \bigg]\dd\tau .
 \label{eq:signed-current}
\end{align}
All spatial integrals in \eqref{eq:signed-current} are over \(\R^3\).  The
pressure contribution is invariant under \(P\mapsto P+c(\tau)\), because
\(\int V\cdot\nabla\psi=0\).

\begin{proposition}[Signed-current additivity]
\label{prop:signed-current-additivity}
For every represented suitable solution there is a nonnegative local
energy-defect functional \(M_I(\psi)\), linear on nonnegative tests, such that
\begin{equation}\label{eq:signed-local-balance}
 K_\psi(t)+D_I(\psi)+M_I(\psi)
 = K_\psi(s)+J_I(\psi).
\end{equation}
For a smooth solution equality holds and \(M_I(\psi)=0\).  If
\(\psi=\sum_{k=1}^N\psi_k\) with \(\psi_k\ge0\), then
\begin{align}
 K_\psi&=\sum_{k=1}^NK_{\psi_k},
 &D_I(\psi)&=\sum_{k=1}^ND_I(\psi_k),
 &J_I(\psi)&=\sum_{k=1}^NJ_I(\psi_k),
 &M_I(\psi)&=\sum_{k=1}^NM_I(\psi_k).
 \label{eq:current-additivity}
\end{align}
Consequently, for an exact parent--child partition
\(\psi_\alpha=\sum_{\beta\in\operatorname{ch}(\alpha)}\psi_\beta\),
the difference between the parent current and the sum of its child currents
is identically zero.  This is refinement invariance, not selective
cancellation: adding the parent balance to all child balances double-counts
every term, whereas subtracting the parent balance cancels every term,
including \(D_I\) and \(M_I\).
\end{proposition}

\begin{proof}
The suitable local energy inequality for the repaired equation defines a
nonnegative defect distribution.  Integrating it between \(s\) and \(t\) and
placing that defect on the left gives
\eqref{eq:signed-local-balance}; the terms generated by the repaired drift are
the last two terms of \eqref{eq:signed-current}.  Every displayed term is linear in
\(\psi\) and its derivatives, which proves
\eqref{eq:current-additivity}.  For smooth solutions the local energy
identity holds with equality.
\end{proof}

\begin{theorem}[Exact critical financing identity]
\label{thm:exact-critical-financing}
Let \(C_0,\ldots,C_{N-1}\) be consecutive normalized spacetime cells along
the canonical observer.  The represented-time interval of \(C_j\) begins at
the incoming face of the \(j\)-th selected window and ends at the incoming
face of the next one; hence it contains both that active window and any
following inactive interval.  Choose compatible smooth localization tests
and oriented endpoint traces.  Write \(K_j^-\) and \(K_j^+\) for the
localized kinetic energy at the incoming and outgoing faces, \(D_j\) for all
viscous dissipation in \(C_j\), \(M_j\ge0\) for all suitable-energy defect in
\(C_j\), and \(H_j\) for the complete signed pressure, transport, dilation,
translation, and cutoff current.  Then
\begin{equation}\label{eq:exact-critical-financing-window}
 K_j^+ +D_j+M_j=K_j^-+H_j
 \qquad(0\le j<N),
\end{equation}
and
\begin{equation}\label{eq:exact-critical-financing-sum}
 \sum_{j=0}^{N-1}(D_j+M_j)
 =K_0^- -K_{N-1}^+ +\sum_{j=0}^{N-1}H_j.
\end{equation}
In particular, \(D_j\) dominates the dissipation on the active unit window
contained in \(C_j\).  The identity is invariant under addition of a function
of time to the pressure.
\end{theorem}

\begin{proof}
Apply the local energy equality with defect measure to each cell.  This is
\eqref{eq:signed-local-balance} with the nonnegative defect placed on the
left; all terms in which a derivative hits the localization or the moving
chart comprise \(H_j\), giving
\eqref{eq:exact-critical-financing-window}.  On a common face, the outgoing
kinetic trace of \(C_j\) is the incoming trace of \(C_{j+1}\).  The two
appear with opposite signs and cancel.  Because each cell already includes
the inactive interval preceding the next selected window, its dissipation and
defect remain in \(D_j+M_j\), rather than being reclassified as current.
Summation yields
\eqref{eq:exact-critical-financing-sum}.  Finally,
\(\int V\cdot\nabla\psi=0\) makes the pressure part invariant under
\(P\mapsto P+c(\tau)\).
\end{proof}

\subsection{Normalized shell decomposition}

Let
\begin{equation}\label{eq:NS-component-set}
 \mathfrak C_{\mathrm{NS}}
 :=\{\mathrm{mat},\mathrm{vis},\mathrm{int},\mathrm{gauge},\mathrm{ext}\}.
\end{equation}
The labels denote material deformation, viscous motion, signed interior
transport including pressure, chart and material-relabeling motion, and true
exterior flux.  This is a decomposition of the actual repaired
Navier--Stokes current, not a decomposition of its absolute value.

\begin{definition}[Normalized singular-target shells]
\label{def:normalized-target-shells}
Let \(\widehat{\mathcal X}_{\mathrm{sing}}\) be a Hausdorff state-or-germ
space containing the regular orbit and the nonremovable singular germs.  A
normalized shell presentation consists of nested neighborhoods
\(U_\varepsilon\) of the singular set and continuous functions
\(\mathscr S_\varepsilon:\widehat{\mathcal X}_{\mathrm{sing}}\to[0,1]\)
such that
\begin{equation}\label{eq:target-shell-values}
 \mathscr S_\varepsilon=0
 \quad\hbox{outside }U_{2\varepsilon},
 \qquad
 \mathscr S_\varepsilon=1
 \quad\hbox{on }U_\varepsilon.
\end{equation}
On every regular precontact segment on which it is used,
\(\mathscr S_\varepsilon(u(\cdot))\) is absolutely continuous and its chain
rule has the exact form
\begin{equation}\label{eq:NS-shell-chain-rule}
 \frac{\dd}{\dd t}\mathscr S_\varepsilon(u(t))
 =\sum_{c\in\mathfrak C_{\mathrm{NS}}}q^c_\varepsilon(t)
 \quad\hbox{for a.e. }t.
\end{equation}
A shell-crossing record is an interval
\(J_\varepsilon=[t^-_\varepsilon,t^+_\varepsilon]\) lying strictly before
contact such that the endpoint values in
\eqref{eq:target-shell-values} are zero and one, respectively.
\end{definition}

\begin{theorem}[Exact shell-crossing identity]
\label{thm:NS-shell-identity}
Every shell-crossing record for which the five functions in
\eqref{eq:NS-shell-chain-rule} are integrable satisfies
\begin{equation}\label{eq:NS-unit-shell-progress}
 1=\sum_{c\in\mathfrak C_{\mathrm{NS}}}Q^c_\varepsilon,
 \qquad
 Q^c_\varepsilon:=\int_{J_\varepsilon}q^c_\varepsilon(t)\dd t.
\end{equation}
Consequently
\begin{equation}\label{eq:NS-positive-shell-component}
 \max_{c\in\mathfrak C_{\mathrm{NS}}}(Q^c_\varepsilon)_+
 \ge\frac15.
\end{equation}
Neither conclusion uses an energy or entropy bound.
\end{theorem}

\begin{proof}
Integrate \eqref{eq:NS-shell-chain-rule} over the crossing interval and use
the endpoint values.  Since the sum of five real numbers is one, the sum of
their positive parts is at least one and one positive part is at least
\(1/5\).
\end{proof}

\begin{proposition}[Persistent contributing component]
\label{prop:persistent-NS-component}
Suppose contact supplies infinitely many shell-crossing records on pairwise
disjoint precontact events.  After passage to an infinite subsequence, one
component \(c_*\in\mathfrak C_{\mathrm{NS}}\) satisfies
\begin{equation}\label{eq:persistent-NS-component}
 Q^{c_*}_j\ge\frac15
 \qquad\hbox{for every index in the subsequence}.
\end{equation}
\end{proposition}

\begin{proof}
Choose on each record a component satisfying
\eqref{eq:NS-positive-shell-component}.  One of the five labels is selected
infinitely often.
\end{proof}

\subsection{Finite-depth concentration decomposition}
\label{sec:finite-depth-PDE-decomposition}

The shell identity becomes useful only after every loss of compactness or
current cancellation is retained as a new Navier--Stokes object.  We use the
following five PDE source families.  They describe the origin of a term and
do not introduce additional equations:
\begin{equation}\label{eq:five-PDE-source-families}
 \mathfrak F_{\rm PDE}
 :=\{\mathrm{vis},\mathrm{tr},\mathrm{quo},
       \mathrm{sc},\mathrm{bd}\}.
\end{equation}
Here \(\mathrm{vis}\) is the localized Stokes form and its Dirichlet
production; \(\mathrm{tr}\) is the signed convection--pressure--moving-frame
current together with kinetic storage; \(\mathrm{quo}\) consists of the
Leray, pressure-gauge, and Galilean quotients; \(\mathrm{sc}\) contains
parabolic rebasing, weak-product defects, and scale current; and
\(\mathrm{bd}\) contains exterior, annular, terminal, and trace currents.
The order in \eqref{eq:five-PDE-source-families} is fixed whenever a first
failure is selected.

\begin{definition}[Finite-depth localized event record]
\label{def:finite-depth-localized-event-record}
A finite-depth localized event record of length \(N\) consists of exact
parabolic cells \(C_0,\ldots,C_N\) cut from one physical suitable solution,
their normalized velocity and pressure representatives, compatible additive
pressure gauges, the scale--translation maps between consecutive cells, and
the local-energy coordinates
\eqref{eq:marked-local-energy-coordinates} on every common interface.  It is
\emph{target retaining} at level \(\eta>0\) when
\begin{equation}\label{eq:finite-record-target-floor}
 \mathcal O_{Q_1}(V_j)\ge\eta,
 \qquad 0\le j\le N.
\end{equation}
Every weak limit used in the record carries its equation residual, Reynolds
stress, normalized pressure, scale--translation parameters, suitable-energy
defect, and exterior current.  A limit is called ordinary only when all these
coordinates converge in the topologies in which their local weak identities
pass to the limit.
Records of increasing depth are \emph{compatible} when the restriction of a
depth-\(N+1\) realization to its first \(N\) cells is a realization of the
depth-\(N\) record with the same physical origin, pressure gauges, oriented
traces, and composed observer.  A terminal family is \emph{cofinal} when the
physical radius of its last cell tends to zero with the depth.
\end{definition}

\begin{theorem}[Finite-depth Navier--Stokes source alternative]
\label{thm:finite-depth-NS-source-alternative}
Let the record-point extraction in
\cref{thm:record-point-ancient-extraction}, together with its exact
descendants, generate a compatible cofinal family of target-retaining
localized event records of arbitrarily large finite depth.  After an ordered
subsequence selection, the first applicable one of the following outputs
occurs.
\begin{enumerate}[label=\textup{(F\arabic*)},leftmargin=2.8em]
\item The Galilean-reduced target is lost.  The selected velocity is spatially
homogeneous on the retained component.
\item One nonnegative target-aligned charge is bounded below by a fixed
positive number on a pairwise disjoint family of events and is dominated by
one finite physical Radon measure.
\item There is a first nonordinary coordinate.  In the order
\eqref{eq:five-PDE-source-families} it is a viscous-domain defect, a
convection or Reynolds defect, a non-affine pressure-quotient defect, a
parabolic scale or realization defect, or an exterior, annular, terminal, or
trace defect.  The coordinate and the shell covector detecting it are
retained at the same physical origin and scale.
\item Every finite coordinate is ordinary, all local identities pass to the
limit, and the inverse system of records gives an actual same-origin
scale-collapsing Navier--Stokes path satisfying
\eqref{eq:finite-record-target-floor} at every level.
\end{enumerate}
No pressure tail, diffuse velocity remainder, or escaping observer is
identified with an ordinary profile in this alternative.
\end{theorem}

\begin{proof}
Enumerate the countable compact-cylinder exhaustion, the tests in
\(\mathscr T\), and the coordinates in
\cref{def:finite-depth-localized-event-record}.  First test the quotient
target.  Its loss gives \textup{(F1)} by the connected-overlap argument in
\cref{prop:backward-concentration-remainder}.  On the retained target, test
the nonnegative charges admitted by the physical local-energy measure.  A
uniform positive value on disjoint events gives \textup{(F2)}.  A normalized
Dirichlet value carrying an additional scale factor is not admitted at this
step and remains in the scale-current test below.

For the remaining coordinates, select the first index in the fixed source
order and countable exhaustion at which ordinary convergence fails.  The
local pressure decomposition identifies failure of pressure compactness with
a non-affine harmonic tail; \cref{prop:harmonic-pressure-tail-accessibility}
retains its exterior ancestry.  Failure of velocity tightness is subjected to
the first-bad-scale selection of
\cref{prop:first-bad-scale-descendant}; hence it gives an actual rescaled
velocity profile, a homogeneous quotient, or a retained pressure tail.
Failure of an observer, defect, current, or realization coordinate is retained
in the corresponding coordinate of the finite-depth record.  These are
precisely the alternatives in \textup{(F3)}.

If no finite coordinate fails, diagonal compactness gives compatible limits
on every member of the exhaustion.  The inverse-limit support coordinate
supplies simultaneous realization by the original terminal sequence, and
compatibility of the observer maps gives one same-origin path.  Strong local
\(L^3\) convergence preserves \eqref{eq:finite-record-target-floor}, and
cofinality gives scale collapse.  This is \textup{(F4)}.  The first-applicable
convention makes the four outputs disjoint, and the countable coordinate
exhaustion makes them exhaustive.
\end{proof}

\begin{proposition}[Closure of the compact nonexpanding output]
\label{prop:compact-nonexpanding-output-closure}
Suppose the ordinary output \textup{(F4)} has a compact covariant closure on
which the normalized scale coefficient is nonpositive and every minimal
covariant component retains the target floor
\eqref{eq:finite-record-target-floor}.  Then \textup{(F4)} is impossible.
The same conclusion holds when the closure carries a stationary monotone
covariant ensemble in the sense of
\cref{def:stationary-monotone-ensemble} whose law retains the target almost
surely.
\end{proposition}

\begin{proof}
In the autonomous case apply
\cref{cor:no-target-recurrent-nonexpanding-limit}.  In the variable-coefficient
case \cref{thm:stationary-monotone-ensemble-rigidity} gives
\(\nabla V=0\) almost surely.  The quotient activity therefore vanishes,
contrary to \eqref{eq:finite-record-target-floor}.
\end{proof}

\begin{proposition}[First unbounded activity scale]
\label{prop:first-unbounded-activity-scale}
Let \(U\) be the bounded ancient profile in
\cref{thm:record-point-ancient-extraction}, and let \(z_j\) be points in the
diffuse remainder of
\cref{prop:backward-concentration-remainder}.  For a fixed
\(0<\varepsilon_0<\varepsilon_{\rm reg}\), let \(\rho_j\) be the first radius
at which
\(C_U^\circ(z_j,\rho_j)=\varepsilon_0\), after the homogeneous mode is
removed.  If \(\rho_j\to\infty\), the stopping-time construction gives one
of the following concrete outputs:
\begin{enumerate}[label=\textup{(T\arabic*)},leftmargin=2.8em]
\item an actual target-retaining parabolic descendant;
\item a spatially homogeneous velocity with affine pressure, which is
target-null;
\item a non-affine harmonic pressure tail with exterior ancestry.
\end{enumerate}
The last output remains a typed member of \textup{(F3)} until it is
converted into an actual target-retaining successor or shown to be
target-null.
\end{proposition}

\begin{proof}
Apply \cref{prop:backward-concentration-remainder}.  Its concentrated output
is \textup{(T1)}, and its homogeneous output is \textup{(T2)}.  On the
diffuse output, apply \cref{prop:first-bad-scale-descendant}; alternatives
\textup{(a)} and \textup{(b)} there give \textup{(T1)}, while alternative
\textup{(c)} gives \textup{(T3)}.  Those alternatives are exhaustive under
the hypotheses of that proposition.  Accessibility of the pressure tail follows from
\cref{prop:harmonic-pressure-tail-accessibility}.  None of these steps
replaces a tail coordinate by an unproved velocity limit.
\end{proof}

\begin{definition}[Collapsed same-origin residual]
\label{def:collapsed-same-origin-residual}
A collapsed same-origin residual is an output of \textup{(F4)} which
\begin{enumerate}[label=\textup{(J\arabic*)},leftmargin=2.8em]
\item retains a fixed positive quotient activity on a cofinal sequence of
shrinking exact parabolic cells;
\item satisfies the Navier--Stokes equation, incompressibility, the compatible
pressure equations, and the signed suitable local-energy identities on every
finite block;
\item has no terminal \(L^3\) endpoint, no homogeneous quotient, and no
positive-density compact covariant component retaining the target;
\item retains every nonzero Reynolds, pressure, scale, exterior, trace, and
realization coordinate as a first-failure successor; and
\item is not a uniformly record-bounded separated marked cascade, since that
case is excluded by
\cref{thm:no-bounded-marked-regenerated-cascade}; and
\item remains in the normalized window class
\cref{def:normalized-target-window} with the fixed bounds \(\mathbf M\), and
on the record-based realization retains the Type~I pointwise envelope
\eqref{eq:local-typeI-envelope} on every normalized domain used by the marked
projective family.  By \cref{lem:typeI-envelope-scale-escape}, loss of that
envelope is a scale/Type~II coordinate in \textup{(J4)}, while loss of any
other normalized bound is an observer, pressure, or compactness coordinate;
neither is a member of the residual.
\end{enumerate}
Thus the residual consists precisely of amplitude-renormalizing scale
collapse together with any current needed to preserve the target through its
successive changes of observer.
\end{definition}

\begin{lemma}[Type~I envelope versus scale escape]
\label{lem:typeI-envelope-scale-escape}
Let \(C_j=Q_{r_j}(z_j)\) be the exact cells in an ordinary same-origin path,
and fix \(\kappa>1\) so that \(Q_{\kappa r_j}(z_j)\) lies in the realizing
domain.  After passage to a subsequence, exactly one of the following holds:
\begin{enumerate}[label=\textup{(I\arabic*)},leftmargin=2.8em]
\item the scale-invariant pointwise envelope
\begin{equation}\label{eq:local-typeI-envelope}
 \sup_j r_j\|u\|_{L^\infty(Q_{\kappa r_j}(z_j))}<\infty
\end{equation}
holds on every normalized domain used by the marked projective family;
\item there are cells for which
\begin{equation}\label{eq:local-typeII-scale-escape}
 r_j\|u\|_{L^\infty(Q_{\kappa r_j}(z_j))}\longrightarrow\infty.
\end{equation}
The second alternative is the scale/Type~II output in \textup{(F3)}.  It is
selected before entry into the collapsed residual.
\end{enumerate}
\end{lemma}

\begin{proof}
Apply the bounded/unbounded dichotomy to the nonnegative sequence
\(r_j\|u\|_{L^\infty(Q_{\kappa r_j}(z_j))}\).  In the bounded case, parabolic
change of variables gives one common \(L^\infty\) bound on the normalized
cells.  In the unbounded case, select a subsequence tending to infinity and
retain the point, scale, and realizing pressure atlas.  This is precisely a
failure of the scale envelope in the \(\mathrm{sc}\) coordinate of
\eqref{eq:five-PDE-source-families}, hence \textup{(F3)}.  The ordered
finite-depth alternative selects that failure before the ordinary output
\textup{(F4)}.
\end{proof}

\begin{corollary}[Reduction to one explicit residual mechanism]
\label{cor:reduction-to-collapsed-residual}
For the finite-depth event construction, target loss, the compact
nonexpanding output, finite ordinary descendant chains, and uniformly record-bounded
separated cascades are closed by
\cref{prop:compact-nonexpanding-output-closure,cor:finite-descendant-chains-close,thm:no-bounded-marked-regenerated-cascade}.
Every non-tight or nonordinary failure is retained by
\cref{thm:finite-depth-NS-source-alternative,prop:first-unbounded-activity-scale}.
Consequently the only output not closed by those results is the collapsed
same-origin residual of
\cref{def:collapsed-same-origin-residual}, including any infinite chain of its
typed first-failure successors.  Its event-indexed closure is proved in
\cref{thm:collapsed-same-origin-residual-exclusion}.
\end{corollary}

\begin{proof}
Apply the alternatives in
\cref{thm:finite-depth-NS-source-alternative}.  Alternative \textup{(F1)} is
target-null.  Alternative \textup{(F2)} is impossible by countable
additivity of its finite dominating physical measure.  Alternative
\textup{(F3)} is treated by
\cref{prop:first-unbounded-activity-scale}; it is not discarded.  In
\textup{(F4)}, compact nonexpanding recurrence is excluded by
\cref{prop:compact-nonexpanding-output-closure}, finite chains by
\cref{cor:finite-descendant-chains-close}, and the uniformly record-bounded
separated value by
\cref{thm:no-bounded-marked-regenerated-cascade}.  Apply
\cref{lem:active-event-nesting-antichain} to the remaining realized active
events.  Unbounded nesting gives the longitudinal path in
\cref{def:collapsed-same-origin-residual}; large antichains are reduced by
\cref{thm:transverse-active-family-reduction} to a first typed output or to a
coalescent marked path governed by the same signed-current identity.  Removing
these disjoint subcases leaves exactly the properties in
\cref{def:collapsed-same-origin-residual}; a non-affine intermediate pressure
source remains one of its typed successor coordinates.
\end{proof}

The shell layer identifies the signed current which finances repeated progress
toward the singular target.  It is kept in full: pressure, convection,
translation, dilation, cutoff motion, kinetic storage, dissipation, and the
suitable-energy defect are never separated into unrelated estimates.

\begin{theorem}[Persistent-current compactification]
\label{thm:persistent-current-compactification}
Let the canonical orbit of \cref{thm:canonical-critical-orbit} contain
infinitely many target-retaining unit windows, and let their complete signed
currents be defined by \eqref{eq:signed-current}.  If the cumulative current
is not sublinear in the number of selected windows, then, after one ordered
subsequence selection, exactly one of the following occurs:
\begin{enumerate}[label=\textup{(\roman*)},leftmargin=2.2em]
\item a first noncompact center, pressure, scale, exterior, trace, or
weak-product coordinate produces a sequence-realized critical-tail datum;
\item every finite coordinate is compact, the successor observers and their
oriented currents converge as marks, and recurrence of the marked successor
path produces a nonzero sequence-realized logarithmic critical tail.
\end{enumerate}
In both cases the tail is realized by exact translations, Galilean changes,
and parabolic rescalings of the original solution.  No diffuse measure or
pressure remainder is declared to be an ordinary velocity profile.
\end{theorem}

\begin{proof}
Apply the finite-depth local-energy decomposition to longer and longer initial
blocks of the canonical orbit.  Target loss contradicts the retained
oscillation floor, while infinitely many disjoint positive charges contradict
the finiteness of the local dissipation and defect measures.  If a compact
coordinate first fails, its first detecting test and physical realization
survive the diagonal selection and give alternative~\textup{(i)}.

Otherwise the complete event records are compact.  The nesting--antichain
alternative reduces transverse families to a first failed coordinate or a
coalescent marked path; longitudinal nesting gives the same marked path
directly.  The exact moving-root identity excludes sublinear current.  For
persistent current, the event-shift closure retains both activity and the
oriented boundary current.  Recurrence of this compact marked path gives the
sequence-realized tail in alternative~\textup{(ii)}.  The complete derivation,
including the sparse moving-root case, is given in
\cref{lem:active-event-nesting-antichain,thm:transverse-active-family-reduction,lem:compact-closed-marked-path-space,thm:sparse-moving-root-path-exclusion}.
\end{proof}

\begin{theorem}[Critical-tail rigidity]
\label{thm:critical-tail-rigidity}
Every sequence-realized logarithmic critical tail produced by persistent
signed current is target-null.  This includes diffuse, homogeneous,
non-Dirac defect, log-periodic, and aperiodic tails.
\end{theorem}

\begin{proof}
Decompose the tail on logarithmic shells and retain every compact shell
carrying positive quotient activity.  Shell-support realization turns every
support point into an actual translated-cylinder limit of the original
solution.  Failure of logarithmic tightness is the diffuse alternative; on a
tight shell, local compactness either records the first pressure, observer,
Young, or viscous-defect coordinate, or gives strong \(L^3\) convergence.
Every recorded coordinate is rescaled at its first visible parabolic scale and
therefore returns to the same marked-current alternative, whose sparse case is
closed by the exact critical-currency identity.  If no coordinate survives,
the tail is coherent.  Bounded-origin decay then makes homogeneous,
log-periodic, and aperiodic minimal coherent tails zero.  In every case the
quotient activity vanishes.  The detailed exhaustion is proved in
\cref{lem:logarithmic-shell-support-realization,lem:realized-shell-strong-compactness,lem:concentration-free-tail-closure,lem:bounded-origin-coherent-tail-rigidity,thm:sparse-moving-root-path-exclusion}.
\end{proof}

\section{Covariant adjoint geometry and the moving-cutoff energy}
\label{sec:corrected-entropy}

The adjoint entropy \(\mathcal W_{\mathrm{ad}}\), the kinetic variance
\(\OscEnergy_\Gamma\), and the ray defect \(\RayDefect\) have distinct
roles.  The first controls the localization geometry, the second detects the
retained velocity core, and the third measures the component of viscous
motion transverse to the target-gradient ray.  None of them separately has
the monotonicity and coercivity needed below.

\subsection{The explicit nonnegative components}

For a smooth positive adjoint density, define the Gaussian production
\begin{equation}\label{eq:gaussian-production}
 \mathfrak D_{\mathrm G}(\rho)
 :=2\nu^2\theta\int_{B_R}\chi^2
 \left|\nabla^2f-\frac{1}{2\nu\theta}I\right|^2\rho\dd y.
\end{equation}
Define also the localized viscous production
\begin{equation}\label{eq:localized-viscous-production}
 \mathfrak D_{\nu,\Gamma}(V)
 :=\nu\int_{B_R}m|\nabla V|^2\dd y
\end{equation}
and the instantaneous ray production
\begin{equation}\label{eq:instantaneous-ray-production}
 \mathfrak D_{\mathrm{ray},\Gamma}(V)
 :=\widehat{\mathfrak D}_{\mathrm{ray},\Gamma}(V),
\end{equation}
where the hatted, scale-critical density is defined by
\eqref{eq:critical-ray-radon-nikodym}.
At weak regularity \(\mathfrak D_{\nu,\Gamma}\) is defined directly from the
weighted Dirichlet integral.  If the Gaussian or ray terms are retained in a
weak production measure, their identification and compact stability require
additional adjoint compactness beyond the localized energy argument; neither is inferred merely
from local weak velocity convergence.  The coercive floor used in the
compactness argument is independent of the adjoint density:
\begin{equation}\label{eq:explicit-production-floor}
 \mathfrak D_{0,\zeta}(V)
 :=\nu\int_{\R^3}\zeta^2|\nabla V|^2\dd y.
\end{equation}
The Gaussian production remains a favorable nonnegative term in the smooth
Perelman identity, but it may be discarded before taking a weak limit.  The
ray production is likewise an additional nonnegative covariant diagnostic on
each finite-energy representative.  Neither term is needed to separate
retained compact limits from equality.  Inserting the ray term into
\eqref{eq:explicit-production-floor} would in addition require a localization
theorem for the whole-space Helmholtz projection under merely local blow-up
convergence.

Formula \eqref{eq:adjoint-entropy-repaired-drift} contains a raw
velocity-strain term with no separate sign.  It becomes
coercive after coupling it to the viscous and Gaussian terms in
\cref{prop:strain-gaussian-completion}.  The moving-cutoff kinetic energy below
handles pressure and cutoff transport directly and dominates
\eqref{eq:explicit-production-floor} because \(\Phi=1\) on
\(\supp\zeta\).

\subsection{The uncoupled entropy balance}

The whole-space smooth calculation identifies every term in the correction.
Put
\begin{equation}\label{eq:repaired-material-operator}
 \mathcal L_cF:=\partial_\tau F-\nu\Delta F+c\cdot\nabla F,
 \qquad c=V-ay-b.
\end{equation}
For the kinetic density \(e=|V|^2/2\), the repaired equation gives
\begin{equation}\label{eq:kinetic-density-balance}
 \mathcal L_ce
 =-\nu|\nabla V|^2+2ae-\diver(PV).
\end{equation}
Indeed, after moving the dilation and translation derivatives to the left,
\eqref{eq:renormalized-NS} becomes
\[
 \partial_\tau V+c\cdot\nabla V+\nabla P
 =\nu\Delta V+aV.
\]
Taking its scalar product with \(V\) proves
\eqref{eq:kinetic-density-balance}.

\begin{proposition}[Scale-covariant kinetic--entropy balance]
\label{prop:uncoupled-entropy-balance}
Let \((V,P)\) and a positive probability density \(\rho\) be smooth on
\(\R^3\), with sufficient decay, and suppose that \(\rho\) solves the
unnormalized whole-space adjoint equation.  Set
\[
 \bar V_\rho(\tau):=\int_{\R^3}V\rho\dd y,
 \qquad
 \mathcal K_\rho^\circ(\tau)
 :=\frac12\int_{\R^3}|V-\bar V_\rho|^2\rho\dd y.
\]
Then
\begin{equation}\label{eq:weighted-kinetic-balance}
 \frac{\dd}{\dd\tau}\mathcal K_\rho^\circ
 =-\nu\int|\nabla V|^2\rho
   +2a\mathcal K_\rho^\circ
   -\int P(V-\bar V_\rho)\cdot\nabla f\,\rho.
\end{equation}
Let \(\theta\) solve \eqref{eq:covariant-heat-scale}.  For every
\(\beta>0\), define
\begin{equation}\label{eq:scale-covariant-functional}
 \mathcal G_\beta
 :=\theta\mathcal K_\rho^\circ-\beta\mathcal W_{\mathrm{ad}}.
\end{equation}
Then
\begin{align}
 \frac{\dd}{\dd\tau}\mathcal G_\beta
 ={}&-\mathcal K_\rho^\circ
 -\theta\nu\int|\nabla V|^2\rho
 -2\beta\nu^2\theta\int
 \left|\nabla^2f-\frac{1}{2\nu\theta}I\right|^2\rho
 \notag\\
 &-\theta\int P(V-\bar V_\rho)\cdot\nabla f\,\rho
 +2\beta\nu\theta\int
 (\nabla f)^T(\nabla V)_{\mathrm{sym}}
 \nabla f\,\rho.
 \label{eq:uncoupled-entropy-balance}
\end{align}
\end{proposition}

\begin{proof}
Apply the whole-space form of \cref{lem:adjoint-transport} componentwise to
the represented velocity equation.  Since
\(\mathcal L_cV=aV-\nabla P\),
\begin{equation}\label{eq:weighted-mean-momentum}
 \bar V_\rho'=a\bar V_\rho-
 \int_{\R^3}\nabla P\,\rho\dd y.
\end{equation}
Applying the same transport identity to
\eqref{eq:kinetic-density-balance} and integrating by parts gives
\[
 -\int\diver(PV)\rho
 =\int PV\cdot\nabla\rho
 =-\int PV\cdot\nabla f\,\rho.
\]
Subtracting \(\frac12|\bar V_\rho|^2\) and using
\eqref{eq:weighted-mean-momentum} proves
\eqref{eq:weighted-kinetic-balance}.  In particular,
\[
 (\theta\mathcal K_\rho^\circ)'
 =-\mathcal K_\rho^\circ
  -\theta\nu\int|\nabla V|^2\rho
  -\theta\int P(V-\bar V_\rho)\cdot\nabla f\,\rho,
\]
because \(\theta'+2a\theta=-1\).  Subtracting \(\beta\) times
\eqref{eq:adjoint-entropy-repaired-drift} proves
\eqref{eq:uncoupled-entropy-balance}.
\end{proof}

Thus the covariant heat scale removes the scalar modulation from the kinetic
variance as well as from the adjoint entropy.  No sign or integrability
condition on \(a\) is used in this cancellation.  The strain term in
\eqref{eq:uncoupled-entropy-balance} combines exactly with the viscous and
Gaussian terms.

\begin{proposition}[Strain--Gaussian square completion]
\label{prop:strain-gaussian-completion}
Under the hypotheses of \cref{prop:uncoupled-entropy-balance}, set
\[
 S:=\frac12\bigl(\nabla V+(\nabla V)^T\bigr),
 \qquad
 \Omega:=\frac12\bigl(\nabla V-(\nabla V)^T\bigr),
 \qquad
 A_f:=\nabla^2f-\frac{1}{2\nu\theta}I.
\]
Then
\begin{equation}\label{eq:strain-hessian-transmutation}
 \int_{\R^3}\rho\,S:(\nabla f\otimes\nabla f)\dd y
 =\int_{\R^3}\rho\,S:A_f\dd y.
\end{equation}
Consequently the entropy balance has the exact form
\begin{align}
 \frac{\dd}{\dd\tau}\mathcal G_\beta
 +\mathcal K_\rho^\circ
 +\mathfrak Q_{\beta,\theta}(V,f)
 ={}&-\theta\int_{\R^3}
 P(V-\bar V_\rho)\cdot\nabla f\,\rho\dd y,
 \label{eq:completed-entropy-balance}
\end{align}
where
\begin{align}
 \mathfrak Q_{\beta,\theta}(V,f)
 :={}&\theta\nu\int_{\R^3}|\Omega|^2\rho\dd y
 +\theta\left(\nu-\frac{\beta}{2}\right)
   \int_{\R^3}|S|^2\rho\dd y
\notag\\
&+2\beta\nu^2\theta\int_{\R^3}
 \left|A_f-\frac{1}{2\nu}S\right|^2\rho\dd y.
 \label{eq:coupled-coercive-form}
\end{align}
If \(0<\beta\le\nu\), then
\begin{equation}\label{eq:coupled-dirichlet-floor}
 \mathfrak Q_{\beta,\theta}(V,f)
 \ge \frac{\theta\nu}{2}\int_{\R^3}|\nabla V|^2\rho\dd y.
\end{equation}
In particular, the fixed choice
\begin{equation}\label{eq:beta-smallness}
 0<\beta\le\nu
\end{equation}
ensures \eqref{eq:coupled-dirichlet-floor}.  On every retained window,
\eqref{eq:theta-buffer} further gives
\begin{equation}\label{eq:coupled-unweighted-dirichlet-floor}
 \mathfrak Q_{\beta,\theta}(V,f)
 \ge \frac{\nu\kappa_*e^{-2M_*}}{2}
       \int_{\R^3}|\nabla V|^2\rho\dd y.
\end{equation}
\end{proposition}

\begin{proof}
Since \(\rho=(4\pi\nu\theta)^{-3/2}e^{-f}\),
\begin{equation}\label{eq:rho-hessian-identity}
 \nabla^2\rho
 =\rho\bigl(\nabla f\otimes\nabla f-\nabla^2f\bigr).
\end{equation}
Incompressibility gives
\[
 \partial_i\partial_jS_{ij}
 =\Delta\diver V=0.
\]
The decay hypotheses therefore justify two integrations by parts and give
\[
 \int_{\R^3}S:\nabla^2\rho\dd y
 =\int_{\R^3}\rho\,\partial_i\partial_jS_{ij}\dd y=0.
\]
Using \eqref{eq:rho-hessian-identity} and
\(\operatorname{tr}S=\diver V=0\) yields
\begin{align*}
 \int\rho\,S:(\nabla f\otimes\nabla f)
 &=\int\rho\,S:\nabla^2f
 =\int\rho\,S:A_f.
\end{align*}
This proves \eqref{eq:strain-hessian-transmutation}.  Substitution into
\eqref{eq:uncoupled-entropy-balance} gives
\eqref{eq:completed-entropy-balance}.  Finally,
\begin{align*}
 &\theta\nu|\nabla V|^2+2\beta\nu^2\theta|A_f|^2
       -2\beta\nu\theta S:A_f\\
 &\qquad=\theta\nu|\Omega|^2
  +\theta\left(\nu-\frac{\beta}{2}\right)|S|^2
  +2\beta\nu^2\theta
       \left|A_f-\frac{S}{2\nu}\right|^2,
\end{align*}
because \(|\nabla V|^2=|S|^2+|\Omega|^2\).  This proves
\eqref{eq:coupled-coercive-form}.  When \(\beta\le\nu\), the first
two terms on its right control
\((\theta\nu/2)(|S|^2+|\Omega|^2)\), proving
\eqref{eq:coupled-dirichlet-floor}.  Combining this with the lower bound in
\eqref{eq:theta-buffer} proves
\eqref{eq:coupled-unweighted-dirichlet-floor}.
\end{proof}

\begin{lemma}[Localized strain transmutation]
\label{lem:localized-strain-transmutation}
Let \(V,\rho,f\) be smooth on \(B_R\), with \(\diver V=0\), let
\(\chi\in C_c^\infty(B_R)\), and assume
\(\rho=(4\pi\nu\theta)^{-3/2}e^{-f}\) and
\(\int_{B_R}\chi^2\rho=1\).  With \(A_f\) as in
\cref{prop:strain-gaussian-completion},
\begin{align}
 \int_{B_R}\chi^2\rho\,S:(\nabla f\otimes\nabla f)\dd y
 ={}&\int_{B_R}\chi^2\rho\,S:A_f\dd y
 +\mathfrak C_\chi(S,\rho),
 \label{eq:localized-strain-transmutation}\\
 \mathfrak C_\chi(S,\rho)
 :={}&-\int_{B_R}\rho\,S:\nabla^2(\chi^2)\dd y
 -2\int_{B_R}S:
       \bigl(\nabla(\chi^2)\otimes\nabla\rho\bigr)\dd y.
 \label{eq:strain-cutoff-commutator}
\end{align}
Writing
\(\mathcal A_\chi
 :=\supp\bigl(|\nabla\chi|+|\nabla^2\chi|\bigr)\), one has
for every \(\widehat\chi\in C_c^\infty(B_R)\) satisfying
\(\widehat\chi=1\) on a neighborhood of \(\mathcal A_\chi\),
\begin{align}
 |\mathfrak C_\chi(S,\rho)|
 \le{}&\|\nabla^2(\chi^2)\|_\infty
       \int_{\mathcal A_\chi}\rho|S|\dd y
 \notag\\
 &+2\|\nabla(\chi^2)\|_\infty
 \left(\int_{\mathcal A_\chi}\rho|S|^2\dd y\right)^{1/2}
 \left(\int_{B_R}\widehat\chi^2
       \frac{|\nabla\rho|^2}{\rho}\dd y\right)^{1/2}
 \label{eq:strain-cutoff-estimate}
\end{align}
 Moreover, with \(m=\chi^2\rho\), the localized quadratic form
 \begin{align}
  \mathfrak Q_{\beta,\theta}^{\chi}(V,f)
 :={}&\theta\nu\int_{B_R}m|\Omega|^2\dd y
 +\theta\left(\nu-\frac{\beta}{2}\right)
   \int_{B_R}m|S|^2\dd y
 \notag\\
 &+2\beta\nu^2\theta\int_{B_R}m
       \left|A_f-\frac{S}{2\nu}\right|^2\dd y
 \label{eq:localized-coupled-form}
\end{align}
satisfies, whenever \(0<\beta\le\nu\),
\begin{equation}\label{eq:localized-coupled-floor}
 \mathfrak Q_{\beta,\theta}^{\chi}(V,f)
 \ge\frac{\theta\nu}{2}\int_{B_R}m|\nabla V|^2\dd y
 =\frac{\theta}{2}\mathfrak D_{\nu,\Gamma}(V).
\end{equation}
\end{lemma}

\begin{proof}
Identity \eqref{eq:rho-hessian-identity} holds on \(B_R\).  Moreover
\(\partial_i\partial_jS_{ij}=0\), as in the proof of
\cref{prop:strain-gaussian-completion}.  Testing this distributional identity
against the compactly supported function \(\chi^2\rho\) gives
\[
 0=\int_{B_R}S:\nabla^2(\chi^2\rho)\dd y.
\]
Expanding the Hessian and using the symmetry of \(S\) yields
\[
 \int_{B_R}\chi^2S:\nabla^2\rho\dd y
 =-\int_{B_R}\rho S:\nabla^2(\chi^2)\dd y
  -2\int_{B_R}S:
       \bigl(\nabla(\chi^2)\otimes\nabla\rho\bigr)\dd y.
\]
Substitute \eqref{eq:rho-hessian-identity}, then use
\(\operatorname{tr}S=0\).  This proves
\eqref{eq:localized-strain-transmutation} and
\eqref{eq:strain-cutoff-commutator}.  The first term in
\eqref{eq:strain-cutoff-estimate} follows directly from support and the
\(L^\infty\) norms of the cutoff derivatives.  For the second, use
\[
 |S|\,|\nabla\rho|
 =(\sqrt\rho\,|S|)
   \left(\frac{|\nabla\rho|}{\sqrt\rho}\right)
\]
and Cauchy--Schwarz on \(\mathcal A_\chi\).
The last integral over \(\mathcal A_\chi\) is bounded by the displayed
\(\widehat\chi\)-weighted Fisher information because
\(\widehat\chi=1\) there.
The identity
\(|\nabla V|^2=|S|^2+|\Omega|^2\) and the same pointwise square completion
as in \cref{prop:strain-gaussian-completion} prove
\eqref{eq:localized-coupled-form}--\eqref{eq:localized-coupled-floor}.
\end{proof}

Thus velocity strain requires no independent correction in the interior
smooth identity.  The only localized discrepancy is the annular commutator
\(\mathfrak C_\chi\).  The scalar scale term has disappeared exactly from
\eqref{eq:completed-entropy-balance}.  The sole unsigned interior term in the
whole-space scale-covariant identity is the pressure transfer
\[
 -\theta\int P(V-\bar V_\rho)\cdot\nabla f\,\rho.
\]
For a localized adjoint density there are, in addition, the explicitly
displayed cutoff and normalization terms from
\eqref{eq:adjoint-transport}.  We now compare this structural identity with
the local kinetic-energy balance, where pressure occurs as a signed cutoff
flux.

In particular, the adjoint pairing
\(-\theta\int P(V-\bar V_\rho)\cdot\nabla f\,\rho\) is not identified with the moving-cutoff
flux \(\int PV\cdot\nabla\Phi\).  The pressure-tail estimates apply to the
latter.  The proof therefore discards the adjoint balance before invoking
monotonicity; no cancellation between these two distinct pressure terms is
claimed.

\subsection{The moving-cutoff kinetic energy}

Fix a radial nonincreasing \(\phi\in C_c^\infty(B_2)\), with
\(0\le\phi\le1\) and \(\phi=1\) on \(B_1\).  Let
\(R:[\tau_0,\infty)\to[R_*,\infty)\) be locally absolutely continuous and
nondecreasing, and put
\begin{equation}\label{eq:moving-cutoff-cutoff}
 \Phi(\tau,y):=\phi(y/R(\tau)),\qquad
 \mathcal E_R(\tau):=\frac12\int_{\R^3}|V|^2\Phi\dd y.
\end{equation}
Write \(a_+=\max\{a,0\}\), \(a_-=\max\{-a,0\}\), and define
\begin{equation}\label{eq:moving-cutoff-dissipation}
 \widetilde{\mathfrak D}_R(\tau)
 :=\nu\int|\nabla V|^2\Phi\dd y
   +a_+(\tau)\int|V|^2\Phi\dd y.
\end{equation}
The full unsigned boundary-flux density is
\begin{align}
 B_{\mathrm{mov}}(\tau):={}&
 \frac12\left|\int |V|^2\partial_\tau\Phi\right|
 +\frac\nu2\left|\int |V|^2\Delta\Phi\right|
 +\frac12\left|\int |V|^2V\cdot\nabla\Phi\right|
 +\left|\int PV\cdot\nabla\Phi\right| \notag\\
 &+\frac12|b(\tau)|\left|\int |V|^2\nabla\Phi\right|
 +\frac12a_-(\tau)\int |V|^2\Phi
 +\frac12a_+(\tau)\int |V|^2(-y\cdot\nabla\Phi).
 \label{eq:moving-cutoff-error}
\end{align}
Every integral in \eqref{eq:moving-cutoff-dissipation}--
\eqref{eq:moving-cutoff-error} is over \(\R^3\).  The pressure term is
unchanged by adding a function of time to \(P\), because
\(\int V\cdot\nabla\Phi=0\).

\begin{theorem}[Corrected moving-cutoff monotonicity]
\label{thm:moving-cutoff-monotonicity}
Let \((V,P)\) be a represented suitable weak solution of
\eqref{eq:renormalized-NS}.  If \(B_{\mathrm{mov}}\in
L^1(\tau_0,\infty)\), then
\begin{equation}\label{eq:moving-cutoff-entropy-definition}
 \EntropyNS(\tau):=\mathcal E_R(\tau)
   +\int_\tau^\infty B_{\mathrm{mov}}(s)\dd s
\end{equation}
has a nonincreasing representative and, in distributions,
\begin{equation}\label{eq:moving-cutoff-entropy-production}
 \frac{\dd}{\dd\tau}\EntropyNS
 +\frac12\widetilde{\mathfrak D}_R\le0.
\end{equation}
In particular \(\EntropyNS\ge0\), and for \(\tau_0\le s<t\),
\begin{equation}\label{eq:global-entropy-error-balance}
 \EntropyNS(t)+\ProductionNS((s,t])\le\EntropyNS(s),
 \qquad
 \ProductionNS((s,t]):=\frac12\int_s^t
 \widetilde{\mathfrak D}_R(\tau)\dd\tau.
\end{equation}
\end{theorem}

\begin{proof}
Testing the suitable local energy inequality with \(\Phi\) gives
\begin{align*}
 \frac{\dd}{\dd\tau}\mathcal E_R
 \le{}&-\nu\int|\nabla V|^2\Phi
 +\frac12\int|V|^2\partial_\tau\Phi
 +\frac\nu2\int|V|^2\Delta\Phi
 +\frac12\int|V|^2V\cdot\nabla\Phi\\
 &+\int PV\cdot\nabla\Phi
 -\frac a2\int|V|^2\Phi
 -\frac a2\int|V|^2y\cdot\nabla\Phi
 -\frac12b\cdot\int|V|^2\nabla\Phi .
\end{align*}
Add \(\widetilde{\mathfrak D}_R/2\).  The viscous term leaves
\(-\frac\nu2\int|\nabla V|^2\Phi\le0\), while
\(-aM/2+a_+M/2=a_-M/2\).  Since \(\phi\) is radial nonincreasing,
\(y\cdot\nabla\Phi\le0\); hence the scale-shell term is bounded above by
the last term of \eqref{eq:moving-cutoff-error}.  Taking absolute values of
the other fluxes proves
\[
 \frac{\dd}{\dd\tau}\mathcal E_R
 +\frac12\widetilde{\mathfrak D}_R\le B_{\mathrm{mov}}.
\]
The suitable energy-defect measure has the favorable sign and was discarded
in this inequality.  Differentiating the finite tail in
\eqref{eq:moving-cutoff-entropy-definition} subtracts \(B_{\mathrm{mov}}\), proving
\eqref{eq:moving-cutoff-entropy-production}.  Both terms in
\eqref{eq:moving-cutoff-entropy-definition} are nonnegative.
\end{proof}

\subsection{Finite-block use of the moving cutoff}

The proof uses \cref{thm:moving-cutoff-monotonicity} on the finite blocks
selected by the canonical observer.  No global summable-radius schedule is
needed.  The complete error \(B_{\mathrm{mov}}\) on a block is part of its
signed current \(H_j\) in \cref{thm:exact-critical-financing}; if these errors
are sublinear, the critical price closes the path, and if they persist, their
first noncompact coordinate or recurrent support is retained by
\cref{thm:persistent-current-compactification}.  Thus the adjoint
localization and the eventwise current are two expressions of the same
financing identity.

\subsection{Equality states}

\begin{definition}[Oscillation-removable equality state]
\label{def:removable-equality-state}
A represented state is oscillation-removable when its velocity is spatially
constant on the connected retained region \(\{\zeta>0\}\) for almost every
represented time.  The constant may depend on time.  Such a state has zero
Galilean-reduced oscillatory concentration.  No assertion that a localized
Dirichlet adjoint density equals the whole-space Gaussian is included in this
definition; the Gaussian remains the equality model only for the
whole-space identity \eqref{eq:adjoint-entropy-production}.
\end{definition}

\begin{theorem}[Inner-Dirichlet and zero-production rigidity]
\label{thm:zero-production-rigidity}
Let \(\mathfrak w=(V,P,a,b;[0,1])\) be a normalized window or a compact
limit of normalized windows.  If
\begin{equation}\label{eq:zero-dirichlet-production}
 \int_0^1\mathfrak D_{0,\zeta}(V(\tau))\dd\tau=0,
\end{equation}
then \(\mathfrak w\) is oscillation-removable and
\begin{equation}\label{eq:zero-state-no-oscillation}
 \OscMass_\Gamma(V;[0,1])=0.
\end{equation}
In particular, zero moving-localized dissipation on the window implies
\eqref{eq:zero-state-no-oscillation} whenever \(\Phi=1\) on
\(\supp\zeta\).
\end{theorem}

\begin{proof}
Since \(\nu>0\), \eqref{eq:zero-dirichlet-production} gives
\(\nabla V=0\) almost everywhere
on the connected set \(\{\zeta>0\}\) for almost every \(\tau\).  Hence
\(V(\cdot,\tau)=c(\tau)\) there.  Its \(\zeta^3\)-weighted mean is
the same vector \(c(\tau)\), so the integrand in
\eqref{eq:oscillatory-mass} vanishes.  This proves
\eqref{eq:zero-state-no-oscillation}.  The final assertion follows from
the bound
\(\ProductionNS(I)\ge\frac12\int_I\mathfrak D_{0,\zeta}\) and
nonnegativity.
\end{proof}

\section{The positive critical price}
\label{sec:entropy-gap}

This section identifies the abstract covariant dissipation of
\cref{sec:dissipation-reduction} with the moving-localized dissipation.  It is
a concrete realization of the structural balance.  No numerical estimate for
the gap is needed: compactness upgrades the qualitative equality theorem to a
uniform statement.

\begin{lemma}[Compactness--rigidity gap principle]
\label{lem:abstract-gap-principle}
Let \(\mathcal K\) be a sequentially compact space, let
\(\mathcal B\subset\mathcal K\) be sequentially closed, and let
\(\mathfrak d:\mathcal K\to[0,\infty]\) be sequentially lower
semicontinuous.  If
\begin{equation}\label{eq:abstract-gap-separation}
 \{\mathfrak d=0\}\cap\mathcal B=\varnothing,
\end{equation}
then
\begin{equation}\label{eq:abstract-positive-gap}
 \inf_{x\in\mathcal B}\mathfrak d(x)>0.
\end{equation}
\end{lemma}

\begin{proof}
If the infimum were zero, there would be \(x_j\in\mathcal B\) with
\(\mathfrak d(x_j)\to0\).  Sequential compactness gives
\(x_j\to x_\infty\) after extraction.  Closedness gives
\(x_\infty\in\mathcal B\), while lower semicontinuity gives
\(\mathfrak d(x_\infty)=0\), contradicting
\eqref{eq:abstract-gap-separation}.
\end{proof}

For the Navier--Stokes application, \(\mathcal K\) is the sequential closure
of the selected windows from one canonical orbit in the full topology of
clause~\textup{(iii)} of \cref{thm:canonical-critical-orbit};
\(\mathcal B\) is the subset satisfying
\eqref{eq:target-window-floor}; and
\begin{equation}\label{eq:dirichlet-gap-functional}
 \mathfrak d_\nu(\mathfrak w)
 :=\int_0^1\mathfrak D_{0,\zeta}(V(\tau))\dd\tau.
\end{equation}
Strong local \(L^3\) convergence makes \(\mathcal B\) closed, while
\eqref{eq:inner-production-lsc} makes \(\mathfrak d_\nu\) lower
semicontinuous.  Both statements are proved in
\cref{prop:entropy-compactification}.  The first assertion of
\cref{thm:zero-production-rigidity} supplies
\eqref{eq:abstract-gap-separation}.  Thus the compactness gap uses no
semicontinuity assertion about the corrected localized energy, the Gaussian Hessian,
or the nonlocal ray distance.

\begin{theorem}[Positive critical dissipation price]
\label{thm:uniform-production-gap}
Fix the canonical orbit generated by one putative terminal loss.  There
exists \(d_*=d_*(\text{orbit})>0\) such that every selected
target-retaining window satisfies
\begin{equation}\label{eq:uniform-production-gap}
 D_j^{\rm act}:=\nu\int_{I_j}\int_{\R^3}\zeta^2|\nabla V|^2
 \dd y\dd\tau\ge d_*.
\end{equation}
\end{theorem}

\begin{proof}
Apply \cref{lem:abstract-gap-principle} to the closure of the selected
windows with \(\mathfrak d=\mathfrak d_\nu\), as described above.  It gives
a number \(\delta_\nu>0\), depending on this orbit, such that
\(\mathfrak d_\nu(\mathfrak w_j)\ge\delta_\nu\) for every selected window.
By \eqref{eq:explicit-production-floor}, the left side of
\eqref{eq:uniform-production-gap} is
\(\mathfrak d_\nu(\mathfrak w_j)\).  Take \(d_*=\delta_\nu\).
\end{proof}

\begin{corollary}[Infinite normalized expenditure]
\label{cor:retained-target-production-divergence}
For the infinite pairwise disjoint selected windows of the canonical orbit,
\begin{equation}\label{eq:retained-target-production-divergence}
 \sum_{j=1}^{\infty}D_j^{\rm act}=\infty.
\end{equation}
More generally, any nonnegative countably additive measure \(\mu\) satisfying
\eqref{eq:critical-measure-dominates-dirichlet} on those windows obeys
\(\sum_j\mu(I_j)=\infty\).
\end{corollary}

\begin{proof}
The first conclusion follows from
\eqref{eq:uniform-production-gap}.  The general conclusion is
\cref{cor:target-visible-dissipation-divergence}.  Only the component selected
by the retained target is used; no estimate of target-null transport or
material-frame motion enters the sum.
\end{proof}

\begin{corollary}[One-window entropy estimate]
\label{cor:one-window-entropy-drop}
If the moving-cutoff entropy satisfies
\(\ProductionNS(I_j)\ge\tfrac12D_j^{\rm act}\), every selected window satisfies
\begin{equation}\label{eq:one-window-entropy-drop}
 \EntropyNS(\tau_j+1)
 \le \EntropyNS(\tau_j)-\tfrac12d_*.
\end{equation}
\end{corollary}

\begin{proof}
Apply \eqref{eq:global-entropy-error-balance} on \(I_j\) and then use
\eqref{eq:uniform-production-gap}.
\end{proof}

\section{Covariant exclusion of singularities}
\label{sec:global-exclusion}

\begin{proof}[Proof of \cref{thm:intro-main}]
Let \(T_*\) be the maximal classical existence time and suppose that
\(T_*<\infty\).  Choose a global finite-energy suitable Leray solution.  By
weak--strong uniqueness it agrees with the classical solution on
\([0,T_*)\) and gives a suitable continuation through \(T_*\).

Canonical observation, \cref{thm:canonical-critical-orbit}, supplies one
coherent scale--center path and pairwise disjoint unit windows
\(I_j\), all satisfying
\[
 \OscMass_\Gamma(V;I_j)\ge\eta_0.
\]
By the positive-price theorem \cref{thm:uniform-production-gap},
\[
 D_j^{\rm act}\ge d_*>0
 \qquad\text{for every }j.
\]
Apply the exact financing identity
\cref{thm:exact-critical-financing} to the first \(N\) consecutive cells:
\[
 \sum_{j<N}(D_j+M_j)
 =K_0^- -K_{N-1}^+ +\sum_{j<N}H_j .
\]
The kinetic endpoints are uniformly bounded.  Since each full-cell
dissipation satisfies \(D_j\ge D_j^{\rm act}\), one has
\(\sum_{j<N}(D_j+M_j)\ge Nd_*\), and the cumulative signed current cannot be
sublinear.

The remaining possibility is persistent critical current.  By
\cref{thm:persistent-current-compactification}, it either loses a first
center, pressure, scale, exterior, trace, or defect coordinate, or repeatedly
returns through active centers.  The first alternative produces a
sequence-realized critical tail at its first visible parabolic scale.  In the
second, the observer and the complete current are retained as marks; recurrence
of the marked successor path again produces a nonzero sequence-realized
logarithmic critical tail.  Both alternatives contradict
\cref{thm:critical-tail-rigidity}.  Hence \(T_*=\infty\).

Standard parabolic regularity then gives smoothness on every compact positive
time interval, and weak--strong uniqueness identifies the Leray solution with
this global classical solution.
\end{proof}

This is the entire global argument.  The appendices below prove its five
structural theorems; all local configurations enter only through
persistent-current compactification and critical-tail rigidity.

\clearpage
\appendix

\section{Weak covariant calculus and compact stability}
\label{app:weak-calculus}
\label{sec:weak-extension}

This appendix records the standard weak covariant identities and proves the
local stability statements used in the critical-price theorem.

\subsection{Regularization of the equation}

Let \(S_hv(t)=h^{-1}\int_t^{t+h}v(s)\dd s\) be a Steklov average on an
interior interval, and let \(v_\varepsilon\) denote spatial mollification.
After a diagonal choice,
\begin{equation}\label{eq:regularization-convergence}
 u_{\varepsilon,h}\to u
 \quad\text{strongly in }L^2_{\mathrm{loc}}H^1_{\mathrm{loc}},
\end{equation}
with weak-* convergence in \(L^\infty L^2\).  The nonlinear commutator is
\begin{equation}\label{eq:nonlinear-commutator}
 \mathcal C_\varepsilon
 :=\diver\bigl((u\otimes u)_\varepsilon
              -u_\varepsilon\otimes u_\varepsilon\bigr).
\end{equation}
Since \(u\in L^{8/3}_{\mathrm{loc}}L^4_{\mathrm{loc}}\),
\begin{equation}\label{eq:commutator-convergence}
 \mathcal C_\varepsilon\to0
 \quad\text{in }L^{4/3}_{\mathrm{loc}}H^{-1}_{\mathrm{loc}}.
\end{equation}
On a compact repaired interval, spatial mollification commutes with the
translation derivative, while
\begin{equation}\label{eq:dilation-commutator}
 a[\,y\cdot\nabla,\varphi_\varepsilon*\,]V\to0
 \quad\text{in }L^1_{\mathrm{loc}}L^2_{\mathrm{loc}}.
\end{equation}
The spatial commutator is uniformly \(L^2\)-bounded and converges strongly to
zero; dominated convergence uses \(a\in L^1_{\mathrm{loc}}\) and the local
\(L^\infty_\tau L^2_y\) bound.  Moreover,
\[
 \|(b_1-b_2)\cdot\nabla V\|_{H^{-1}(B_R)}
 \le c_R|b_1-b_2|\,\|V\|_{L^2(B_R)},
\]
with the analogous estimate for
\((a_1-a_2)(V+y\cdot\nabla V)\).

\begin{proposition}[Weak validity of the covariant identities]
\label{prop:weak-covariant-identities}
For every suitable weak solution:
\begin{enumerate}[label=\textup{(\roman*)},leftmargin=2.2em]
\item \(R(u)=0\) in distributions, and
\(\Dcov_\tau^{X,\Lambda}V+\nu A_{\mathrm S}V=0\) on every repaired compact
cylinder;
\item the global energy balance is \eqref{eq:weak-energy-balance};
\item the localized identity is \eqref{eq:weak-probe-balance} for every
admissible smooth compactly supported weight;
\item the Stokes-metric projection inequalities hold for almost every weak time.
\end{enumerate}
\end{proposition}

\begin{proof}
For item~\textup{(i)}, write the mollified equation with residual
\(\mathcal C_\varepsilon\).  The linear and pressure terms pass by
distributional convergence, the nonlinear residual vanishes by
\eqref{eq:commutator-convergence}, and the modulation terms pass by
\eqref{eq:dilation-commutator} and \(L^1\)-approximation of \(a,b\).
Item~\textup{(ii)} is the Leray--Hopf energy inequality in the measure form
proved in \cref{thm:weak-energy-balance}.  For item~\textup{(iii)}, test the
suitable local energy inequality by \(w_\Gamma\), and test the weak momentum
equation componentwise by the same scalar weight.  The latter identity gives
the distributional derivative of
\(M_\Gamma^{-1}|\int w_\Gamma u|^2/2\); subtracting it from the weighted
kinetic-energy relation gives exactly \eqref{eq:weak-probe-balance}, including
the same nonnegative local defect measure.  Approximation of temporal
indicators gives the asserted one-sided representatives.  Item~\textup{(iv)}
is Cauchy--Schwarz in the solenoidal Gelfand triple.
\end{proof}

\subsection{Compactness of the variational data}

Suppose, on compact cylinders,
\begin{equation}\label{eq:canonical-convergence}
 V_j\to V\quad\text{strongly in }L^2_{\mathrm{loc}}
 \text{ and }L^3_{\mathrm{loc}},
 \qquad
 \nabla V_j\rightharpoonup\nabla V\quad\text{in }L^2_{\mathrm{loc}},
\end{equation}
with the corresponding pressure and modulation convergence.  For the
additional adjoint-weighted conclusions, assume
\begin{equation}\label{eq:weight-convergence}
 m_j\to m\quad\text{strongly in }L^\infty(B_R\times I),
 \qquad m_j\ge m_*>0\quad\text{on }\supp\zeta\times I,
\end{equation}
and, for normalized adjoint densities,
\begin{equation}\label{eq:adjoint-compactness-hypotheses}
 \inf_j\inf_\tau Z_j(\tau)>0,\quad
 \sup_j\|\widetilde\rho_j\|_{L^\infty L^2\cap L^2H^1_0}<\infty,
 \quad
 \sup_j\int_I\int_{B_R}\chi_j^2
 |\nabla\sqrt{\rho_j}|^2\dd y\dd\tau<\infty.
\end{equation}
\begin{proposition}[Stability of the local variational data]
\label{prop:entropy-compactification}
Under \eqref{eq:canonical-convergence}:
\begin{enumerate}[label=\textup{(\roman*)},leftmargin=2.2em]
\item the fixed-cutoff means \(\bar V_{j,\zeta}\) converge strongly in
\(L^3(I)\), and \(\OscMass_{\Gamma_j}(V_j;I)\to
\OscMass_\Gamma(V;I)\);
\item the adjoint-weighted kinetic variances converge in \(L^1(I)\) whenever
\(m_j\to m\) strongly on the full support of the probe;
\item the fixed inner Dirichlet floor is lower semicontinuous:
\begin{equation}\label{eq:inner-production-lsc}
 \int_I\mathfrak D_{0,\zeta}(V)\dd\tau
 \le\liminf_{j\to\infty}
 \int_I\mathfrak D_{0,\zeta}(V_j)\dd\tau;
\end{equation}
if in addition \eqref{eq:weight-convergence}--
\eqref{eq:adjoint-compactness-hypotheses} hold, then the adjoint-weighted
viscous production is also lower semicontinuous:
\begin{equation}\label{eq:viscous-production-lsc}
 \int_I\mathfrak D_{\nu,\Gamma}(V)\dd\tau
 \le\liminf_{j\to\infty}
 \int_I\mathfrak D_{\nu,\Gamma_j}(V_j)\dd\tau;
\end{equation}
\item under \eqref{eq:weight-convergence}, the positivity bound for \(m\) on
\(\supp\zeta\) passes to the limit.
\end{enumerate}
No lower-semicontinuity claim for the Gaussian Hessian square, the varying-ray
distance, or the corrected production is made here, and none is used in the
gap theorem.  Retaining such terms in a stronger weak-production statement
would require identification of the adjoint logarithmic Hessian and all
associated localization terms; that refinement is not part of the moving-cutoff
argument.
\end{proposition}

\begin{proof}
The first assertion follows from strong \(L^3\) convergence and H\"older's
inequality on the fixed support of \(\zeta\).  The same argument, together
with strong convergence of the weights, gives the second assertion.  For the
third, \(\zeta\nabla V_j\rightharpoonup\zeta\nabla V\) weakly in
\(L^2\), which proves \eqref{eq:inner-production-lsc}.  Under the additional
weight convergence, \(m_j^{1/2}\nabla V_j\rightharpoonup
m^{1/2}\nabla V\) weakly in \(L^2\); weak lower semicontinuity of the
\(L^2\)-norm gives \eqref{eq:viscous-production-lsc}.  Uniform convergence of
the weights gives the final assertion.
\end{proof}

\section{Singular-germ extraction}
\label{app:critical-orbit}

This appendix derives \cref{thm:canonical-critical-orbit} directly from the
suitable local-energy inequality, pressure reconstruction, and parabolic
compactness.

\subsection{Velocity concentration}
\label{app:velocity-concentration}

The local concentration assertion was proved in
\cref{thm:velocity-concentration}.  For the global regularity problem one can
start more directly from a maximal classical solution.  The following
record-point construction removes the existence of an ancient PDE profile
from the list of assumptions.

\begin{theorem}[Record-point ancient-profile extraction]
\label{thm:record-point-ancient-extraction}
Let \(u\) be a classical finite-energy solution of \eqref{eq:NS} on
\([0,T)\), with divergence-free initial datum in \(H^m(\R^3)\),
\(m>5/2\), and suppose that \(T<\infty\) is its maximal classical time.
Then there are \(t_j\uparrow T\), \(x_j\in\R^3\), and \(r_j\downarrow0\)
such that
\begin{equation}\label{eq:record-point-choice}
 r_j^{-1}=M_j:=\sup_{0\le t\le t_j}\|u(t)\|_{L^\infty},
 \qquad |u(x_j,t_j)|\ge\frac12M_j,
\end{equation}
and the rescaled solutions
\begin{equation}\label{eq:record-point-rescaling}
 U_j(y,s):=r_j u(x_j+r_jy,t_j+r_j^2s),
 \qquad
 Q_j(y,s):=r_j^2p(x_j+r_jy,t_j+r_j^2s)
\end{equation}
are defined on \(\R^3\times(-t_j/r_j^2,0]\), satisfy
\begin{equation}\label{eq:record-profile-bound}
 \|U_j\|_{L^\infty(\R^3\times(-t_j/r_j^2,0])}\le1,
 \qquad |U_j(0,0)|\ge\frac12,
\end{equation}
and have a subsequence converging locally smoothly in velocity, and locally
in pressure modulo functions of time, to a bounded mild ancient solution
\((U,Q)\) on \(\R^3\times(-\infty,0]\).  The limit satisfies
\begin{equation}\label{eq:record-profile-nonzero}
 \|U\|_{L^\infty}\le1,
 \qquad |U(0,0)|\ge\frac12.
\end{equation}
Every further spacetime translation with translation time tending to
\(-\infty\) has, after pressure normalization and passage to a subsequence,
a bounded mild \emph{entire} Navier--Stokes limit.
\end{theorem}

\begin{proof}
The standard \(L^\infty\) continuation criterion implies
\(\limsup_{t\uparrow T}\|u(t)\|_\infty=\infty\); otherwise the mild
solution can be continued a uniform positive time beyond \(T\).  Choose
record times \(t_j\uparrow T\) for which \(M_j\to\infty\), and choose
\(x_j\) as in \eqref{eq:record-point-choice}.  Such approximate maximizers
are sufficient; for \(H^m\)-data the velocity is continuous and vanishes at
spatial infinity at every time.  Then \(r_j=M_j^{-1}\downarrow0\), and
parabolic invariance gives the equation for \((U_j,Q_j)\).  The record
property gives \eqref{eq:record-profile-bound}.  Moreover
\(t_j/r_j^2=t_jM_j^2\to\infty\).

Fix \(R,L<\infty\).  For all large \(j\), the cylinder
\(B_{2R}\times[-2L,0]\) lies in the domain of \(U_j\).  Bounded mild
estimates, applied from a slightly earlier time, give uniform interior
derivative bounds for the velocity on \(B_R\times[-L,0]\).  Equivalently,
one may combine the local energy inequality, the local
Calder\'on--Zygmund/harmonic pressure decomposition, and interior parabolic
regularity.  Normalize \(Q_j\) by its spatial mean on each member of a
countable nested cylinder exhaustion.  The local pressure bounds and a
diagonal Arzel\`a--Ascoli extraction give a smooth velocity limit and a
compatible local pressure atlas.  Passing to the mild formula proves that the
limit is mild.  The pointwise normalization passes to the limit and gives
\eqref{eq:record-profile-nonzero}.

Finally, translate \((U,Q)\) by times \(s_j\to-\infty\).  Every fixed
compact time interval then lies in the translated domain.  The same uniform
bounded-mild estimates and diagonal argument give a solution on all of
\(\R^3\times\R\).  Spatial translations and constant parabolic rescalings
commute with this extraction and preserve the assertion.
\end{proof}

The nonzero value in \eqref{eq:record-profile-nonzero} is not yet the
Galilean-invariant target: a locally constant ancient limit can carry that
normalization.  The concentration--remainder selection below removes the
homogeneous mode and chooses the first larger parabolic scale carrying fixed
positive oscillation.  Finite energy excludes a nonzero constant physical
representative on all of \(\R^3\).

\subsection{Compactness selection}
\label{app:compactness-selection}

The endpoint of a finite descendant chain is controlled by an existing
ancient-solution Liouville theorem.  We record its exact form because it is the
reason that only infinite or diffuse successor chains require the new
covariant argument.

\begin{theorem}[Endpoint ancient Liouville theorem]
\label{thm:AB-endpoint}
Let \(U\) be a mild ancient solution of the three-dimensional
Navier--Stokes equations on \(\R^3\times(-\infty,0)\).  If there are times
\(s_k\downarrow-\infty\) such that
\begin{equation}\label{eq:AB-backward-L3-sequence}
 \sup_k\|U(\cdot,s_k)\|_{L^3(\R^3)}<\infty,
\end{equation}
then \(U\equiv0\).
\end{theorem}

\begin{proof}
This is the sequence-\(L^3\) Liouville theorem of Albritton and Barker
\cite[Theorem~1.2]{AlbrittonBarker2019}.  Their statement requires mildness
but no terminal \(L^3\) hypothesis.
\end{proof}

For a bounded ancient velocity define the unit quotient activity
\begin{equation}\label{eq:unit-quotient-activity}
 \mathfrak a_U(z)
 :=\inf_{c\in\R^3}\int_{Q_1(z)}|U-c|^3\dd x\dd t,
 \qquad z=(x,t),\quad t<-1.
\end{equation}
The pressure is not part of this target.  On a bounded mild family it is
reconstructed after subtracting spatial means, and ordinary spacetime
translations are compact on every fixed cylinder whose translated times tend
to \(-\infty\).

\begin{proposition}[Backward concentration--remainder alternative]
\label{prop:backward-concentration-remainder}
Let \(U\) be a bounded mild ancient solution on
\(\R^3\times(-\infty,0)\).  After passage to subsequences, at least one of
the following conclusions can be selected; the first applicable conclusion
in the displayed order may be used.
\begin{enumerate}[label=\textup{(\roman*)},leftmargin=2.2em]
\item There are \(s_k\downarrow-\infty\) satisfying
\eqref{eq:AB-backward-L3-sequence}.
\item After removal of one spatially homogeneous mild mode, there are
\(z_k=(x_k,t_k)\), \(t_k\to-\infty\), and \(\eta>0\) such that
\(\mathfrak a_U(z_k)\ge\eta\).  The translates about \(z_k\) have a
locally smooth subsequence converging to a bounded mild entire solution
\(W\) with \(\mathfrak a_W(0,0)\ge\eta\).
\item There are \(t_k\to-\infty\), \(R_k\to\infty\), and measurable
sets \(E_k\subset\{|x|>R_k\}\) such that
\begin{equation}\label{eq:diffuse-tail-mass}
 1\le\int_{E_k}|U(x,t_k)|^3\dd x\le2,
\end{equation}
whereas every unit parabolic recentering meeting
\(E_k\times(t_k-1,t_k+1)\) has quotient activity tending to zero.  This is
the diffuse remainder.
\item The selected exterior fields converge locally, after subtraction of
time-dependent constants, to a spatially homogeneous velocity.  The momentum
equation then pairs it only with an affine pressure, so the limit is a
Galilean/parasitic quotient state and has zero target activity.
\end{enumerate}
In particular, a bounded ancient profile which is terminal under nontrivial
spacetime descendants, has no diffuse remainder, and is not a homogeneous
quotient state satisfies alternative~\textup{(i)} and is zero by
\cref{thm:AB-endpoint}.
\end{proposition}

\begin{proof}
Put \(A(t)=\int_{\R^3}|U(x,t)|^3\dd x\in[0,\infty]\).  If
\(\liminf_{t\to-\infty}A(t)<\infty\), choose a backward sequence on which
\(A\) is bounded; this is \textup{(i)}.  Otherwise
\(A(t)\to\infty\) as \(t\to-\infty\).  Since \(U\) is bounded, its
\(L^3\)-mass on every fixed ball is bounded uniformly in time.  We may
therefore choose \(t_k\to-\infty\), radii \(R_k\to\infty\), and, by
absolute continuity of the integral, exterior sets \(E_k\) satisfying
\eqref{eq:diffuse-tail-mass}.

Cover \(E_k\times(t_k-1,t_k+1)\) by a fixed bounded-overlap family of unit
parabolic cylinders.  If the supremum of the quotient activity on those
cylinders has a positive limsup, choose one such cylinder for each \(k\).
Bounded mild interior estimates give uniform velocity derivative bounds on
every compact subcylinder after recentering.  The local
Calder\'on--Zygmund/harmonic decomposition gives pressure compactness modulo
time functions.  Because \(t_k\to-\infty\), the translated time domains
exhaust \(\R\).  A diagonal extraction therefore gives an actual bounded
mild entire solution.  Strong local \(L^3\) convergence preserves the
positive quotient activity and proves \textup{(ii)}.

If the supremum tends to zero, retain the normalized exterior restrictions
in \eqref{eq:diffuse-tail-mass}.  When their quotient oscillation remains
visible only after aggregation over an increasing number of unit cylinders,
they give exactly \textup{(iii)}.  Otherwise the local Poincar\'e quotients
vanish on every cylinder in each connected overlap component.  The associated
local constants agree up to an error tending to zero; diagonal passage gives
a spatially homogeneous limit.  Substitution in the momentum equation shows
that its pressure is affine in space modulo a function of time.  This is
\textup{(iv)}.  The four outcomes exhaust the covering dichotomy.

Under the terminal assumptions in the last sentence, alternatives
\textup{(ii)--(iv)} are unavailable, so \textup{(i)} holds and
\cref{thm:AB-endpoint} applies.
\end{proof}

The proposition is deliberately phrased in terms of actual translated PDE
limits.  The diffuse alternative is measure-valued at this stage and is not
declared an ordinary solution.  The next subsections select its first scale of
nonzero quotient oscillation and retain the observer and current through an
infinite iteration.

The first part of that realization is a stopping-time argument.  For a
suitable solution \(W\), a center \(z=(x_0,t_0)\), and a constant vector
\(c\), let
\[
 Q_r^c(z):=\{(x,t):t_0-r^2<t<t_0,
             |x-x_0-c(t-t_0)|<r\}
\]
and write
\begin{equation}\label{eq:galilean-reduced-C}
 C_W^\circ(z,r)
 :=\inf_{c\in\R^3}r^{-2}\int_{Q_r^c(z)}|W-c|^3\dd x\dd t.
\end{equation}
This is the velocity criterion applied after an exact constant Galilean
transformation; it is invariant under translations, Galilean changes, and
parabolic scaling.

\begin{proposition}[First bad scale gives an ordinary descendant]
\label{prop:first-bad-scale-descendant}
Fix \(0<\varepsilon_0<\varepsilon_{\mathrm{reg}}\), where
\(\varepsilon_{\mathrm{reg}}\) is the velocity regularity threshold used in
\cref{thm:velocity-concentration}.  Let \(U\) be a bounded mild ancient
solution and suppose that a diffuse remainder as in
\eqref{eq:diffuse-tail-mass} contains points at which \(U\ne0\), after the
homogeneous quotient has been removed.  Then one of the following occurs.
\begin{enumerate}[label=\textup{(\alph*)},leftmargin=2.2em]
\item A bounded-radius unit recentering carries a fixed positive quotient
activity, and alternative~\textup{(ii)} of
\cref{prop:backward-concentration-remainder} applies.
\item There are parabolic centers \(z_j\), radii \(\rho_j\to\infty\), and
constant Galilean parameters \(c_j\) such that the exact rescalings
\begin{equation}\label{eq:first-bad-rescaling}
 W_j(y,s):=\rho_j
   \bigl(U(x_j+\rho_jy+c_j\rho_j^2s,
                 t_j+\rho_j^2s)-c_j\bigr)
\end{equation}
with the correspondingly transformed pressure are suitable on expanding
backward cylinders, have a locally smooth subsequence modulo homogeneous
velocity and affine-pressure modes, and converge to an ordinary bounded mild
ancient solution \(W_\infty\).  On the nonhomogeneous branch,
\begin{equation}\label{eq:first-bad-descendant-floor}
 \int_{Q_1}|W_\infty-\bar W_{\infty,Q_1}|^3\ge c(\varepsilon_0)>0.
\end{equation}
Every compact subset of the limit is obtained by an exact parabolic
rescaling of the original solution; no measure-valued state is used.
\item On the selected first-bad cylinders the local velocity bounds hold but
the mean-subtracted pressures are not precompact in
\(L^{3/2}_{\mathrm{loc}}\).  In the local decomposition
\(P=P_{\mathrm{loc}}+H\), the Calder\'on--Zygmund part is compact, so this
failure is carried by a non-affine harmonic pressure tail sourced outside the
selected cylinder.
\end{enumerate}
If the rescaled forward time depth tends to infinity, the limit is entire.
If it stays bounded, the standard terminal self-similar change of variables
places the limit in \eqref{eq:renormalized-NS} with a constant coefficient
\(a_*<0\).  Thus both time-depth alternatives have nonpositive scale drift.
\end{proposition}

\begin{proof}
For a smooth point \(z\), one has
\(C_U^\circ(z,r)\to0\) as \(r\downarrow0\).  If
\(C_U^\circ(z,r)<\varepsilon_0\) on arbitrarily large backward cylinders,
apply quantitative velocity \(\varepsilon\)-regularity after an almost
minimizing Galilean transformation.  Its derivative estimate gives
\[
 |\nabla U(z)|\le \frac{C(\varepsilon_0)}{r^2}.
\]
Thus a solution for which this happens at every point is spatially
homogeneous.  On the nonhomogeneous quotient branch there is consequently a
point and a scale at which \(C_U^\circ=\varepsilon_0\).  Select the first
such scale.  If these first scales stay bounded, a bounded-radius cylinder
has a positive oscillatory subsequence and gives \textup{(a)}.

Assume that the first scales tend to infinity.  Apply the standard parabolic
point-selection procedure.  If a chosen cylinder contains a smaller cylinder
at which the threshold is first reached and whose distance to the boundary,
in its own units, is larger, replace the chosen cylinder by that one.  Each
replacement decreases the radius by a fixed factor while increasing the
available normalized boundary distance.  An infinite replacement chain in a
fixed compact physical cylinder would have radii tending to zero and centers
converging to a smooth point, contradicting
\(C_U(z,r)\to0\).  Thus, for the \(j\)-th selected cylinder, the procedure
terminates with a radius \(\rho_j\) such that every subcylinder of radius at
most \(\rho_j\) lying within normalized parabolic distance \(j\) has
Galilean-reduced velocity quantity at most \(\varepsilon_0\).  Absence of alternative
\textup{(a)} forces \(\rho_j\to\infty\).

Rescale by \eqref{eq:first-bad-rescaling}.  Cover a fixed compact cylinder in
the new variables by finitely many subunit cylinders.  On each of them the
velocity criterion is below \(\varepsilon_{\mathrm{reg}}\), so interior
regularity gives bounds for every velocity derivative on a slightly smaller
cover.  The local pressure decomposition
\[
 P=P_{\mathrm{loc}}+H,
 \qquad -\Delta P_{\mathrm{loc}}
 =\partial_i\partial_j(\chi W_iW_j),
 \qquad \Delta H=0,
\]
makes the local part compact after subtracting spatial means.  If the harmonic
parts are not compact modulo affine functions, alternative~\textup{(c)}
holds.  Otherwise the affine part is exactly the pressure paired with the
removed homogeneous velocity, and diagonal compactness produces an ordinary
suitable mild limit.  The threshold equality at the selected cylinder passes
by strong local \(L^3\) convergence.  If its value were carried only by the
homogeneous mode, the limit belongs to the quotient alternative.  Otherwise
Poincar\'e's inequality and the first-scale normalization give
\eqref{eq:first-bad-descendant-floor} after decreasing the constant once.

Let \(d_j^+\) be the forward time depth of the rescaled physical domain.  If
\(d_j^+\to\infty\), the diagonal limit is entire.  If \(d_j^+\) is bounded,
translate the terminal time to zero and use
\(t=-e^{-\tau}\), \(\Lambda=e^{-\tau/2}\).  The resulting equation is
\eqref{eq:typeI-centered-equation}, up to a constant rescaling of \(\tau\),
and hence has \(a_*<0\).  These are the only two subsequential values of the
time-depth ratio.
\end{proof}

\begin{proposition}[Finite energy makes the non-affine pressure tail
parabolically accessible]
\label{prop:harmonic-pressure-tail-accessibility}
Let \(u\) be a finite-energy solution on \((0,T)\), and let
\begin{equation}\label{eq:finite-energy-exact-rescalings}
 \widetilde W_j(y,s)
 :=\lambda_j u(x_j+\lambda_jy,t_j+\lambda_j^2s),
 \qquad
 d_j:=\frac{t_j}{\lambda_j^2}\longrightarrow\infty,
\end{equation}
where \(t_j\ge t_*>0\) and \(\lambda_j\downarrow0\).  A constant
Galilean velocity may be subtracted when local oscillation is measured; the
finite-energy representative \(\widetilde W_j\) is retained for the
whole-space pressure formula.  Put \(R_j=d_j^{1/4}\), and let
\[
 H_j^{>R_j}
 :=\mathcal R_i\mathcal R_k
   \bigl(\mathbf 1_{\{|y|>R_j\}}
         \widetilde W_{j,i}\widetilde W_{j,k}\bigr).
\]
For every fixed
\(R<\infty\) and compact interval \(I\Subset(-\infty,0]\), the part of the
pressure generated by \(|y|>R_j\), after subtraction of its constant and
linear Taylor terms on \(B_R\), satisfies
\begin{equation}\label{eq:far-pressure-tail-decay}
 \|\nabla^2 H_j^{>R_j}\|_{L^\infty(B_R\times I)}
 \le C_{R,I,t_*,\|u_0\|_2}\,d_j^{-3/4}\longrightarrow0
\end{equation}
at every time at which the energy inequality holds.  Moreover
\(R_j^2=d_j^{1/2}=o(d_j)\).  Consequently a failure of local pressure
compactness modulo affine functions cannot be produced at spatial distances
beyond the parabolic time depth.  It must be sourced in the intermediate
region
\begin{equation}\label{eq:accessible-pressure-annulus}
 2R<|y|<R_j,
 \qquad -R_j^2<s\le0,
\end{equation}
up to a term converging to zero locally and an affine pressure mode.
\end{proposition}

\begin{proof}
The global energy inequality and \eqref{eq:finite-energy-exact-rescalings}
give, uniformly at admissible rescaled times,
\begin{equation}\label{eq:rescaled-energy-depth-bound}
 \|\widetilde W_j(s)\|_2^2
 \le\lambda_j^{-1}\|u_0\|_2^2
 =\sqrt{\frac{d_j}{t_j}}\,\|u_0\|_2^2
 \le t_*^{-1/2}\|u_0\|_2^2d_j^{1/2}.
\end{equation}
The pressure Poisson equation is unchanged by subtraction of a spatially
constant velocity: the constant--constant term has zero second derivative,
and the two cross terms vanish because the velocity is divergence free.
Thus the tail may be estimated using the finite-energy representative.

The kernel of \(\mathcal R_i\mathcal R_k\) is homogeneous of degree
\(-3\), and its second spatial derivatives are bounded by
\(C|y|^{-5}\) away from the origin.  Taylor's formula on \(B_R\) therefore
gives
\begin{align*}
 \sup_{x\in B_R}|\nabla^2H_j^{>R_j}(x,s)|
 &\le C_R\int_{|y|>R_j}|y|^{-5}
          |\widetilde W_j(y,s)|^2\dd y\\
 &\le C_RR_j^{-5}\|\widetilde W_j(s)\|_2^2
 \le C_{R,t_*,\|u_0\|_2}d_j^{-3/4},
\end{align*}
because \(R_j^5=d_j^{5/4}\).  The bound is uniform on a fixed compact
rescaled time interval by the energy inequality.  Interior estimates for
harmonic functions control all lower local norms after the affine Taylor jet
has been removed.  Finally, \(R_j^2=d_j^{1/2}<d_j\) for large \(j\), so
the backward parabolic cylinders associated with the intermediate region lie
inside the rescaled time domain.  This proves the assertion.
\end{proof}

\begin{remark}[What the pressure estimate does not prove]
The accessibility estimate rules out a genuinely super-parabolic pressure
source.  It does not by itself turn the intermediate stress in
\eqref{eq:accessible-pressure-annulus} into a fixed positive
Galilean-reduced \(L^3\) oscillation on one cylinder.  The ordered argument
does not require that conversion: it retains the intermediate source as a
pressure-tail coordinate.  If a later first-bad-scale selection produces
velocity oscillation, that ordinary descendant is then recorded; otherwise
the pressure coordinate remains part of the marked critical tail treated by
\cref{thm:critical-tail-rigidity}.
\end{remark}

\begin{remark}[What the stopping time proves]
\label{rem:first-bad-scale-scope}
\Cref{prop:first-bad-scale-descendant} removes diffuse velocity mass as a
terminal label: it becomes an actual velocity descendant or a homogeneous
quotient state.  It also isolates the only pressure obstruction as a
non-affine harmonic tail with exterior ancestry; it does not silently call
that tail compact.
An infinite sequence of such actual descendants need not have a stationary
law satisfying all invariances in
\cref{def:stationary-monotone-ensemble}.  The marked construction in
\cref{app:marked-regeneration} retains the successor witnesses and their
currents and converts the resulting recurrent path into a critical tail.
\end{remark}

\begin{definition}[Actual target-retaining descendant]
\label{def:actual-target-descendant}
An actual target-retaining descendant of a bounded ancient represented
solution is a locally smooth limit, with compatible local pressure gauges, of
exact spacetime translations, constant Galilean changes, and constant
parabolic rescalings of that solution, for which a fixed quotient-oscillation
floor holds on \(Q_1\).  A first-bad-scale limit is included only through
\cref{prop:first-bad-scale-descendant}.  Thus every descendant is an ordinary
Navier--Stokes solution; a diffuse measure, a failed pressure chart, or a
formal boundary value is not a descendant.
\end{definition}

\begin{corollary}[Finite ordinary descendant chains close]
\label{cor:finite-descendant-chains-close}
Let \(U_0\) be the record-point profile in
\cref{thm:record-point-ancient-extraction}, after removal of the homogeneous
quotient.  Iterate \cref{prop:backward-concentration-remainder} and, on its
diffuse branch, the ordinary compact alternative of
\cref{prop:first-bad-scale-descendant}.  If the resulting target-retaining
chain terminates after finitely many actual descendants, then its last profile
is impossible or its first non-affine intermediate pressure source is a
typed pressure-tail output of \textup{(F3)}.  Thus every finite chain is
closed within the ordered source alternative; outside a typed first failure,
only an infinite chain of actual target-retaining descendants can remain.
\end{corollary}

\begin{proof}
At a terminal profile there is no further nonhomogeneous descendant and no
diffuse remainder.  If pressure compactness modulo affine functions fails,
\cref{prop:harmonic-pressure-tail-accessibility} retains the first non-affine
source with its exterior ancestry, which is the asserted \textup{(F3)}
output.  Otherwise the homogeneous alternative has zero quotient target and
the last node satisfies the backward sequence-\(L^3\) alternative of
\cref{prop:backward-concentration-remainder}.  It vanishes by
\cref{thm:AB-endpoint}, contradicting its retained quotient-oscillation floor.
The same conclusion holds after a terminal self-similar change because its
physical pullback is a bounded mild ancient solution and the critical
\(L^3\)-norm is invariant.
\end{proof}

\begin{lemma}[Event law for an infinite realized chain]
\label{lem:compact-descendant-path-law}
Let \(K_\eta\) be a compact metric space of ordinary normalized represented
solutions on which the quotient target \(\mathcal O\) is continuous and
\(\mathcal O\ge\eta>0\).  Let
\(\mathcal R\subset K_\eta\times K_\eta\) be an actual-descendant
relation.  If
\begin{equation}\label{eq:actual-infinite-descendant-chain}
 w_0\mathcal Rw_1\mathcal Rw_2\mathcal R\cdots,
\end{equation}
then the event-count measures
\begin{equation}\label{eq:actual-descendant-event-law}
 \nu_N:=\frac1N\sum_{j=0}^{N-1}\delta_{w_j}
\end{equation}
have weak-* limit points, every one of which is supported on
\(\{\mathcal O\ge\eta\}\).  The closure of the shift orbit of the given
one-sided sequence carries a shift-invariant probability measure whose zeroth
marginal has the same property.  No closedness of the profile-only relation is
asserted.
\end{lemma}

\begin{proof}
Compactness of \(K_\eta\) gives weak-* compactness of
\(\mathcal P(K_\eta)\).  Every measure in
\eqref{eq:actual-descendant-event-law} is supported on the closed set
\(\{\mathcal O\ge\eta\}\), and the Portmanteau theorem preserves this support
condition under weak-* limits.  The closure of the shift orbit of
\((w_j)_{j\ge0}\) is compact in \(K_\eta^{\mathbb N}\).  Weak-* limits of the
shift empirical measures are shift invariant because, for every continuous
\(F\),
\[
 \frac1N\sum_{j=0}^{N-1}
 \bigl(F(S^{j+1}w)-F(S^jw)\bigr)
 =\frac{F(S^Nw)-F(w)}N\longrightarrow0.
\]
Their zeroth marginals are limits of
\eqref{eq:actual-descendant-event-law}.  The successor edge can be lost when
its recentering witness escapes, so the conclusion deliberately concerns the
event law and not a closed path relation.
\end{proof}

\begin{remark}[The profile-only path space]
If the unmarked relation happens to be closed, then the path space
\begin{equation}\label{eq:actual-descendant-path-space}
 \mathscr P_\eta
 :=\{(v_j)_{j\ge0}\in K_\eta^{\mathbb N}:
           v_j\mathcal Rv_{j+1}\text{ for all }j\}
\end{equation}
is nonempty and compact, is invariant under the left shift \(S\), and carries
an \(S\)-invariant Borel probability measure.  Every coordinate marginal of
that measure is target-retaining.  The hypothesis is not available for the
profile-only descendant relation: local convergence on fixed cylinders does
not control observers whose centers escape.  The marked construction below
restores closedness by retaining those observers.
\end{remark}

\begin{proposition}[A rooted target is not preserved by an invariant hull]
\label{prop:rooted-target-not-hull-invariant}
Let \(Q\) be a fixed bounded cylinder and define
\[
 \mathcal O_Q(V):=\inf_{c\in\R^3}\int_Q|V-c|^3.
\]
Suppose that spatial translates \(T_{x_j}V\) converge strongly in
\(L^3(Q)\) to a spatially constant field as \(|x_j|\to\infty\).  Then
\begin{equation}\label{eq:rooted-target-lost-at-infinity}
 \mathcal O_Q(T_{x_j}V)\longrightarrow0.
\end{equation}
Consequently the closed covariant hull of a locally concentrated profile need
not remain in \(\{\mathcal O_Q\ge\eta\}\), even when the profile itself has
\(\mathcal O_Q(V)\ge\eta\).  In particular, invariance of a probability law
under the successor shift does not imply the spatial invariance required in
\cref{def:stationary-monotone-ensemble} while retaining the root target.
\end{proposition}

\begin{proof}
The map \(V\mapsto\mathcal O_Q(V)^{1/3}\) is \(1\)-Lipschitz in
\(L^3(Q)\): apply the triangle inequality before taking the infimum over
constants, and then interchange the two fields.  Its value on a spatially
constant field is zero, which proves
\eqref{eq:rooted-target-lost-at-infinity}.  The remaining assertions follow
because a covariantly invariant hull contains all finite spatial translates
and their local limit points, whereas the successor shift acts on a sequence
of separately recentered profiles.
\end{proof}

\subsection{Marked descendant structure for regenerated current}
\label{app:marked-regeneration}

The successor shift and the covariant translation action are different
dynamical systems.  The preceding proposition prevents replacing one by the
other.  We instead retain the geometric and local-energy data which realize
each successor.  The resulting relation is closed because an escaping witness
is recorded as an exterior, pressure, diffuse, or scale-tail output rather than
discarded.

Fix once and for all a countable family \(\mathscr T\) of nonnegative smooth
parabolic cutoffs which is dense in the test-function Fr\'echet topology on
each member of a compact-cylinder exhaustion.  The stopping radii and times
are chosen from the full-measure set for which the endpoint energy traces and
the coarea limits of every test in \(\mathscr T\) exist.  Equivalently, the
boundary currents below may be defined as limits of smooth collars.  The local
energy distribution, including its nonnegative Radon defect, then defines for
every physical parabolic cell \(C\) the coordinates
\begin{equation}\label{eq:marked-local-energy-coordinates}
 \mathbf q_C(\psi):=
 \bigl(K_C^-(\psi),K_C^+(\psi),D_C(\psi)+M_C(\psi),J_C(\psi)\bigr),
 \qquad \psi\in\mathscr T,
\end{equation}
where the endpoint energy traces and the complete signed current are oriented with
respect to \(C\).  The suitable-energy defect is included in the third
coordinate, so the balance is an equality.

\begin{definition}[Marked successor record]
\label{def:marked-successor-record}
Let \(C_j=Q_{r_j}(z_j)\) and
\(C_{j+1}=Q_{r_{j+1}}(z_{j+1})\Subset Q_{r_j/2}(z_j)\) be two actual cells cut from the
same suitable solution.  A marked successor record is
\begin{equation}\label{eq:marked-successor-tuple}
 \mathfrak m_j=
 (\mathfrak w_j,\mathfrak w_{j+1},C_j,
   \lambda_j,\xi_j,\theta_j,\gamma_j,
   \mathbf q_j,\mathfrak s_j),
\end{equation}
where
\begin{equation}\label{eq:marked-observer-increments}
 \lambda_j=\frac{r_{j+1}}{r_j},\qquad
 \xi_j=\frac{x_{j+1}-x_j}{r_j},\qquad
 \theta_j=\frac{t_{j+1}-t_j}{r_j^2}.
\end{equation}
Here \(\gamma_j\) is the additive pressure gauge on the overlap,
\(\mathbf q_j\) is the countable vector
\((\mathbf q_{C_j}(\psi))_{\psi\in\mathscr T}\), and
\(\mathfrak s_j\) is the inverse-limit support coordinate: for every finite
block ending at \(j\), every first finite family of cylinders and tests, and
every \(m\ge1\), it records a terminal-sequence index at least \(m\), a
composed observer, and a realization error at most \(m^{-1}\) in the strong
velocity and pressure coordinates.  Compatibility is restriction of the
\((m+1)\)-st record to the first \(m\) coordinates.  Concatenation
requires equality of the outgoing window, incoming window, pressure gauge,
and oriented current trace on the common overlap.
\end{definition}

The topology of the marked records is strong local \(L^3\) for velocity,
strong local \(L^{3/2}\) modulo the recorded gauge for pressure, weak local
\(L^2\) for gradients, weak-* convergence for the suitable defect, Euclidean
convergence for \((\lambda,\xi,\theta)\), and coordinatewise convergence of
\(\mathbf q\) and \(\mathfrak s\).  The countable coordinates make this
topology metrizable on bounded sets.

\begin{definition}[Uniformly record-bounded terminal extraction]
\label{def:record-bounded-terminal-extraction}
A terminal extraction \((U_n,P_n)\) on domains \(\mathcal D_n\) is
\emph{uniformly record-bounded} if each pair is an exact parabolic rescaling
of the same physical solution, the domains exhaust every compact cylinder
used by the extraction, and
\begin{equation}\label{eq:uniform-record-bound}
 \sup_n\|U_n\|_{L^\infty(\mathcal D_n)}\le M_*
\end{equation}
for one \(M_*<\infty\).  The record-point sequence in
\cref{thm:record-point-ancient-extraction} has this property with \(M_*=1\)
by \eqref{eq:record-profile-bound}.
\end{definition}

\begin{definition}[Marked regenerated-current cascade]
\label{def:marked-regenerated-cascade}
A marked regenerated-current cascade at level \(\eta>0\) is an infinite
concatenated sequence \((\mathfrak m_j)_{j\ge0}\) with the following
properties.  Write \(V_j\) for the velocity component of \(\mathfrak w_j\).
\begin{enumerate}[label=\textup{(R\arabic*)},leftmargin=2.8em]
\item The cells are strictly nested and
\begin{equation}\label{eq:marked-cascade-scale-separation}
 0<\lambda_j\le\frac14,
 \qquad
 \mathcal O_{Q_1}(V_j)\ge\eta .
\end{equation}
Every finite block is simultaneously realized in one member of the original
terminal sequence by the recorded composed observers.
\item With \(w_j=r_j^{-1}\), let \(\mathscr J_j\) be the complete oriented
local-energy current through the stopping surface between \(C_j\) and
\(C_{j+1}\), including its endpoint energy trace, and let
\(\mathscr P_j\ge0\) be viscous production plus suitable defect in that shell.
The scale-critical shell dissipation is
\begin{equation}\label{eq:marked-critical-dissipation}
 \mathscr D_j^{\rm crit}:=w_j\mathscr P_j .
\end{equation}
It is invariant under Navier--Stokes scaling because physical local-energy
production has homogeneity one in length.
There are \(N_k\uparrow\infty\) and \(\delta_D>0\) such that
\begin{equation}\label{eq:marked-positive-cesaro-dissipation}
 \frac1{N_k}\sum_{j=0}^{N_k-1}\mathscr D_j^{\rm crit}\ge\delta_D
 \qquad(k\ge1).
\end{equation}
\item Every nonzero sequence-realized support profile of the remaining
weighted current regenerates a child satisfying
\eqref{eq:marked-cascade-scale-separation}.
\item The marked records have compact closure; in particular there are
\(\lambda_*>0\) and \(B<\infty\)
such that
\begin{equation}\label{eq:compact-mark-bounds}
 \lambda_*\le\lambda_j\le\frac14,
 \qquad |\xi_j|+|\theta_j|\le B.
\end{equation}
\item The original terminal extraction is uniformly record-bounded in the
sense of \cref{def:record-bounded-terminal-extraction}.  Every finite initial
block of length \(N\) has a simultaneous realization whose ungauged root velocity
\(U^{(N)}\) satisfies
\begin{equation}\label{eq:bounded-root-finite-realization}
 \|U^{(N)}\|_{L^\infty(\mathcal Q_N)}\le M_*.
\end{equation}
Here \(\mathcal Q_N\) contains the images, in root coordinates, of all unit
cylinders in the block.  The normalized velocity at every descendant is the
exact parabolic rebase of \(U^{(N)}\), up to an additive constant Galilean
velocity.  The realization accuracy on the terminal descendant may be chosen
so that its velocity \(V_N^{(N)}\) satisfies
\begin{equation}\label{eq:bounded-terminal-realization-accuracy}
 \|V_N^{(N)}-V_N\|_{L^3(Q_1)}\le\varepsilon_N,
 \qquad \varepsilon_N\downarrow0.
\end{equation}
\end{enumerate}
\end{definition}

\begin{proposition}[Extraction of the marked alternative]
\label{prop:marked-alternative-extraction}
Consider an infinite nested sequence of target-retaining current events
furnished by one finite-stage concentration--remainder construction from a
uniformly record-bounded terminal extraction.  Suppose that each event is
accompanied by its actual successor witness and the coordinates in
\eqref{eq:marked-successor-tuple}.  Apply the following tests in order:
\begin{enumerate}[label=\textup{(E\arabic*)},leftmargin=2.8em]
\item if the Ces\`aro mean of the scale-critical shell dissipations has a
subsequence tending to zero, select shells whose normalized dissipation tends
to zero;
\item if a nonzero sequence-realized support profile of the remaining current
has no target-retaining child, select that profile at its first failed
compactness, pressure, exterior, diffuse, affine, or critical-tail test;
\item if the radii fail to tend to zero, select the compact
non-collapsing scale limit; otherwise retain a subsequence with
\(r_{j+1}/r_j\le1/4\);
\item if, on this subsequence, the ratios have no positive lower bound, the
normalized observer increments are unbounded, or the velocity, pressure
gauge, suitable-defect, current, or realization coordinates have no compact
closure, select the corresponding scale-tail, exterior, pressure/noncompact,
defect, or realization-loss profile.
\end{enumerate}
If none of these selected outputs occurs, the remaining events contain a
marked regenerated-current cascade in the sense of
\cref{def:marked-regenerated-cascade}.  Thus failure of the compact marked
cascade is an explicit first-failure output; it is not an additional hypothesis hidden in
the conclusion.
\end{proposition}

\begin{proof}
Put \(A_N=N^{-1}\sum_{j<N}\mathscr D_j^{\rm crit}\).  If \(\liminf_NA_N=0\), choose
\(N_k\to\infty\) with \(A_{N_k}\to0\).  Since these dissipations are nonnegative,
for every \(k\) some \(j_k<N_k\) satisfies
\(\mathscr D_{j_k}^{\rm crit}\le A_{N_k}\); this is test \textup{(E1)}.
Otherwise choose \(0<\delta_D<\liminf_NA_N\); then
\eqref{eq:marked-positive-cesaro-dissipation} holds on every sufficiently long
initial block.

Starting from the first retained event, follow a target-retaining child of
each nonzero current-support profile.  If no such child exists, the first
failed test is exactly \textup{(E2)}.  If a child always exists, recursive
selection gives an infinite realized chain.  If its radii do not tend
to zero, a subsequence has a positive limiting radius; after normalization
this is the compact non-collapsing scale output in \textup{(E3)}.  If the
radii tend to zero, choose successively the first later radius not exceeding
one quarter of the preceding selected radius.  Relabeling this subsequence
gives the upper bound in \eqref{eq:marked-cascade-scale-separation}; finite
simultaneous realization is preserved under passage to a subsequence and
composition of the intervening observers.

On the separated chain, failure of any lower scale bound, observer bound, or
coordinate compactness has a first coordinate in the countable cylinder and
test exhaustion on which compactness fails.  Selecting that coordinate gives
the output in \textup{(E4)}.  Diagonal compactness gives the compact closure
and the bounds in
\eqref{eq:compact-mark-bounds}.  The retained activity floor and the
regeneration property are closed because velocity convergence is strong in
the fixed local \(L^3\) topology and because the complete successor witness is
one of the marks.  The record bound and the inverse-limit realization
coordinate give \textup{(R5)} by
\cref{lem:record-bound-inherited-by-marked-blocks} below.  Thus properties
\textup{(R1)}--\textup{(R5)} hold, so the
remaining chain is a marked regenerated-current cascade.
\end{proof}

\begin{lemma}[Exact scale-critical shell balance]
\label{lem:marked-critical-shell-balance}
For every finite initial block of a marked regenerated-current cascade,
\begin{equation}\label{eq:marked-critical-shell-balance}
 \sum_{j=0}^{N-1}\mathscr D_j^{\rm crit}
 =w_0\mathscr J_0-w_{N-1}\mathscr J_N
  +\sum_{j=1}^{N-1}(w_j-w_{j-1})\mathscr J_j .
\end{equation}
Common parent--child currents cancel with opposite orientations.  Viscous
production and suitable defect remain nonnegative on the disjoint shell
interiors.
\end{lemma}

\begin{proof}
Apply the suitable local energy identity, with its nonnegative defect included,
to the disjoint regions between consecutive stopping surfaces.  The normal
trace on a common surface occurs once with each orientation and cancels.
Writing the resulting shell identity as
\(\mathscr P_j=\mathscr J_j-\mathscr J_{j+1}\), multiplying by \(w_j\), and
summing gives
\[
 \sum_{j<N}w_j(\mathscr J_j-\mathscr J_{j+1})
 =w_0\mathscr J_0-w_{N-1}\mathscr J_N
  +\sum_{j=1}^{N-1}(w_j-w_{j-1})\mathscr J_j.
\]
This is \eqref{eq:marked-critical-shell-balance}.  The sign assertion follows
because the dissipation and defect measures are restricted to disjoint
interiors; they are not obtained by subtracting a parent localization from its
children.
\end{proof}

\begin{lemma}[Finite simultaneous realization]
\label{lem:marked-finite-simultaneous-realization}
Every finite set of marked events and every finite set of compact normalized
cylinders on which their activity and current coordinates are tested are
realized, with compatible pressure gauges, by one member of the original
terminal extraction.
\end{lemma}

\begin{proof}
This is precisely the finite-block clause in \textup{(R1)}.  More explicitly,
the inverse-limit coordinate \(\mathfrak s_j\) records, for each block length
and each finite cylinder-and-test family, one terminal-sequence index and the
composed observer realizing all of those coordinates.  Restriction
compatibility of \(\mathfrak s_j\) preserves the same index when coordinates
are discarded, and the overlap coordinate \(\gamma_j\) makes the pressure
representatives agree up to their recorded additive gauge.  Hence all members
of the prescribed finite family occur in one realization.  A diagonal choice
over the countable exhaustion gives the final assertion.
\end{proof}

\begin{lemma}[The record bound is inherited by finite marked blocks]
\label{lem:record-bound-inherited-by-marked-blocks}
Let the marked events be generated from a uniformly record-bounded terminal
extraction.  If every finite block has the inverse-limit realization
coordinate of \cref{def:marked-successor-record}, then clause \textup{(R5)}
of \cref{def:marked-regenerated-cascade} holds with the same constant \(M_*\)
as in \eqref{eq:uniform-record-bound}.
\end{lemma}

\begin{proof}
Fix \(N\).  The inverse-limit coordinate supplies terminal-sequence indices
realizing the first \(N\) edges on the finite collection of cylinders needed
for that block, with error tending to zero in every recorded strong local
velocity coordinate.  Choose one such index far enough along the sequence
that the terminal velocity error on \(Q_1\) is at most \(N^{-1}\), and call
the corresponding extraction member \(U^{(N)}\).  This is the ungauged root
field: Galilean constants are retained in the composed observer and are not
added to \(U^{(N)}\).  Since every image cylinder in the finite block lies in
the domain of this realization,
\(\mathcal Q_N\subset\mathcal D_{k(N)}\), and hence
\[
 \|U^{(N)}\|_{L^\infty(\mathcal Q_N)}
 \le \|U_{k(N)}\|_{L^\infty(\mathcal D_{k(N)})}
 \le M_* .
\]
Exact composition of translations, constant Galilean changes, and parabolic
dilations gives the required descendant rebases.  Taking
\(\varepsilon_N=N^{-1}\) proves
\eqref{eq:bounded-terminal-realization-accuracy} and therefore \textup{(R5)}.
\end{proof}

\begin{lemma}[Quotient activity under a bounded parabolic rebase]
\label{lem:bounded-rebase-activity-decay}
Let \(Q\) be a bounded spacetime cylinder, let \(U\) be defined on the image
of an affine parabolic embedding \(\Psi:Q\to\mathbb R^3\times\mathbb R\), and
suppose \(\|U\|_{L^\infty(\Psi(Q))}\le M\).  For \(0<\lambda\le1\) and
\(h\in\mathbb R^3\), put
\begin{equation}\label{eq:bounded-parabolic-rebase}
 V(y,s):=\lambda U(\Psi(y,s))+h .
\end{equation}
Then
\begin{equation}\label{eq:bounded-rebase-activity-bound}
 \mathcal O_Q(V)^{1/3}
 \le |Q|^{1/3}M\lambda .
\end{equation}
Moreover, for any \(V_1,V_2\in L^3(Q)\),
\begin{equation}\label{eq:quotient-activity-lipschitz}
 \left|\mathcal O_Q(V_1)^{1/3}
       -\mathcal O_Q(V_2)^{1/3}\right|
 \le\|V_1-V_2\|_{L^3(Q)}.
\end{equation}
\end{lemma}

\begin{proof}
Choose the constant \(c=h\) in the infimum defining \(\mathcal O_Q(V)\).
Then
\[
 \mathcal O_Q(V)
 \le\int_Q|V-h|^3
 =\lambda^3\int_Q|U\circ\Psi|^3
 \le |Q|M^3\lambda^3,
\]
which proves \eqref{eq:bounded-rebase-activity-bound}.  For every constant
\(c\), Minkowski's inequality gives
\[
 \|V_1-c\|_{L^3(Q)}
 \le\|V_1-V_2\|_{L^3(Q)}+\|V_2-c\|_{L^3(Q)}.
\]
Take the infimum over \(c\), and then interchange \(V_1,V_2\), to obtain
\eqref{eq:quotient-activity-lipschitz}.
\end{proof}

\begin{lemma}[Composition of marked descendant rebases]
\label{lem:marked-rebase-composition}
For a finite concatenated block of marked successors, let
\begin{equation}\label{eq:marked-composed-scale}
 \Lambda_N:=\prod_{j=0}^{N-1}\lambda_j .
\end{equation}
In every simultaneous realization of that block, its terminal normalized
velocity has the form
\begin{equation}\label{eq:marked-composed-rebase}
 V_N^{(N)}(y,s)
 =\Lambda_N U^{(N)}(\Psi_N(y,s))+h_N
 \qquad ((y,s)\in Q_1),
\end{equation}
where \(U^{(N)}\) is the ungauged root representative, \(h_N\) is a constant
Galilean velocity, and \(\Psi_N(Q_1)\subset\mathcal Q_N\).  In particular,
if the block satisfies \textup{(R5)}, then
\begin{equation}\label{eq:terminal-block-activity-bound}
 \mathcal O_{Q_1}(V_N^{(N)})^{1/3}
 \le |Q_1|^{1/3}M_*\Lambda_N .
\end{equation}
\end{lemma}

\begin{proof}
For one edge, the normalized child velocity is obtained from the normalized
parent by an exact translation, parabolic dilation of ratio \(\lambda_j\),
and a constant Galilean change.  It therefore has the form
\[
 V_{j+1}(y,s)
 =\lambda_jV_j(\Psi_j(y,s))+g_j,
\]
where \(\Psi_j\) is the affine parabolic embedding determined by
\((\xi_j,\theta_j)\) and the Galilean spatial drift.  Composition multiplies
the amplitude factors and again gives an affine parabolic embedding; all
spatially constant velocity terms combine into one constant vector.  Induction
proves \eqref{eq:marked-composed-rebase}, with amplitude factor
\(\Lambda_N\).  The cells are nested, so the image of the terminal unit
cylinder lies in \(\mathcal Q_N\).  Apply
\cref{lem:bounded-rebase-activity-decay} with \(M=M_*\) to obtain
\eqref{eq:terminal-block-activity-bound}.
\end{proof}

\begin{theorem}[Record-bounded marked cascades are impossible]
\label{thm:no-bounded-marked-regenerated-cascade}
There is no marked regenerated-current cascade generated from a uniformly
record-bounded terminal extraction with a positive activity floor.
\end{theorem}

\begin{proof}
Suppose that a cascade exists at level \(\eta>0\).  For each \(N\), use
\textup{(R5)} to realize the first \(N\) edges with terminal velocity
\(V_N^{(N)}\) satisfying
\eqref{eq:bounded-terminal-realization-accuracy}.
By \textup{(R1)},
\(\mathcal O_{Q_1}(V_N)\ge\eta\).  On the other hand,
\cref{lem:bounded-rebase-activity-decay,lem:marked-rebase-composition} give
\begin{align*}
 \eta^{1/3}
 &\le \mathcal O_{Q_1}(V_N)^{1/3}\\
 &\le \mathcal O_{Q_1}(V_N^{(N)})^{1/3}+\varepsilon_N\\
 &\le |Q_1|^{1/3}M_*\Lambda_N+\varepsilon_N.
\end{align*}
Since \(\lambda_j\le1/4\),
\(\Lambda_N\le4^{-N}\to0\).  Letting \(N\to\infty\) gives
\(\eta^{1/3}\le0\), a contradiction.
\end{proof}

\begin{lemma}[Compact closed marked path space]
\label{lem:compact-closed-marked-path-space}
Let \((\mathfrak m_j)_{j\ge0}\) be a marked regenerated-current cascade.
The closure of its shift orbit is a compact shift-invariant space
\(X_{\mathrm{mark}}\).  Concatenation and the activity floor are closed there.
Every weak-* limit of the shift empirical measures is shift invariant and has
zeroth marginal supported on \(\{\mathcal O_{Q_1}\ge\eta\}\).
\end{lemma}

\begin{proof}
Property \textup{(R4)} gives relative compactness of the complete marked
records, not merely of their profile coordinates.  Take the closure in the
countable product of the compact coordinate closures.  The inverse-limit
support coordinate is closed by
\cref{lem:marked-finite-simultaneous-realization}, so every point of this
closure is still sequence-realized.

Concatenation is equality of outgoing and incoming overlap coordinates and is
therefore closed.  Strong local \(L^3\) convergence makes the activity floor
closed.  Hence the shift-orbit closure is a compact shift-invariant space of
concatenated marked paths.  The endpoint calculation in the proof of
\cref{lem:compact-descendant-path-law} gives shift invariance of every
empirical weak-* limit, and the Portmanteau theorem preserves the activity
floor.
\end{proof}

\begin{definition}[Sequence-realized logarithmic critical tail]
\label{def:sequence-realized-log-tail}
A sequence-realized logarithmic critical tail is a compatible projective
family of annular limits on
\[
 \{(y,s):e^{-L}<|y|<1,\ -1<s<0\},\qquad L=1,2,\ldots,
\]
obtained from one diagonal subsequence of the original terminal extraction.
The family includes velocity, pressure modulo its compatible additive gauge,
the suitable defect, and all current coordinates in \(\mathscr T\).  It is
nonzero when one fixed annular cross-section has quotient oscillation bounded
below by a positive constant.  Loss of tightness as \(L\to\infty\), survival
of a pressure or defect coordinate, and strong coherent convergence are,
respectively, its diffuse, defect, and coherent alternatives.
\end{definition}

\begin{theorem}[Marked recurrence produces a realized critical tail]
\label{thm:marked-recurrence-critical-tail}
Every marked regenerated-current cascade produces a nonzero
sequence-realized logarithmic critical-tail profile.  More precisely, there
are a recurrent path \(p\in X_{\mathrm{mark}}\) and integers
\(n_k\to\infty\) for which the same non-affine marked profile returns through
observers whose total scale tends to zero.  The associated annular blocks are
simultaneously realized by the original terminal sequence.
\end{theorem}

\begin{proof}
Let \(\varpi\) be a shift-invariant weak-* limit on
\(X_{\mathrm{mark}}\).  Poincar\'e recurrence gives a recurrent path in the
support of \(\varpi\), still satisfying the activity floor, and integers
\(n_k\to\infty\) such that
\begin{equation}\label{eq:marked-recurrence-return}
 S^{n_k}p\longrightarrow p\quad\hbox{in }X_{\mathrm{mark}}.
\end{equation}
This convergence includes velocity, pressure gauge, defect, current, observer,
and realization coordinates.

Composing the first \(n\) observer maps gives
\begin{align}
 \Lambda_n&=\prod_{j=0}^{n-1}\lambda_j,
 \label{eq:marked-total-scale}\\
 X_n-X_0&=r_0\sum_{j=0}^{n-1}
   \left(\prod_{i=0}^{j-1}\lambda_i\right)\xi_j,
 \label{eq:marked-total-center}\\
 T_n-T_0&=r_0^2\sum_{j=0}^{n-1}
   \left(\prod_{i=0}^{j-1}\lambda_i^2\right)\theta_j.
 \label{eq:marked-total-time}
\end{align}
The bounds \eqref{eq:compact-mark-bounds} make the last two series convergent,
whereas
\begin{equation}\label{eq:marked-scale-collapse}
 0<\Lambda_{n_k}\le4^{-n_k}\longrightarrow0.
\end{equation}
After recentering once at the limits in
\eqref{eq:marked-total-center}--\eqref{eq:marked-total-time},
\eqref{eq:marked-recurrence-return} is a return of the same marked profile at
logarithmic displacements tending to infinity.

Retain the complete blocks between scales \(1\) and \(\Lambda_{n_k}\).
Their component shells have disjoint logarithmic interiors.  By
\cref{lem:marked-finite-simultaneous-realization}, every finite family of
annular velocity, pressure, defect, and current tests is realized by one member
of the original terminal extraction.  A diagonal extraction therefore gives a
sequence-realized logarithmic annular state.  It is nonzero because the
recurrent cross-section retains \(\mathcal O_{Q_1}\ge\eta\).

Loss of logarithmic tightness defines the diffuse critical-tail component.
Under logarithmic tightness, failure of strong annular compactness or survival
of a pressure or suitable-defect coordinate defines the defect component.  In
the remaining case the annular convergence is strong, the pressure gauges are
compatible, and the limiting state is a coherent sequence-realized critical
tail.  These alternatives are exhaustive by weak-* compactness of the
logarithmic mass, local compactness on each fixed annulus, and the dichotomy
between zero and nonzero defect coordinates.  In every case the output is the
nonzero critical-tail profile asserted in the theorem.
\end{proof}

\subsection{Closure of the amplitude-renormalizing residual}
\label{app:collapsed-residual-closure}

The preceding recurrence theorem is useful only after the current which moves
the root has been put in the same scale-critical units as the retained
activity.  This subsection performs that accounting.  It is important that
the argument uses the successor shift on actual marked events.  It never
replaces that shift by the full spatial translation hull, and hence never
loses a sparse active root by ordinary spatial averaging.

\begin{lemma}[A retained window has a fixed critical viscous price]
\label{lem:target-window-critical-price}
Let \(\mathfrak w=(V,P,a,b;I)\) be an
\((\eta_0,\mathbf M)\)-normalized target-retaining window.  There is
\(d_*=d_*(\eta_0,M_{R_*},R_*,\zeta)>0\) such that
\begin{equation}\label{eq:target-window-critical-price}
 \int_I\int_{B_{R_*}}|\nabla V|^2\,dy\,d\tau\ge d_*.
\end{equation}
The number \(d_*\) is dimensionless under Navier--Stokes scaling.
\end{lemma}

\begin{proof}
Put \(f=\zeta(V-\bar V_\zeta)\).  The support of \(\zeta\) is connected,
and the weighted mean has been removed.  The local Poincar\'e--Sobolev
inequality, followed by interpolation between \(L^2\) and \(L^6\), gives for
almost every \(\tau\in I\)
\[
 \|f(\tau)\|_3^3
 \le C_\zeta\|V(\tau)\|_{L^2(B_{R_*})}^{3/2}
       \|\nabla V(\tau)\|_{L^2(B_{R_*})}^{3/2}.
\]
The harmless terms containing \(\nabla\zeta\) are absorbed by applying the
weighted Poincar\'e inequality to \(V-\bar V_\zeta\) before multiplication by
\(\zeta\).  Integrating in time and using H\"older's inequality and
\eqref{eq:target-window-bounds} yields
\[
 \eta_0
 \le \OscMass_\Gamma(V;I)
 \le C_\zeta M_{R_*}^{3/2}|I|^{1/4}
       \left(\int_I\int_{B_{R_*}}|\nabla V|^2\right)^{3/4}.
\]
Since \(|I|=1\), rearrangement proves \eqref{eq:target-window-critical-price},
for example with
\(d_*=(\eta_0/(C_\zeta M_{R_*}^{3/2}))^{4/3}\).
Every quantity in the displayed inequality is written in normalized
variables, so the lower bound is scale critical.
\end{proof}

\begin{definition}[Sparse marked current path]
\label{def:sparse-marked-current-path}
A sparse marked current path is an infinite concatenated path of the records
in \cref{def:marked-successor-record} which satisfies clauses
\textup{(R1)}--\textup{(R4)} of
\cref{def:marked-regenerated-cascade}, the uniform normalized window bounds
\eqref{eq:target-window-bounds}, and the following ordered conditions.
\begin{enumerate}[label=\textup{(S\arabic*)},leftmargin=2.8em]
\item Every ordinary occupation average over the normalized spatial
translation subgroup assigns zero mass to the set
\(\{\mathcal O_{Q_1}\ge\eta\}\).
\item Event-count averages along the successor path are not covariantly
balanced.  Equivalently, a fixed element of the countable test family detects
a nonzero translation, scale, pressure, or moving-boundary current along a
subsequence.
\item The first loss of velocity compactness, pressure-gauge compactness,
observer tightness, defect tightness, or simultaneous realization, and every
nonzero residual pressure or suitable-energy defect coordinate, is retained
as a marked output and is not identified with an ordinary velocity profile.
\end{enumerate}
Thus sparsity describes loss of the root under volume averaging, whereas the
marked path continues to see every active event.
\end{definition}

\begin{lemma}[Positive density and vanishing event current are terminal]
\label{lem:density-current-sparse-reduction}
An infinite target-retaining successor path either is impossible by covariant
rigidity or has a subsequence which is a sparse marked current path.
\end{lemma}

\begin{proof}
Let \(\nu_N=N^{-1}\sum_{j<N}\delta_{\mathfrak w_j}\) be the event-count
measures.  Compactness of the complete marked records gives a weak-* limit
after passage to a subsequence, and the endpoint calculation makes its path
law invariant under the successor shift.  Strong local \(L^3\) convergence
preserves the activity floor.

If active roots have positive upper density in a F\o lner sequence of the
normalized spatial translation group, average the event-count law and its
spatial translates.  The endpoint calculation gives stationarity under the
successor shift, while the F\o lner calculation gives spatial covariance.  A
weak-* limit assigns positive mass to the closed active set.  The pressure
cancellation in \cref{lem:recurrent-pressure-cancellation}, followed by
\cref{thm:stationary-monotone-ensemble-rigidity}, makes its velocity spatially
homogeneous almost surely, contradicting that positive mass.

If the active roots have zero covariant density but a subsequence of
\(\nu_N\) is covariantly balanced on the countable dense test algebra, its
weak-* limit is again stationary under the successor shift and invariant under
normalized spatial translations: first pass the balance to the limit on the
dense algebra and then use continuity of the action.  The same rigidity
contradiction applies.  In the remaining case ordinary averages
lose the target and some test detects a persistent covariance defect.  The
infinitesimal pressure component of that defect is exactly the rooted
correction in \cref{lem:rooted-pressure-current-identity}; all other components are
the translation, dilation, cutoff, and exterior coordinates already present
in \(\mathbf q_j\).  Retaining the first noncompact coordinate gives
\textup{(S3)}.  Hence the remaining path satisfies
\textup{(S1)}--\textup{(S3)}.
\end{proof}

\begin{lemma}[Nesting--antichain alternative for realized active events]
\label{lem:active-event-nesting-antichain}
Partially order simultaneously realized active parabolic cylinders by
strict inclusion after shrinking each cylinder by a fixed factor.  For an
arbitrarily long finite family, exactly one of the following occurs along a
subsequence:
\begin{enumerate}[label=\textup{(\roman*)},leftmargin=2.2em]
\item the lengths of inclusion chains tend to infinity;
\item the family contains antichains of arbitrarily large cardinality.
\end{enumerate}
Every fixed finite chain or antichain is realized in one member of the
original terminal sequence, with compatible pressure gauges and observer
maps.
\end{lemma}

\begin{proof}
If the heights of the finite partially ordered sets are unbounded, choose a
chain whose length tends to infinity and diagonalize.  Otherwise their heights
are bounded by some integer \(L\).  Assign to each cylinder the maximum length
of an inclusion chain ending there.  This partitions a family of \(m\)
cylinders into at most \(L\) antichains, one of which has cardinality at least
\(m/L\).  Hence the second alternative occurs.

For simultaneous realization, proceed inductively through the chosen finite
family.  Compose the exact translation, Galilean, and parabolic observer maps
along the ancestry of each event.  Only finitely many compact cylinders are
prescribed, so local convergence permits one common terminal index on which
all approximations hold.  Normalize the pressures on overlaps and extend the
gauges over their finite union.  This gives one realization of the whole
finite family and proves the final assertion.
\end{proof}

\begin{lemma}[Noncompact observer marks give a first geometric output]
\label{lem:noncompact-observer-mark-output}
Consider a simultaneously realized family satisfying the normalized bounds
\eqref{eq:target-window-bounds}.  Local velocity, normalized-pressure,
suitable-defect, and endpoint-energy coordinates are precompact on the fixed
countable cylinder exhaustion.  Consequently, loss of compactness of the
complete marked family has a first coordinate of one of the following types:
\begin{enumerate}[label=\textup{(\roman*)},leftmargin=2.2em]
\item normalized centers escape, producing an exterior, diffuse,
pressure/noncompact, or separated recentering datum;
\item logarithmic scale increments escape every bounded interval, producing
a sequence-realized logarithmic critical-tail datum;
\item center and logarithmic-scale increments are tight, and the shortest
bounded-variation interpolation in the finite-dimensional observer group
makes the translation and dilation currents tight signed marks.
\end{enumerate}
\end{lemma}

\begin{proof}
The local energy inequality and \eqref{eq:target-window-bounds} give compact
endpoint-energy coordinates and weak-* compact suitable-defect measures.
Strong local \(L^3\) compactness gives the local nonlinear terms, while the
Calder\'on--Zygmund/harmonic decomposition gives the normalized-pressure
coordinates.  Thus only center and logarithmic-scale parameters can fail
tightness before one of these PDE coordinates has already supplied the first
failed coordinate.

Unbounded center displacement is recorded at its first escaping cylinder and
gives alternative~\textup{(i)}.  Unbounded logarithmic displacement is tested
on the annular cylinder exhaustion; diagonal extraction with its realization
coordinate gives alternative~\textup{(ii)}.  In the complementary case join
successive observer parameters by shortest bounded-variation paths.  Their
distributional derivatives are finite signed measures with total variation
controlled by the endpoint displacements.  Pairing them with the fixed smooth
cutoff gradients therefore gives tight translation and dilation currents.
Diagonal compactness over the test family proves alternative~\textup{(iii)}.
\end{proof}

\begin{theorem}[Reduction of transverse active families]
\label{thm:transverse-active-family-reduction}
Let \(\mathcal A_N\) be simultaneously realized active antichains with
\(|\mathcal A_N|\to\infty\).  After an ordered subsequence selection, one of
the following occurs:
\begin{enumerate}[label=\textup{(\roman*)},leftmargin=2.2em]
\item a first center, pressure, exterior, defect, or logarithmic-scale
coordinate produces a sequence-realized datum of the corresponding type;
\item for every \(\delta>0\), all but a uniformly bounded number of events
belong to finitely many \(\delta\)-clusters in center--scale space; each
nested sequence of occupied clusters has one local profile limit, and its
canonical observer interpolation produces a compact coalescent marked path.
\end{enumerate}
Thus the transverse alternative yields a first classified output or returns
to the same marked-current analysis as a longitudinal active path.
\end{theorem}

\begin{proof}
Apply \cref{lem:noncompact-observer-mark-output}.  Its first two alternatives
give clause~\textup{(i)}, so assume that centers and logarithmic scales lie in
a fixed compact parameter set.  After one normalization, all cylinders lie in
a fixed parent cylinder and their radii lie in \([\lambda,1]\) for some
\(\lambda>0\).

Fix \(\delta>0\) and choose a maximal subfamily whose core cylinders are
pairwise separated by at least \(\delta\) in the normalized parabolic metric.
The radii lower bound gives a uniform overlap constant.  If the quotient
activity of a selected sequence tends to zero, its compact limit is spatially
homogeneous and hence target-null.  Otherwise every selected member has a
common positive activity floor.  The local Poincar\'e estimate in
\cref{lem:target-window-critical-price} assigns each such core a fixed
positive normalized viscous price.  Since the cores coexist in one
realization and have bounded overlap, their prices are bounded by the local
energy on one fixed enlargement of the parent cylinder.  Therefore the
cardinality of the separated subfamily is bounded independently of \(N\).

Maximality places every remaining event in one of finitely many
\(\delta\)-clusters.  Let \(\delta_k\downarrow0\) and diagonalize.  Local
compactness and pressure-gauge compatibility show that each nested sequence
of occupied clusters has one local profile limit; two inequivalent limits
would remain separated at some fixed \(\delta_k\).  The bounded-variation
observer interpolation from
\cref{lem:noncompact-observer-mark-output} retains the complete oriented
currents and produces the compact coalescent marked path in clause
\textup{(ii)}.  The two clauses exhaust the antichain alternative.
\end{proof}

\begin{lemma}[Exact signed balance along the moving root]
\label{lem:exact-signed-moving-root-balance}
Let \(I_j=[\tau_j,\tau_{j+1})\) be consecutive represented intervals of one
absolutely continuous chart, and let \(\phi\in C_c^\infty\) be nonnegative
and equal to one on the retained core.  Set
\[
 K_j^-=\lim_{\tau\downarrow\tau_j}
    \int\frac{|V(y,\tau)|^2}{2}\phi(y)\,dy,
 \qquad
 K_j^+=\lim_{\tau\uparrow\tau_{j+1}}
    \int\frac{|V(y,\tau)|^2}{2}\phi(y)\,dy,
 \qquad
 D_j=\nu\int_{I_j}\int|\nabla V|^2\phi,
\]
and let \(M_j\ge0\) be the suitable-energy defect on the represented cell.
Then
\begin{equation}\label{eq:one-moving-cell-balance}
 K_j^+-K_j^-+D_j+M_j=H_j,
\end{equation}
where the signed supply is
\begin{align}\label{eq:complete-moving-root-supply}
 H_j={}&\int_{I_j}\int
 \left(\frac{|V|^2}{2}+P\right)V\cdot\nabla\phi
 +\nu\int_{I_j}\int\frac{|V|^2}{2}\Delta\phi \notag\\
 &-\int_{I_j}a(\tau)\int\frac{|V|^2}{2}
       \bigl(\phi+y\cdot\nabla\phi\bigr)\,dy\,d\tau
 -\int_{I_j}\int\frac{|V|^2}{2}b(\tau)\cdot\nabla\phi.
\end{align}
All pressure, translation, dilation, and cutoff terms are signed.  If the
endpoint trace of one cell is the initial trace of the next, then
\begin{equation}\label{eq:signed-moving-root-gluing}
 \sum_{j=0}^{N-1}(D_j+M_j)
 =K_0^- -K_{N-1}^+ +\sum_{j=0}^{N-1}H_j .
\end{equation}
\end{lemma}

\begin{proof}
Test \eqref{eq:renormalized-NS} against \(V\phi\).  Incompressibility turns
convection and pressure into the first flux in
\eqref{eq:complete-moving-root-supply}; integration by parts gives the
viscous production and cutoff term.  Finally,
\[
 \int V\cdot(V+y\cdot\nabla V)\phi
 =-\int\frac{|V|^2}{2}(\phi+y\cdot\nabla\phi),
 \qquad
 \int V\cdot(b\cdot\nabla V)\phi
 =-\int\frac{|V|^2}{2}b\cdot\nabla\phi.
\]
The suitable local-energy inequality becomes equality after adjoining its
nonnegative Radon defect.  This proves
\eqref{eq:one-moving-cell-balance}.  Summing consecutive cells cancels each
common energy trace with the opposite orientation and gives
\eqref{eq:signed-moving-root-gluing}.  No partition of test functions is used,
so dissipation is not canceled in this gluing.
\end{proof}

\begin{lemma}[Scale-critical balance on nested regenerated cells]
\label{lem:scale-critical-currency-balance}
Let \(\psi_0\ge\psi_1\ge\cdots\ge\psi_N\ge0\) be compatible physical
localizations for a nested regenerated family with radii
\(r_0>r_1>\cdots>r_N\), and put \(w_j=r_j^{-1}\).  Denote by
\(\mathcal J_j\) the complete oriented local-energy carrier through the
stopping surface of \(\psi_j\), and set
\[
 \mathcal P_j:=w_j\Bigl[
   \nu\iint |\nabla u|^2(\psi_j-\psi_{j+1})
   +\mu_{\rm loc}(\psi_j-\psi_{j+1})\Bigr]\ge0.
\]
Then \(\mathcal P_j\) is invariant under Navier--Stokes scaling and
\begin{equation}\label{eq:scale-critical-currency-balance}
 \sum_{j=0}^{N-1}\mathcal P_j
 =w_0\mathcal J_0-w_{N-1}\mathcal J_N
  +\sum_{j=1}^{N-1}(w_j-w_{j-1})\mathcal J_j.
\end{equation}
Consequently, a positive mean normalized price has precisely the same linear
growth as the weighted scale current on the right.  If the mean price tends
to zero, a selected normalized shell has zero gradient and zero suitable
defect, unless a first compactness, pressure, observer, or realization
coordinate has already failed.
\end{lemma}

\begin{proof}
The physical local-energy equality on the shell
\(\psi_j-\psi_{j+1}\) gives
\[
 \nu\iint |\nabla u|^2(\psi_j-\psi_{j+1})
 +\mu_{\rm loc}(\psi_j-\psi_{j+1})
 =\mathcal J_j-\mathcal J_{j+1},
\]
with the common interface carrying opposite orientations.  Multiply by
\(w_j\), sum in \(j\), and apply discrete summation by parts to obtain
\eqref{eq:scale-critical-currency-balance}.

Under the parabolic change of variables based at radius \(r_j\), physical
viscous dissipation and the suitable defect both equal \(r_j\) times their
normalized values.  Multiplication by \(w_j=r_j^{-1}\) therefore makes
\(\mathcal P_j\) scale invariant.  If
\(N^{-1}\sum_{j<N}\mathcal P_j\to0\), choose a shell whose price does not
exceed the mean.  Normalize that shell.  Compactness and lower
semicontinuity give zero gradient and zero defect in its limit.  Failure of
any limit interface is, by the fixed first-coordinate order, the stated
marked output.  If the mean stays positive, the identity itself gives the
same linear growth for the weighted scale current.
\end{proof}

\begin{remark}[Physical energy and critical expenditure]
For a cell of radius \(r_j\), a fixed normalized price \(d_*\) costs only
\(r_jd_*\) in physical energy.  When \(r_{j+1}\le r_j/4\), the series
\(\sum_jr_jd_*\) converges.  Equation
\eqref{eq:scale-critical-currency-balance} retains the missing factor
\(r_j^{-1}\): persistent normalized expenditure appears as persistent scale
current, which is compactified together with the observer and converted into
a realized logarithmic tail.
\end{remark}

\begin{lemma}[Sublinear critical current cannot sustain the sparse path]
\label{lem:sublinear-critical-current-closure}
Let \((C_j)\) be a simultaneously realized finite block of a sparse marked
current path, selected so that
\(C_{j+1}\Subset C_j/2\) and \(r_{j+1}\le r_j/4\).  If the complete signed
moving-root supply in \eqref{eq:complete-moving-root-supply} is \(o(N)\) on
blocks of length \(N\), then the path is impossible or has a first marked output of
\textup{(S3)}.
\end{lemma}

\begin{proof}
Interpolate the selected cells by the absolutely continuous chart carried by
the marked observer coordinates, and choose consecutive half-open represented
intervals.  The endpoint traces and pressure gauges agree because
concatenation is part of the marked record.  Choose \(\phi=1\) on the fixed
target core.  By \cref{lem:target-window-critical-price},
\(D_j+M_j\ge D_j\ge\nu d_*\) on every interval.  The normalized local energy
bounds give \(0\le K_j^\pm\le K_*\), uniformly in \(j\).

Apply \eqref{eq:signed-moving-root-gluing}.  If the complete signed supply is
sublinear, then
\[
 N\nu d_*
 \le\sum_{j<N}(D_j+M_j)
 \le2K_*+o(N),
\]
which is impossible as \(N\to\infty\).  If the moving chart, endpoint trace,
pressure gauge, or interpolation needed for the identity loses compactness,
its first failure is exactly the marked output in \textup{(S3)}.  Thus there
is no unrecorded nested-core alternative: the price is paid on disjoint
represented time cells, and all transport across their moving spatial
boundaries is retained in \(H_j\).
\end{proof}

\begin{definition}[Concentration-free realized critical tail]
\label{def:concentration-free-critical-tail}
A sequence-realized logarithmic critical tail in the sense of
\cref{def:sequence-realized-log-tail} is concentration free when every fixed
positive-threshold active observer has first been retained as an actual
successor.  The residual annular state therefore contains only diffuse
velocity mass, a pressure or suitable-energy defect, or a coherent strong
annular limit.  This is an ordered definition: an active observer is selected
before the residual tail is formed.
\end{definition}

\begin{lemma}[Realization of logarithmic shell support]
\label{lem:logarithmic-shell-support-realization}
Let \((W_n,P_n,a_n,b_n)\) be represented normalized states on time windows
modeled on one compact interval \(J_*\), and suppose that for some
\(R_n\to\infty\) and \(\varepsilon_*>0\),
\begin{equation}\label{eq:nonzero-spacetime-tail-mass}
 M_n:=\int_{J_*}\int_{|y|>R_n}|W_n|^3\,dy\,ds\ge\varepsilon_*.
\end{equation}
The probability measures
\[
 d\nu_n=M_n^{-1}|W_n|^3
 \mathbf 1_{J_*\times\{|y|>R_n\}}\,dy\,ds
\]
have, after dyadic shell decomposition and passage to a subsequence, a weak-*
limit on a compact metrizable logarithmic shell-state space.  Every point in
the support of that limit is realized by actual translated unit-shell
cylinders from the same sequence \((W_n,P_n)\).  The construction uses spatial
translation and bounded time recentering only; the shell radius is a
logarithmic coordinate and is not used as a Navier--Stokes dilation.

For every realized support point, either a first pressure, observer-current,
local-energy, defect, or domain coordinate fails, or the point is represented
on a fixed model cylinder by a suitable local solution with the corresponding
compact-cylinder bounds.
\end{lemma}

\begin{proof}
Decompose \(\{|y|>R_n\}\) into the shells
\[
 A_{n,k}=\{2^kR_n<|y|\le2^{k+1}R_n\},\qquad k\ge0,
\]
and cover each occupied spacetime shell by a family of unit parabolic
cylinders with uniformly bounded overlap.  Translate the center of each
selected cylinder to the origin and recenter its time interval inside a fixed
model cylinder \(Q_{\rm sh}\).  Translation, bounded time recentering, and
subtraction of the pressure mean preserve the equation, incompressibility,
the local energy inequality, and local pressure reconstruction.  In
particular, no factor depending on \(2^kR_n\) enters the local equation.

Record the logarithmic shell index, translated center, local velocity and
normalized pressure on a countable compact-cylinder exhaustion, suitable
defect measures, observer currents, and the physical realization coordinate.
Compactify each scalar or finite-measure coordinate by adjoining its first
failure value.  The represented part is sequentially compact by the local
energy bounds, the pressure decomposition, and diagonal extraction.  The
countable product is therefore compact and metrizable.  Pushing \(\nu_n\)
forward through the shell cover and taking a weak-* subsequence gives a
probability measure \(\Pi_{\rm sh}\).

Fix \(z\in\operatorname{supp}\Pi_{\rm sh}\) and a nested neighborhood basis
\(U_m\downarrow\{z\}\).  Since \(\Pi_{\rm sh}(U_m)>0\), weak-* convergence and
the portmanteau inequality give \(\nu_{n}(U_m)>0\) along a further subsequence.
Choose an actually occupied shell cylinder whose state lies in \(U_m\), with
the indices increasing in \(m\).  These genuine translated states converge to
\(z\), proving sequence realization.  If \(z\) lies on a compactified failure
face, the first such coordinate is the asserted output.  Otherwise all local
coordinates remain represented and their distributional limits give the
suitable local solution claimed in the statement.
\end{proof}

\begin{lemma}[Strong compactness on a realized shell]
\label{lem:realized-shell-strong-compactness}
Let \((W_n,P_n,a_n,b_n)\) be a realized shell sequence on \(Q_{\rm sh}\), and
assume that the pressure gauge, represented equation, local energy inequality,
and compact-cylinder bounds have passed on every \(Q'\Subset Q_{\rm sh}\).
Then exactly one of the following occurs:
\begin{enumerate}[label=\textup{(\roman*)},leftmargin=2.2em]
\item the local modulation currents, an upper local critical bound, a pressure
or defect coordinate, or a domain coordinate has a first failure;
\item after a subsequence, \(W_n\to W\) strongly in
\(L^3_{\rm loc}(Q_{\rm sh})\), and the local Calder\'on--Zygmund pressure parts
converge strongly in \(L^{3/2}_{\rm loc}\).
\end{enumerate}
Thus a non-Dirac Young or viscous-defect limit is recorded at a first concrete
coordinate and cannot remain an unrepresented tail.
\end{lemma}

\begin{proof}
Assume that alternative~\textup{(i)} does not occur and fix
\(Q'\Subset Q''\Subset Q_{\rm sh}\).  The local energy inequality gives
\[
 W_n\ \text{ bounded in }\
 L^\infty_sL^2_y(Q'')\cap L^2_sH^1_y(Q'').
\]
The equation \eqref{eq:renormalized-NS}, the local pressure bound, and the
passed modulation-current coordinate give
\[
 \partial_sW_n\ \text{ bounded in }
 L^{4/3}_sW^{-1,3/2}_y(Q'').
\]
Indeed, viscosity belongs to \(L^2_sH^{-1}_y\); convection is controlled in
\(L^{5/3}_sW^{-1,5/3}_y\) by the energy interpolation
\(W_n\in L^{10/3}(Q'')\); the normalized pressure is bounded in
\(L^{3/2}\) and therefore its gradient in \(W^{-1,3/2}\);
and the two observer terms are exactly the retained modulation-current
coordinate.  Aubin--Lions gives strong convergence in \(L^2(Q')\).  The same
energy interpolation supplies a uniform \(L^{10/3}(Q')\) bound, so interpolation
upgrades the convergence to strong \(L^3(Q')\).  A diagonal choice over
\(Q'\Subset Q_{\rm sh}\) proves the velocity assertion.

On nested balls write \(P_n=P_n^{\rm loc}+H_n+c_n(s)\), where
\(P_n^{\rm loc}\) is the Calder\'on--Zygmund pressure generated by
\(W_n\otimes W_n\), \(H_n\) is harmonic, and \(c_n\) is the pressure gauge.
Strong \(L^3\) convergence gives strong \(L^{3/2}\) convergence of
\(P_n^{\rm loc}\).  Harmonic interior estimates compactify \(H_n\) after the
fixed mean normalization; failure of that compactness is precisely the first
pressure coordinate in alternative~\textup{(i)}.  Finally, lower
semicontinuity in the local energy inequality leaves a nonnegative viscous
defect.  If it is nonzero, the first positive test in the fixed countable test
family records the defect in alternative~\textup{(i)}; if every such test
vanishes, the defect is zero.  This proves the dichotomy.
\end{proof}

\begin{lemma}[Bounded-origin rigidity of coherent logarithmic tails]
\label{lem:bounded-origin-coherent-tail-rigidity}
Let \(W\) be a coherent sequence-realized logarithmic critical tail obtained
from uniformly bounded record-point representatives.  Suppose that all
pressure, defect, observer, and domain-realization failures have already been
retained as first-coordinate outputs.  Then \(W\) has a locally bounded
representative on full macroscopic cylinders, including cylinders meeting the
origin.  Moreover its amplitude-normalized logarithmic field
\[
 Z(\rho,\omega,s):=e^\rho W(e^\rho\omega,s)
\]
has the following property: every sequence-realized element \(Z_*\) of its
logarithmic/time-translation hull either produces a first-coordinate output or
satisfies
\begin{equation}\label{eq:coherent-tail-origin-decay}
 Z_*(\rho,\cdot,\cdot)\longrightarrow0
 \quad\text{locally as }\rho\to-\infty.
\end{equation}
Consequently no nonzero homogeneous, log-periodic, or aperiodic minimal
coherent tail can retain quotient activity.
\end{lemma}

\begin{proof}
Let \(W_n\) be the affine-normalized macroscopic representatives producing
the coherent tail.  The record bound gives \(|W_n|\le2M_*\) after choosing
the minimizing homogeneous velocity gauges in the closed ball of radius
\(M_*\).  The domain-realization coordinate ensures that every fixed full
macroscopic cylinder lies in the domain of \(W_n\) for all sufficiently large
\(n\); otherwise that coordinate is already the asserted first output.
Banach--Alaoglu and a diagonal extraction therefore give a weak-* locally
bounded limit \(W^{\rm full}\) on full cylinders.  Strong convergence on
punctured annuli identifies its restriction there with \(W\).

Consider a sequence-realized amplitude-normalized translate
\(Z(\rho+a_k,\omega,s+b_k)\).  The realization mark either remains compact,
in which case the same full-cylinder argument gives a locally bounded
representative \(W_*\), or its first failed pressure, observer, defect, or
domain coordinate is retained.  In the compact case, with \(r=e^\rho\),
\[
 |Z_*(\rho,\omega,s)|=r|W_*(r\omega,s)|\le C_Kr
\]
on every compact angular--time set \(K\).  This proves
\eqref{eq:coherent-tail-origin-decay}.

If \(Z_*\) is independent of \(\rho\), the decay forces \(Z_*=0\).  If it has
a nonzero period \(L\), then
\(Z_*(\rho)=Z_*(\rho-kL)\); choosing the sign of \(k\) so that
\(\rho-kL\to-\infty\) again gives \(Z_*=0\).  Finally, let \(\mathcal M\) be
a minimal aperiodic subsystem of the sequence-realized logarithmic hull.
Every member obeys the decay alternative unless a first-coordinate output has
already occurred.  Hence zero belongs to the hull of every member.  Since
\(\{0\}\) is a closed invariant subsystem, minimality gives
\(\mathcal M=\{0\}\).  Each case contradicts retained nonzero quotient
activity.
\end{proof}

\begin{lemma}[Persistent current gives a successor or a critical tail]
\label{lem:persistent-current-support-alternative}
Suppose a sparse marked current path has positive mean critical price and is
not covered by \cref{lem:sublinear-critical-current-closure}.  Then, after a
subsequence, exactly one of the following occurs:
\begin{enumerate}[label=\textup{(C\arabic*)},leftmargin=2.8em]
\item positive limiting current is supported on a bounded logarithmic window
which contains a target-retaining active observer;
\item after all such observers have been removed, a nonzero
concentration-free sequence-realized critical tail remains;
\item a first pressure, exterior, defect, observer, or realization coordinate
loses compactness, or a nonzero residual pressure/defect coordinate survives,
and is the marked output in \textup{(S3)}.
\end{enumerate}
\end{lemma}

\begin{proof}
For a block of length \(N\), put mass \(N^{-1}H_j^+\) on the logarithmic
interval between \(r_j\) and \(r_{j+1}\).  By
\eqref{eq:signed-moving-root-gluing},
\[
 \sum_{j<N}H_j^+
 \ge\sum_{j<N}H_j
 =\sum_{j<N}(D_j+M_j)+K_{N-1}^+-K_0^-
 \ge N\nu d_*-2K_*.
\]
Thus these positive measures have total mass bounded below.  Decompose
\(H_j\) into its already recorded pressure, nonlinear transport, cutoff,
translation, and dilation coordinates.  Since the positive part of their sum
is bounded by the sum of their absolute values, one coordinate carries a
fixed fraction of the mass.

If this coordinate is not tight on the countable cylinder-and-test
exhaustion, select its first failure; this is \textup{(C3)}.  Under tightness,
\cref{lem:marked-finite-simultaneous-realization} realizes every finite
collection of its annular tests in one member of the original terminal
sequence.  Diagonal compactness therefore gives a sequence-realized support
state, not merely an abstract current measure.  If positive support meets the
closed active set, the first such window gives \textup{(C1)}.  Otherwise
remove the active windows.  The remaining positive mass prevents the
projective annular state from vanishing and gives \textup{(C2)}.  The three
outputs are disjoint by their order and exhaustive by tightness.
\end{proof}

\begin{lemma}[Concentration-free critical tails have no terminal component]
\label{lem:concentration-free-tail-closure}
A concentration-free sequence-realized critical tail arising from a
finite-energy suitable solution is target-null or produces a marked output
already listed in \textup{(S3)}.  In particular, it cannot be the terminal
support of the persistent current in
\cref{lem:persistent-current-support-alternative}.
\end{lemma}

\begin{proof}
Apply the following alternatives in order on the logarithmic annular
projective family.

If velocity mass is not tight on bounded logarithmic windows, choose the
first scale at which the Galilean-reduced quantity
\eqref{eq:galilean-reduced-C} reaches a fixed number below the
\(\varepsilon\)-regularity threshold.  The point-selection proof of
\cref{prop:first-bad-scale-descendant} gives an actual target-retaining
successor, a homogeneous quotient, or a non-affine harmonic pressure tail.
The first alternative is absent by concentration-free selection, the second
is target-null, and the third is the pressure output in \textup{(S3)}.

Assume next that the velocity is logarithmically tight.  Apply
\cref{lem:logarithmic-shell-support-realization} to the normalized annular
mass.  Thus every point of the limiting shell support is represented by an
actual translated cylinder from the original sequence; no abstract support
measure is substituted for a Navier--Stokes state.  On each represented shell,
\cref{lem:realized-shell-strong-compactness} gives a first pressure,
modulation, defect, or domain coordinate in \textup{(S3)}, or strong local
\(L^3\) compactness.  In the latter case the Young measure is Dirac.  If a
positive compact quotient activity survives, its first value on the fixed
dyadic activity ladder is an actual target-retaining successor, contrary to
the concentration-free selection.  Consequently the remaining strongly
compact shell state has zero compact activity.

Pass now to the pressure and suitable-energy coordinates.  On a fixed
annulus, use the actual pressure in the local energy identity and subtract
only its time-dependent gauge.  Strong velocity convergence passes the
transport term.  Lower semicontinuity leaves a nonnegative viscous defect.  A
positive convex test detecting that defect supplies a first defect coordinate
in \textup{(S3)}.  If all positive convex tests vanish, the matrix-valued
viscous defect vanishes; polarization removes its signed coordinates.  The
local pressure decomposition then writes the remaining pressure as a compact
Calder\'on--Zygmund part plus a harmonic part.  The former converges strongly
in \(L^{3/2}\); a noncompact non-affine harmonic part is precisely the
pressure output in \textup{(S3)}.  Hence outside \textup{(S3)} every pressure
and suitable-energy defect coordinate vanishes.

It remains to consider a coherent strong annular state.  Apply
\cref{lem:bounded-origin-coherent-tail-rigidity}.  A failure of full-cylinder
or amplitude-normalized hull realization is the first observer, pressure,
defect, or domain coordinate in \textup{(S3)}.  When all realization
coordinates pass, that lemma excludes the homogeneous, log-periodic, and
aperiodic minimal coherent alternatives by bounded-origin decay.  Thus the
coherent component is target-null as well.  These alternatives exhaust weak
annular compactness, and the assertion follows.
\end{proof}

\begin{theorem}[Sparse moving-root path closure]
\label{thm:sparse-moving-root-path-exclusion}
Every sparse marked current path arising from the collapsed same-origin
residual reaches a typed local output in \textup{(F3)} or is impossible.
\end{theorem}

\begin{proof}
Sublinear current is excluded by
\cref{lem:sublinear-critical-current-closure}.  The weighted identity
\cref{lem:scale-critical-currency-balance} shows that the surviving quantity is
the scale-critical current: vanishing mean price gives a target-null
zero-gradient shell, while positive mean price gives persistent weighted
current.  On the remaining path apply
\cref{lem:persistent-current-support-alternative}.  Alternative
\textup{(C3)} is the retained typed output required by
\textup{(J4)} and \textup{(S3)}, and gives the first conclusion of the
theorem.  Assume henceforth that it does not occur.  Alternative
\textup{(C2)} is impossible by
\cref{lem:concentration-free-tail-closure}.

Thus every sufficiently long consecutive block has its positive limiting
current supported on regenerated active observers.  Retain each observer,
its complete oriented current, and its realization coordinate.  Compactness
gives the closed marked path space of
\cref{lem:compact-closed-marked-path-space}.  Poincar\'e recurrence for the
successor shift and exact composition of the observer maps then give a
nonzero sequence-realized critical tail by
\cref{thm:marked-recurrence-critical-tail}.

Apply the ordered annular reduction in the proof of
\cref{lem:concentration-free-tail-closure}.  A target-null component is
impossible because the recurrent cross-section retains the activity floor; a
pressure, defect, exterior, or realization failure is \textup{(F3)}.  A
nonzero concentration-free component is excluded by that lemma.  The only
remaining output is a target-retaining first-bad observer.  Append this
observer, with its complete marks, to the successor path and repeat.

If this procedure terminates, one of the preceding alternatives applies.  If
it does not, choose at every stage the first mesoscopic scale at which the
fixed activity threshold is reached.  All smaller scales are then below that
threshold.  The translation-modulus argument in the second paragraph of the
proof of \cref{lem:concentration-free-tail-closure} shows that the associated
minimal-scale blow-downs are strongly \(L^3_{\rm loc}\)-compact; the pressure
and defect alternatives have already been placed in \textup{(F3)}.  Their
diagonal limit is therefore a nonzero coherent logarithmic tail.  The
bounded-origin/minimal-hull argument in the final paragraph of that proof
forces this tail to be zero, a contradiction.  Hence endless regeneration is
also impossible.  This eliminates the persistent-current case and proves the
theorem.
\end{proof}

\begin{theorem}[Closure of the collapsed same-origin residual]
\label{thm:collapsed-same-origin-residual-exclusion}
Every collapsed same-origin residual in
\cref{def:collapsed-same-origin-residual} either reaches a typed local output
in \textup{(F3)} or is impossible.  Consequently no collapsed residual
remains after the typed \textup{(F3)} outputs have been removed in the ordered
current alternative.
\end{theorem}

\begin{proof}
Apply \cref{lem:density-current-sparse-reduction} to its infinite active
successor path.  Positive-density and covariantly balanced paths contradict
covariant rigidity.  The only remaining path is sparse, and is impossible by
\cref{thm:sparse-moving-root-path-exclusion}.  Every failure used in that
proof is one of the typed successors in \textup{(J4)} and therefore reaches
\textup{(F3)}.  If no such output occurs, the sparse-path contradiction
applies.  Thus no pressure, exterior, scale, defect, or realization loss is
discarded, and the two alternatives in the statement are exhaustive.
\end{proof}

\begin{proof}[Proof of \cref{thm:persistent-current-compactification}]
Apply \cref{thm:finite-depth-NS-source-alternative} to the compatible marked
records cut from the canonical orbit.  The retained quotient floor excludes
\textup{(F1)}.  Alternative \textup{(F2)} is impossible because a finite
nonnegative Radon measure cannot dominate a fixed positive charge on
infinitely many pairwise disjoint events.  In \textup{(F3)}, choose the first
failed coordinate in the fixed countable exhaustion.  Its detecting test,
physical realization coordinate, and pressure gauge survive the diagonal
selection; \cref{prop:first-unbounded-activity-scale} and
\cref{lem:persistent-current-support-alternative} therefore give the
sequence-realized datum in clause~\textup{(i)} of the theorem.

It remains to consider the ordinary inverse-limit output \textup{(F4)}.
Compact nonexpanding recurrence is target-null by
\cref{prop:compact-nonexpanding-output-closure}, and sublinear signed current
is impossible by \cref{lem:sublinear-critical-current-closure}.  The ordered
reduction \cref{cor:reduction-to-collapsed-residual} consequently leaves the
collapsed same-origin path, with every observer and oriented current retained
as a mark.  A first loss of a marked coordinate again gives clause
\textup{(i)}.  If no coordinate is lost, the closed path space in
\cref{lem:compact-closed-marked-path-space} carries a recurrent successor
path, and \cref{thm:marked-recurrence-critical-tail} produces the nonzero
sequence-realized logarithmic tail in clause~\textup{(ii)}.  These choices
are ordered, disjoint, and exhaustive, and every limit retains the common
inverse-limit realization by the original solution.
\end{proof}

\subsection{Terminal concentration and exclusion of the record profile}
\label{app:terminal-record-exclusion}

The record-point extraction reduces a finite-time breakdown of a classical
solution to a bounded mild ancient profile.  We now close that profile without
assuming the existence of a separately constructed continuous orbit.  All
limits below are taken with the pressure gauge, suitable-energy defect, and
realization coordinate in \cref{def:marked-successor-record}.

\begin{definition}[Terminally indecomposable profile]
\label{def:terminally-indecomposable-profile}
A bounded mild ancient profile \(U\), represented modulo spatially homogeneous
velocity and time-dependent pressure, is terminally indecomposable if it
retains positive quotient activity on one compact cylinder and has none of the
following outputs:
\begin{enumerate}[label=\textup{(A\arabic*)},leftmargin=2.8em]
\item an actual target-retaining descendant;
\item a second asymptotically separated target-retaining profile;
\item a diffuse exterior remainder;
\item a noncompact velocity or non-affine pressure-gauge remainder;
\item a nonzero sequence-realized logarithmic critical tail.
\end{enumerate}
The alternatives are tested at every positive rational activity threshold.
\end{definition}

\begin{lemma}[Indecomposability produces a backward \(L^3\) sequence]
\label{lem:indecomposable-backward-L3}
Let \(U\) be a terminally indecomposable bounded mild ancient profile obtained
from the record-point extraction.  Then there are \(s_k\downarrow-\infty\)
such that
\begin{equation}\label{eq:indecomposable-backward-L3}
 \sup_k\|U(\cdot,s_k)\|_{L^3(\mathbb R^3)}<\infty.
\end{equation}
\end{lemma}

\begin{proof}
Assume the contrary.  Then
\(\int_{\mathbb R^3}|U(x,s)|^3dx\to\infty\) as \(s\to-\infty\).
Choose \(s_k\downarrow-\infty\) and \(R_k\uparrow\infty\) so that
\[
 1\le\int_{|x|>R_k}|U(x,s_k)|^3dx\le2.
\]
Cover the exterior slice, together with a unit backward time interval, by a
fixed bounded-overlap family of unit parabolic cylinders.  If one cylinder
has quotient activity bounded below, recentering there and using the local
pressure decomposition gives either an actual target-retaining descendant or
a first non-affine pressure/compactness output.  These are
\textup{(A1)} or \textup{(A4)}.

Suppose instead that the supremum of the unit-cylinder activity tends to zero.
The normalized exterior restrictions are then diffuse unless their mass first
becomes visible at a larger multiplicative scale.  The former is
\textup{(A3)}.  In the latter case choose the first dyadic scale at which a
fixed quotient threshold below \(\varepsilon_{\rm reg}\) is reached.  The
point-selection and pressure argument of
\cref{prop:first-bad-scale-descendant} produces an actual descendant, a
spatially homogeneous quotient, a pressure/noncompact output, or a realized
critical-tail window.  The homogeneous quotient carries no activity; the
other three outputs are \textup{(A1)}, \textup{(A4)}, and \textup{(A5)}.
Every alternative contradicts terminal indecomposability.  Hence
\eqref{eq:indecomposable-backward-L3} holds.
\end{proof}

\begin{lemma}[No terminally indecomposable active profile]
\label{lem:no-terminally-indecomposable-profile}
No terminally indecomposable bounded mild ancient profile obtained from the
record-point extraction retains positive quotient activity.
\end{lemma}

\begin{proof}
By \cref{lem:indecomposable-backward-L3}, the profile satisfies the hypothesis
of the endpoint ancient Liouville theorem
\cref{thm:AB-endpoint}.  Hence \(U\equiv0\), contradicting its positive
quotient activity.
\end{proof}

\begin{lemma}[Finite separated families reduce to an infinite chain or an
indecomposable profile]
\label{lem:finite-separated-family-reduction}
Let \(U_1,\ldots,U_J\) be a finite family of pairwise asymptotically separated
target-retaining profiles simultaneously realized from one record-point
sequence.  Then some \(U_j\) is terminally indecomposable, or the family
generates an infinite chain of actual target-retaining descendants.
\end{lemma}

\begin{proof}
Assume that no \(U_j\) is terminally indecomposable.  Choose for each profile
its first output in the order \textup{(A1)}--\textup{(A5)}.  Outputs
\textup{(A3)}--\textup{(A5)} are retained by their own coordinates and enter
the diffuse, pressure, or critical-tail alternatives below.  On the remaining
profiles, every vertex has an actual target-retaining successor.

Whenever such a successor is asymptotically separated from all profiles
already selected, adjoin it to the family.  If this happens indefinitely, the
successors form the asserted infinite chain.  Otherwise one obtains a maximal
finite family.  Draw an arrow from each vertex to the member equivalent to its
selected successor.  Every vertex has an outgoing arrow, so the finite
directed graph contains a cycle.  Periodically repeating that cycle gives an
infinite actual descendant chain.  Simultaneous realization and pressure-gauge
compatibility are preserved under composition by
\cref{lem:marked-finite-simultaneous-realization}.
\end{proof}

\begin{theorem}[Terminal concentration alternative]
\label{thm:terminal-concentration-alternative}
Let \(U\) be a nonzero bounded mild ancient profile produced by
\cref{thm:record-point-ancient-extraction}, with its complete record-point
ancestry.  After removal of its spatially homogeneous quotient, one of the
following occurs:
\begin{enumerate}[label=\textup{(T\arabic*)},leftmargin=2.8em]
\item a terminally indecomposable active profile;
\item a finite separated family of active profiles;
\item an infinite chain of actual target-retaining descendants;
\item a diffuse exterior remainder;
\item a noncompact velocity or non-affine pressure remainder;
\item a nonzero sequence-realized logarithmic critical tail.
\end{enumerate}
Every output is realized by the original record-point sequence; a measure or
pressure remainder is not declared to be an ordinary velocity profile.
\end{theorem}

\begin{proof}
Apply \cref{prop:backward-concentration-remainder} to \(U\).  Its backward
\(L^3\) alternative makes \(U=0\) by \cref{thm:AB-endpoint}.  Its homogeneous
alternative is target-null.  A positive compact oscillation gives an actual
target-retaining profile.  On the diffuse output, apply
\cref{prop:first-bad-scale-descendant}; this gives another actual descendant,
a homogeneous quotient, a pressure/noncompact remainder, or the first
critical-tail window.

Iterate this procedure, always choosing the first output in the displayed
order and retaining the realization marks.  If it stops at one active profile
with no further output, that profile is terminally indecomposable and
\textup{(T1)} holds.  If only finitely many inequivalent active profiles are
produced, they form \textup{(T2)}.  Infinitely many successive active profiles
give \textup{(T3)}.  Loss of exterior tightness, local velocity or pressure
compactness, or logarithmic scale tightness gives respectively
\textup{(T4)}--\textup{(T6)}.  The countable cylinder-and-test exhaustion in
\cref{def:finite-depth-localized-event-record} makes the first-failure
selection exhaustive, and its inverse-limit coordinate gives simultaneous
realization.  Thus no additional output is possible.
\end{proof}

\begin{theorem}[Exclusion of record-point ancient profiles]
\label{thm:record-point-profile-exclusion}
No nonzero bounded mild ancient profile with the complete ancestry and
uniform record bound furnished by
\cref{thm:record-point-ancient-extraction} exists.
\end{theorem}

\begin{proof}
Assume that such a profile exists and apply
\cref{thm:terminal-concentration-alternative}.  Alternative \textup{(T1)} is
excluded by \cref{lem:no-terminally-indecomposable-profile}.  Alternative
\textup{(T2)} reduces by
\cref{lem:finite-separated-family-reduction} to \textup{(T1)} or
\textup{(T3)}.

For \textup{(T3)}, positive-density and covariantly balanced chains are
excluded by \cref{lem:density-current-sparse-reduction}.  The remaining sparse
marked path reaches a typed first-failure output or is impossible by
\cref{thm:sparse-moving-root-path-exclusion}.  A typed first failure is one of
\textup{(T4)}--\textup{(T6)}, so it is retained rather than discarded.

For \textup{(T4)}, normalize a positive exterior restriction.  Positive
unit-scale quotient activity gives an actual descendant and returns to
\textup{(T2)} or \textup{(T3)}.  Vanishing unit activity gives the first bad
scale in \cref{prop:first-bad-scale-descendant}, hence a homogeneous quotient,
\textup{(T5)}, or \textup{(T6)}.  Thus diffuse mass is not terminal.

In \textup{(T5)}, the local pressure decomposition separates a compact
Calder\'on--Zygmund part from a harmonic remainder.  The far harmonic Taylor
jet is affine by \cref{prop:harmonic-pressure-tail-accessibility} and is
target-null.  A non-affine intermediate remainder has a first bad parabolic
scale and therefore gives an actual descendant or a critical-tail window by
\cref{prop:first-bad-scale-descendant}.  Hence \textup{(T5)} is not terminal.

Finally, remove all active observers from \textup{(T6)}.  A remaining nonzero
tail is excluded by \cref{lem:concentration-free-tail-closure}; if an active
observer remains, it is an actual descendant and returns to
\textup{(T2)}--\textup{(T3)}.  Endless return is the marked regenerated case
in the proof of \cref{thm:sparse-moving-root-path-exclusion} and produces the
same excluded coherent tail.  Every alternative is impossible.  This
contradicts the assumed nonzero profile.
\end{proof}

\begin{proof}[Proof of \cref{thm:critical-tail-rigidity}]
First remove every bounded logarithmic window carrying the fixed quotient
activity threshold.  If none remains, the tail is concentration free and
\cref{lem:concentration-free-tail-closure} shows that it is target-null unless
a first pressure, exterior, scale, trace, defect, or realization coordinate
is visible.  Retain that coordinate with its realizing observer.  The
stopping-time construction \cref{prop:first-bad-scale-descendant} and the
pressure accessibility estimate
\cref{prop:harmonic-pressure-tail-accessibility} give, in the fixed order, an
ordinary active successor, a homogeneous or affine target-null state, or the
next typed marked coordinate.  The first alternative enters the active case
below; the second is already target-null; and indefinite repetition of the
third is the sparse marked path excluded by
\cref{thm:sparse-moving-root-path-exclusion}.  Thus no nonordinary coordinate
is silently promoted to a velocity profile.

Suppose instead that active logarithmic windows remain.  Retain their centers,
scales, Galilean parameters, pressure gauges, defects, and oriented currents.
The marked space is closed by
\cref{lem:compact-closed-marked-path-space}.  Apply
\cref{lem:density-current-sparse-reduction} to the resulting infinite active
successor path.  Positive-density or covariantly balanced recurrence is
target-null by the stationary covariant identity.  In the sparse case the
root and its complete boundary current remain marks, so no spatial average is
used.  The support alternative
\cref{lem:persistent-current-support-alternative} gives a first failed
coordinate, a concentration-free tail, or another active successor.  The
first is reduced at its first visible parabolic scale exactly as in the
preceding paragraph; the second is reduced by
\cref{lem:concentration-free-tail-closure} to the target-null case or the same
ordered first-coordinate reduction; and endless active regeneration is
excluded by
\cref{thm:sparse-moving-root-path-exclusion}.  Thus the rooted and unrooted
recurrence cases are both target-null.

For completeness, the weak annular limits used above exhaust the usual five
forms.  By \cref{lem:logarithmic-shell-support-realization}, every support
point is an actual translated-cylinder limit.  Loss of logarithmic tightness
is the diffuse case.  On a tight shell,
\cref{lem:realized-shell-strong-compactness} either records the first Young,
pressure, modulation, or viscous-defect coordinate, or makes the Young measure
Dirac by strong \(L^3\) convergence.  A coherent homogeneous limit has zero
quotient activity, while coherent log-periodic and aperiodic minimal limits
are zero by \cref{lem:bounded-origin-coherent-tail-rigidity}.  Thus no
additional tail remains.
\end{proof}

\begin{theorem}[Exclusion of bounded regenerated current]
\label{thm:marked-regenerated-current-exclusion}
No uniformly record-bounded marked regenerated-current cascade exists.
Before bounded realization, every failed center, pressure, scale, exterior,
defect, or trace coordinate is retained by
\cref{thm:persistent-current-compactification}.
\end{theorem}

\begin{proof}
The bounded cascade is impossible by
\cref{thm:no-bounded-marked-regenerated-cascade}.  The first-failure assertion
is exactly the ordered coordinate selection in the proof of
\cref{thm:persistent-current-compactification}.
\end{proof}

\begin{remark}[Scope of the marked closure]
The finite-depth estimate closes the record-bounded regenerated-current case
before a logarithmic invariant law is needed.  The quotient observable removes
the accumulated constant Galilean velocity, while exact composition retains
the product scale \(\Lambda_N\).  The required root bound is inherited from
\eqref{eq:uniform-record-bound} by
\cref{lem:record-bound-inherited-by-marked-blocks}.  Pressure created by motion of the root
remains in \(\mathbf q_j\) and in the exact shell balance
\eqref{eq:marked-critical-shell-balance}; it is not used in the amplitude
estimate.
\end{remark}

\subsection{A single forward orbit}
\label{app:single-forward-orbit}

\begin{proposition}[Coherent interpolation of selected observers]
\label{prop:coherent-observer-interpolation}
Let \((x_j,t_j,r_j)\) be the chronologically ordered, target-retaining cells
selected from the record-point sequence and its first-bad-scale descendants.
After passing to a subsequence, they are represented by one absolutely
continuous chart \((X,\Lambda,t)\) on \([0,\infty)\) for which
\[
 t'=\Lambda^2,\qquad \Lambda(\tau_j)=r_j\downarrow0,
 \qquad X(\tau_j)=x_j,
\]
the selected unit intervals have disjoint interiors, and
\eqref{eq:chart-transition-composition} holds.  The assertion is unchanged
when \(|x_j|\to\infty\).
\end{proposition}

\begin{proof}
First pass to a subsequence with \(r_{j+1}\le r_j/4\) and with physical time
images strictly ordered.  On a unit neighborhood of each selected time keep
\(X=x_j\), \(\Lambda=r_j\), and use \(t'=\Lambda^2\).  In the intervening
intervals interpolate \(\log\Lambda\) monotonically and \(X\) linearly after
choosing the interval length large enough that
\(a=\Lambda'/\Lambda\) and \(b=X'/\Lambda\) are integrable.  Define
\(t(\tau)\) by integrating \(t'=\Lambda^2\), adjusting only the lengths of
the inactive intervals; their total physical duration can be made equal to
the prescribed gaps because increasing a represented interval continuously
increases its physical duration from zero to infinity.

The change of variables \eqref{eq:chart}--\eqref{eq:chart-fields} is applied to the
same physical solution at every \(\tau\).  Composition of affine parabolic
changes therefore gives \eqref{eq:chart-transition-composition}
algebraically.  Constant Galilean modes are removed only in the observable
\eqref{eq:oscillatory-mass}, so interpolation does not alter the target on a
selected window.  No bound on \(X\) is used, which covers escaping centers.
\end{proof}

\subsection{Compactness and persistence}
\label{app:orbit-compactness-persistence}

\begin{proposition}[Compactness and retention on the canonical orbit]
\label{prop:canonical-orbit-compactness-retention}
The windows in \cref{prop:coherent-observer-interpolation} satisfy
\eqref{eq:target-window-bounds}.  Their closure is sequentially compact in
\eqref{eq:critical-orbit-compactness}, normalized pressures converge strongly
in \(L^{3/2}_{\mathrm{loc}}\), and the target floor
\eqref{eq:target-window-floor} passes to every limit.
\end{proposition}

\begin{proof}
On every fixed normalized cylinder, the record bound and the first-bad-scale
selection give a uniform velocity bound after an inner-cylinder reduction.
The local energy inequality then gives uniform
\(L^\infty_\tau L^2_y\cap L^2_\tau H^1_y\) bounds.  The equation gives
\(\partial_\tau V\) bounded in
\(L^{4/3}_\tau H^{-1}_y\) after the \(L^1\) bounds for \(a,b\) from
\cref{prop:coherent-observer-interpolation}.  Aubin--Lions compactness and
interpolation yield strong local \(L^3\) convergence.

Write \(P=P_{\rm loc}+H\) on nested balls.  Calder\'on--Zygmund estimates and
strong \(L^3\) velocity convergence give strong \(L^{3/2}\) convergence of
\(P_{\rm loc}\).  After subtracting \((P)_{B_1}\), harmonic interior
estimates compactify \(H\).  If a non-affine harmonic component fails these
bounds, \cref{prop:harmonic-pressure-tail-accessibility} locates its first
parabolic scale and it is selected as the next window, rather than remaining
in the closure under consideration.  Diagonal extraction over the
compact-cylinder exhaustion proves sequential compactness of the closure.
Finally, \(\OscMass_\Gamma\) is continuous under strong \(L^3\) convergence
by \cref{prop:entropy-compactification}, so the positive floor persists.
\end{proof}

\begin{proof}[Proof of \cref{thm:canonical-critical-orbit}]
Apply \cref{thm:record-point-ancient-extraction} at successive record times.
At a terminal singular point, apply the velocity regularity criterion after
each exact constant Galilean transformation.  If the Galilean-reduced
quantity \(C^circ\) in \eqref{eq:galilean-reduced-C} were below its regularity
threshold on all sufficiently small tilted cylinders, the corresponding
Galilean representative would be locally bounded and the original point
would be regular.  Hence \(C^circ\) crosses a fixed level below that threshold
on a sequence of shrinking cylinders.  Choose its first crossing on each
dyadic exhaustion.  The point-selection argument of
\cref{prop:first-bad-scale-descendant} supplies target-retaining normalized
windows with a common quotient floor and the stated local bounds; any first
pressure or realization failure is retained in the same ordered marked
record rather than discarded.

Use \cref{prop:coherent-observer-interpolation} to join these selected charts
by absolutely continuous center and scale parameters.  Every chart is an
exact change of variables of the same suitable continuation, so the
composition law holds and the physical time images approach (T).  Finally,
\cref{prop:canonical-orbit-compactness-retention} gives the asserted product
compactness, pressure convergence, and persistence of the quotient floor.
Passing once to the diagonal subsequence in that proposition produces the
single orbit and the pairwise disjoint windows required in clauses
\textup{(i)}--\textup{(iv)}.
\end{proof}

\section{Gauge, pressure, and adjoint estimates}
\label{app:gauge-pressure}

\subsection{Scale--center normalization}
\label{app:scale-center-normalization}

\begin{proposition}[Local scale--center slice]
\label{prop:local-scale-center-slice}
On every retained window for which the Jacobian
\eqref{eq:gauge-jacobian} is invertible, the equations
\(\mathfrak F=0\) determine a unique local scale and center.  Along a
represented solution these parameters are absolutely continuous and
\(a,b\in L^1_{\mathrm{loc}}\).  A degenerating sequence either loses the
quotient target or has a first-bad-scale descendant with a nondegenerate
slice.
\end{proposition}

\begin{proof}
The first assertion is the implicit-function theorem applied to
\eqref{eq:moment-map}.  Differentiate the four moment equations along the
weakly regularized flow.  The coefficient matrix is
\eqref{eq:gauge-jacobian}; its inverse is bounded on a compact
nondegenerate window.  The right-hand side consists of the localized equation
pairings and the derivatives of both fixed-cutoff means in
\eqref{eq:fixed-cutoff-mean}, all in \(L^1\).  Hence the scale and center are
absolutely continuous and their logarithmic velocities \(a,b\) are locally
integrable.  If the determinant tends to zero while the quotient activity
persists, select the first smaller parabolic scale at which the fixed
activity threshold is met.  \Cref{prop:first-bad-scale-descendant} gives an
ordinary nonhomogeneous limit or a retained pressure/scale output.  Repeating
the moment selection on that first scale yields a nondegenerate slice; a
homogeneous limit is target-null.
\end{proof}

\subsection{Local pressure decomposition}
\label{app:local-pressure-decomposition}

\begin{proposition}[Pressure normalization and compactness]
\label{prop:pressure-normalization-compactness}
On nested inner balls one has
\[
 P=P_{\mathrm{loc}}+H,
 \qquad
 -\Delta P_{\mathrm{loc}}
 =\partial_i\partial_j(\zeta V_iV_j),
 \qquad \Delta H=0,
\]
with compatible spatial means.  On compact normalized families the local
part is compact in \(L^{3/2}\), and the harmonic part is compact modulo its
affine Taylor jet.  A non-affine failure occurs at an accessible first
parabolic scale.
\end{proposition}

\begin{proof}
Apply the whole-space pressure formula after multiplying the velocity by a
cutoff equal to one on the inner ball.  Calder\'on--Zygmund estimates give
\(\|P_{\rm loc}\|_{3/2}\lesssim\|V\|_3^2\), and strong local \(L^3\)
convergence gives compactness of this part.  The remainder is harmonic, so
interior derivative estimates control it after subtraction of its constant
and linear Taylor terms.  The constant is the pressure gauge and the linear
term pairs with the removed homogeneous velocity.  If a non-affine harmonic
remainder is not compact, \cref{prop:harmonic-pressure-tail-accessibility}
places its source inside the available parabolic depth, and
\cref{prop:first-bad-scale-descendant} selects its first visible scale.
Affine parabolic chart changes preserve this decomposition and compatible
means.
\end{proof}

\subsection{Adjoint localization on marked blocks}
\label{app:normalized-adjoint-solutions}

\begin{proposition}[Covariant adjoint localization]
\label{prop:covariant-adjoint-localization}
On every finite marked block, the normalized backward adjoint problem
\eqref{eq:normalized-adjoint} has a positive solution of unit mass.
Under the chart changes of the canonical observer it transforms covariantly,
and the localized kinetic identity is
\cref{thm:moving-cutoff-monotonicity}.  Every cutoff, pressure, translation,
and scale term appearing there is exactly a coordinate of the signed current
\(H_j\) in \cref{thm:exact-critical-financing}.
\end{proposition}

\begin{proof}
Existence, positivity, normalization, and weak stability are
\cref{lem:weak-adjoint}.  Covariance follows by changing variables in the
backward equation with \(t'=\Lambda^2\); the mass normalization removes the
Jacobian.  Testing the suitable local-energy inequality by the transported
cutoff and subtracting the weighted mean equation gives
\cref{thm:moving-cutoff-monotonicity}.  Comparing its displayed terms with
\eqref{eq:signed-current} identifies them one by one with the complete
moving-root current.  Pressure-gauge invariance follows from
\(\int V\cdot\nabla\Phi=0\).
\end{proof}

\section{Critical flux--dissipation balance}
\label{app:scale-space-current}

This appendix supplies the localization algebra behind
\cref{thm:exact-critical-financing}.  The local energy equation separates
positive dissipation and defect from signed transport.  Compatible smooth
localizations give an exact refinement identity, while disjoint spacetime
cells give cancellation of oppositely oriented internal traces.  The outer
scale-space boundary remains in the complete signed current \(H_j\).

\subsection{Compatible nested parabolic cells}
\label{app:nested-parabolic-cells}

Let \(\{\psi_\alpha\}_{\alpha\in\mathcal T_0}\) be a finite family of
nonnegative smooth local-energy tests on one represented time interval.  It
is compatible when every nonterminal test is the sum of its children.  The
family may be obtained from a smooth partition of a finite nested collection
of parabolic cells; only the partition identity is used below.

\subsection{Summed local energy identity}
\label{app:summed-local-energy}

\begin{proposition}[Refinement identity for local-energy tests]
\label{prop:finite-scale-space-cancellation}
For a compatible finite family, the sum of the signed currents of the
terminal tests equals the signed current of the root test.  The same equality
holds separately for kinetic energy, dissipation, and suitable-energy defect.
Consequently replacing a parent by its children changes none of the terms in
the balance and introduces no additional pressure or convection contribution.
The identity is invariant under \(P\mapsto P+c(\tau)\).
\end{proposition}

\begin{proof}
Iterate \eqref{eq:current-additivity} from the leaves to the root.  Each term
in \eqref{eq:signed-current} is linear in the test and its derivatives, so
the asserted equality holds term by term.  Pressure-gauge invariance follows
from \(\int V\cdot\nabla\psi=0\) for every compactly supported test.
\end{proof}

This proposition identifies the aggregate root flux but does not discard its
outer boundary term.  In the single proof spine that term is precisely the
signed financing current.  Sublinear growth is excluded by the critical
price; persistent growth is compactified by
\cref{thm:persistent-current-compactification}.  Thus no separate
flux-domination alternative is required.

\section{Target-visible coercivity and zero-dissipation limits}
\label{app:material-rigidity}

The dissipation need not control every material direction uniformly.  The
singular target detects only the Galilean-invariant oscillatory component of
the velocity.  Compactness and the local Dirichlet form are coercive on that
quotient: a sequence on which the form vanishes converges to a spatially
constant velocity and therefore disappears from the target.  This gives the
required coercivity-or-reduction alternative without estimating the rotation
of the material eigenframe.

\subsection{Zero-dissipation limits are target-null}
\label{app:compact-zero-dissipation-limits}

\begin{theorem}[Target-nullity of zero-dissipation limits]
\label{thm:target-null-zero-dissipation}
Let \(\mathfrak w_j=(V_j,P_j,a_j,b_j;[0,1])\) be normalized windows in the
compact closure specified by
\cref{thm:canonical-critical-orbit}.  If
\begin{equation}\label{eq:vanishing-window-dirichlet}
 \int_0^1\int_{\R^3}\zeta^2|\nabla V_j|^2\dd y\dd\tau
 \longrightarrow0,
\end{equation}
then, after passage to a subsequence,
\begin{equation}\label{eq:constant-zero-dissipation-limit}
 V_j\longrightarrow c(\tau)
 \quad\hbox{strongly in }L^3(\supp\zeta\times[0,1])
\end{equation}
for a spatially constant vector field \(c\).  Consequently
\begin{equation}\label{eq:zero-limit-target-null}
 \OscMass_\Gamma(c;[0,1])=0.
\end{equation}
In particular, no sequence satisfying the fixed target floor
\eqref{eq:target-window-floor} can satisfy
\eqref{eq:vanishing-window-dirichlet}.
\end{theorem}

\begin{proof}
The compactness in \eqref{eq:critical-orbit-compactness} gives, after
extraction,
\[
 V_j\to V_\infty
 \quad\hbox{strongly in }L^3(\supp\zeta\times[0,1]),
 \qquad
 \nabla V_j\rightharpoonup\nabla V_\infty
 \quad\hbox{weakly in }L^2_{\mathrm{loc}}.
\]
Weak lower semicontinuity and
\eqref{eq:vanishing-window-dirichlet} imply
\[
 \int_0^1\int\zeta^2|\nabla V_\infty|^2\dd y\dd\tau=0.
\]
The set \(\{\zeta>0\}\) is connected, so for almost every \(\tau\) there is
a vector \(c(\tau)\) with
\(V_\infty(y,\tau)=c(\tau)\) on that set.  This proves
\eqref{eq:constant-zero-dissipation-limit}.  The weighted mean of a constant
is the same constant, hence the Galilean-invariant oscillatory integrand
vanishes and \eqref{eq:zero-limit-target-null} follows.  Strong
\(L^3\)-convergence passes the fixed-cutoff oscillatory functional to the
limit by \cref{prop:entropy-compactification}.  The positive target floor
would therefore pass to a limit with zero target value, which is impossible.
\end{proof}

\subsection{Coercivity on the target quotient}
\label{app:target-quotient-coercivity}

\begin{theorem}[Target-visible coercivity-or-reduction]
\label{thm:target-visible-coercivity-reduction}
Let \(\mathcal K\) be the compact closure of the selected windows of one
canonical orbit, and let
\[
 \mathcal K_{\eta_0}
 :=\{\mathfrak w\in\mathcal K:
       \OscMass_\Gamma(V;[0,1])\ge\eta_0\}.
\]
Exactly one of the following alternatives holds:
\begin{enumerate}[label=\textup{(\roman*)},leftmargin=2.2em]
\item \(\mathcal K_{\eta_0}=\varnothing\), so the compact orbit closure has
lost the singular target in the Galilean quotient;
\item \(\mathcal K_{\eta_0}\ne\varnothing\), and there exists
\(\delta_\nu>0\) such that
\begin{equation}\label{eq:target-quotient-dirichlet-gap}
 \nu\int_0^1\int\zeta^2|\nabla V|^2\dd y\dd\tau
 \ge\delta_\nu
 \qquad(\mathfrak w\in\mathcal K_{\eta_0}).
\end{equation}
\end{enumerate}
For the canonical target-retaining orbit, alternative~\textup{(i)} is
excluded by \eqref{eq:target-window-floor}; hence
\eqref{eq:target-quotient-dirichlet-gap} holds on every selected window.
\end{theorem}

\begin{proof}
The target functional is continuous under the strong local \(L^3\)
convergence, so \(\mathcal K_{\eta_0}\) is compact.  Assume first that it is
nonempty.  If the infimum of the
left side of \eqref{eq:target-quotient-dirichlet-gap} were zero, a minimizing
sequence would satisfy \eqref{eq:vanishing-window-dirichlet}.  By
\cref{thm:target-null-zero-dissipation}, its limit would have zero
Galilean-invariant oscillation, contradicting membership in
\(\mathcal K_{\eta_0}\).  The infimum is therefore positive.  If this set is
empty, alternative~\textup{(i)} holds; otherwise the positive infimum gives
alternative~\textup{(ii)}.
\end{proof}

The preceding theorem is the material reduction needed by the regularity
argument.  The null space consists of spatially constant velocities on the
retained connected core, which are precisely the local Galilean directions
removed by the target observable.  Material stretching, eigenvalue crossing,
and rotation of compressed directions may be arbitrarily large inside this
null class; none remains visible to the singular target.

\subsection{Divergence on a persistent target sequence}
\label{app:target-visible-dissipation-divergence}

\begin{corollary}[Divergence of normalized target-visible dissipation]
\label{cor:target-visible-dissipation-divergence}
For the infinite pairwise disjoint target-retaining windows of a canonical
orbit,
\begin{equation}\label{eq:divergent-normalized-dirichlet-series}
 \sum_{j=1}^{\infty}
 \nu\int_{I_j}\int\zeta^2|\nabla V|^2\dd y\dd\tau
 =\infty.
\end{equation}
If a nonnegative measure \(\mu\) satisfies
\begin{equation}\label{eq:critical-measure-dominates-dirichlet}
 \mu(I_j)\ge c_0\nu
 \int_{I_j}\int\zeta^2|\nabla V|^2\dd y\dd\tau,
 \qquad c_0>0,
\end{equation}
then \(\mu\) cannot have finite mass on the terminal tail.
\end{corollary}

\begin{proof}
Every summand in \eqref{eq:divergent-normalized-dirichlet-series} is at least
\(\delta_\nu\) by
\cref{thm:target-visible-coercivity-reduction}; there are infinitely many
windows.  If \eqref{eq:critical-measure-dominates-dirichlet} holds,
countable additivity on the disjoint half-open windows gives
\[
 \mu([\tau_0,\infty))
 \ge\sum_j\mu(I_j)
 \ge c_0\sum_j\nu
       \int_{I_j}\int\zeta^2|\nabla V|^2\dd y\dd\tau
 =\infty.
\]
\end{proof}

For the corrected moving-cutoff production,
\eqref{eq:critical-measure-dominates-dirichlet} holds with
\(c_0=1/2\) because \(\Phi=1\) on \(\supp\zeta\).  Thus this realization
needs no estimate of the material condition number and no
classification of rotating eigendirections.  A different physical measure
may carry scale-dependent coefficients; its lower-bound series must then be
proved divergent separately.  The critical localized-energy balance avoids that loss because
its inward feeding and dissipation are compared before returning to physical
Stokes units.

\section{Localized equality and the role of the Gaussian}
\label{app:rigidity}

For the moving-cutoff realization, the proof of
\cref{thm:zero-production-rigidity} uses only the coercive fixed-cutoff
Dirichlet floor and connectedness of \(\{\zeta>0\}\).  These properties are
closed under the convergence proved for the canonical orbit.  The
self-contained signed-current realization has a larger apparent equality class;
its reduction to target-null states is the content of
Appendix~\ref{app:material-rigidity}.

The Gaussian statement has a separate and more limited role.  For a smooth
positive whole-space probability density, vanishing of
\(\mathfrak D_{\mathrm G}\) gives
\(\nabla^2f=(2\nu\theta)^{-1}I\), and normalization and integrability then
give the heat kernel \eqref{eq:heat-kernel} up to translation.  A localized
Dirichlet adjoint density is not that whole-space Gaussian.  The manuscript
therefore uses Gaussian rigidity only as the equality model for the formal
adjoint entropy, not as a premise in the local oscillation contradiction.

\begin{thebibliography}{99}

\bibitem{AlbrittonBarker2019}
D.~Albritton and T.~Barker,
\emph{On local Type I singularities of the Navier--Stokes equations and
Liouville theorems},
J. Math. Fluid Mech. \textbf{21} (2019), Article~43.

\bibitem{Arnold1966}
V.~I. Arnold,
\emph{Sur la g\'eom\'etrie diff\'erentielle des groupes de Lie de dimension
infinie et ses applications \`a l'hydrodynamique des fluides parfaits},
Ann. Inst. Fourier (Grenoble) \textbf{16} (1966), no.~1, 319--361.

\bibitem{CaffarelliKohnNirenberg1982}
L.~Caffarelli, R.~Kohn, and L.~Nirenberg,
\emph{Partial regularity of suitable weak solutions of the Navier--Stokes
equations}, Comm. Pure Appl. Math. \textbf{35} (1982), no.~6, 771--831.

\bibitem{EbinMarsden1970}
D.~G. Ebin and J.~Marsden,
\emph{Groups of diffeomorphisms and the motion of an incompressible fluid},
Ann. of Math. (2) \textbf{92} (1970), no.~1, 102--163.

\bibitem{GustafsonKangTsai2007}
S.~Gustafson, K.~Kang, and T.-P.~Tsai,
\emph{Interior regularity criteria for suitable weak solutions of the
Navier--Stokes equations},
Comm. Math. Phys. \textbf{273} (2007), no.~1, 161--176.

\bibitem{Leray1934}
J.~Leray,
\emph{Sur le mouvement d'un liquide visqueux emplissant l'espace},
Acta Math. \textbf{63} (1934), 193--248.

\bibitem{Lin1998}
F.-H.~Lin,
\emph{A new proof of the Caffarelli--Kohn--Nirenberg theorem},
Comm. Pure Appl. Math. \textbf{51} (1998), no.~3, 241--257.

\bibitem{Onsager1931}
L.~Onsager,
\emph{Reciprocal relations in irreversible processes. I},
Phys. Rev. \textbf{37} (1931), 405--426.

\bibitem{Perelman2002}
G.~Perelman,
\emph{The entropy formula for the Ricci flow and its geometric applications},
arXiv:math/0211159.

\end{thebibliography}

\end{document}
